\ifx\danslelivre\undefined
\documentclass[10pt]{../fichiers-configuration-latex/smfart}
\usepackage{ifthen}
\newboolean{english}
\newboolean{francais}
\setboolean{english}{true}
\setboolean{francais}{false}
\input{../fichiers-configuration-latex/paquets}
\selectlanguage{english}
\input{../fichiers-configuration-latex/lemme-proposition-theoreme}
\input{../fichiers-configuration-latex/numerotation}
\input{../fichiers-configuration-latex/macros}

\externaldocument{../00-introduction/introduction}
\externaldocument{../01-quasi-excellents/excellent}
\externaldocument{../02-topologies/topologies}
\externaldocument{../03-approximation/approximation}
\externaldocument{../04-Cohen-Gabber/Cohen-Gabber}
\externaldocument{../05-algebrisation-partielle/algebrisation}
\externaldocument{../06-log-lissite/log}
\externaldocument{../07-uniformisation-locale/uniformisation-faible}
\externaldocument{../09-uniformisation-premiere-a-l/uniformisation-locale-premiere-a-l}
\externaldocument{../11-produits-orientes/produits-orientes}
\externaldocument{../12-descente-cohomologique-orientee/descente-orientee}
\externaldocument{../13-theoreme-de-finitude/finitude}
\externaldocument{../14-fonctions-dimension/fonctions-dimension}
\externaldocument{../15-Lefschetz-affine/Lefschetz-affine}
\externaldocument{../16-classes-de-Chern/gysin}
\externaldocument{../17-complexes-dualisants/dualite}
\externaldocument{../19-exemples-contre-exemples/contre-exemples}
\externaldocument{../20-rigidite/rigidite}
\externaldocument{../21-finitude-non-abelienne/finitude-non-abelienne}

\thanks{The research of M.T. was partially supported by the
European Union Seventh Framework Programme (FP7/2007-2013) under
grant agreement 268182.}

\begin{document}

\begin{center} XVIII. Gabber's modification theorem (absolute case)\\
Luc Illusie and Michael Temkin\footnote{The research of M.T. was partially supported by the
European Union Seventh Framework Programme (FP7/2007-2013) under grant agreement 268182.}
\end{center}

\tableofcontents

\else
\setboolean{english}{true}
\setboolean{francais}{false}
\selectlanguage{english}
\chapter{Gabber's modification theorem
(absolute case)}
\addtocontents{toc}{par Luc Illusie et Michael Temkin}
\begin{center}
Luc Illusie and Michael Temkin
\footnote{The research of M.T. was partially supported by the European Union Seventh
Framework Programme (FP7/2007-2013) under grant agreement 268182.}

\end{center}
\fi

In this expos\'e we state and prove Gabber's modification theorem
mentioned in the introduction (see step (C)). Its main application is
to Gabber's refined -- i.e. prime to $\ell$ -- local uniformization
theorem. This is treated in expos\'e IX. A relative variant of
the modification theorem, also due to Gabber, has applications to
prime to $\ell$ refinements of theorems of de Jong on alterations of
schemes of finite type over a field or a trait. This is discussed
in expos\'e X.  In \S1, we state Gabber's modification theorem in
its absolute form (Theorem \ref{8-TH. MODIFICATION}). The proof
of this theorem occupies \S\S\ref{8-prelimsec}--\ref{8-endsec}. A
key ingredient is the existence of functorial (with respect to
regular morphisms) resolutions in characteristic zero; the relevant
material is collected in \S\ref{8-functorressec}. We apply it in
\S\ref{8-logregsec} to get resolutions of log regular log schemes,
using the language of Kato's fans and Ogus's monoschemes. The main
results, on which the proof of \ref{8-TH. MODIFICATION} is based,
are Theorems \ref{8-monoidaldesth} and \ref{8-Morel 2.4.2}. \S\S
\ref{8-functorressec} and \ref{8-logregsec} can be read independently
of \S\S \ref{8-mainsec}, \ref{8-prelimsec}, \ref{8-endsec}.

We wish to thank Sophie Morel for sharing with us her notes
on resolution of log regular log schemes and Gabber's magic box.
They were quite useful.

\section{Statement of the main theorem}\label{8-mainsec}
\begin{theorem}[Gabber]\label{8-TH. MODIFICATION} Let $X$ be a
noetherian, qe, separated, log regular fs log scheme (\refext{VI}{6-1.2}),
endowed with an action of a finite group $G$. We assume that $G$
acts tamely (\refext{VI}{6-3.1}) and generically freely on $X$
(i.e. there exists a $G$-stable, dense open subset of $X$ where
the inertia groups $G_{\overline{x}}$ are trivial). Let $Z$ be the
complement of the open subset of triviality of the log structure of
$X$, and let $T$ be the complement of the largest $G$-stable open
subset of $X$ over which $G$ acts freely. Then there exists an fs log
scheme $X'$ and a  $G$-equivariant morphism $f=f_{(G,X,Z)}{\colon}
X' \to X$ of log schemes having the following properties:

(i) As a morphism of schemes, $f$ is a projective modification,
i.e. $f$ is projective and induces an isomorphism of dense open
subsets.

(ii) $X'$ is log regular and $Z' = f^{-1}(Z \cup T)$ is the
complement of the open subset of triviality of the log structure
of $X'$.

(iii) The action of $G$ on $X'$ is very tame (\refext{VI}{6-3.1}).
\end{theorem}

When proving the theorem we will construct $f_{(G,X,Z)}$ that
satisfies a few more nice properties that will be listed in Theorem
\ref{8-addendum}. We remark that Gabber also proves the theorem, 
more generally, when $X$ is not assumed to be qe. However, 
the quasi-excellence assumption simplifies the proof so we impose it here.
Most of the proof works for a general noetherian $X$, so we will assume 
that $X$ is qe only when this will be needed in \S\ref{8-endsec}.

\subsection{}\label{8-remarks main th.} (a) Note that we do
not demand that $f$ is log smooth. In general, it is not. Here
is an example. Let $k$ be an algebraically closed field of
characteristic $\ne 2$. Let $G = \{\pm 1\}$ act on the affine plane
$X = \mathbf{A}^2_k$, endowed with the trivial log structure, by
$x \mapsto \pm x$. Then $X$ is regular and log regular, and $T=
\{0\}$. The action of $G$ on $X$ is tame, but not very tame, as
$G_{\{0\}}$ (= $G$) does not act trivially on the (only) stratum $X$
of the stratification by the rank of $M^{gp}/\OO^*$. Let
$f\colon X' \to X$ be the blow up of $T$, with its natural action
of $G$. Then the pair $(X',Z' = f^{-1}(T))$ is log regular, $f$
is a $G$-equivariant morphism of log schemes, $X'-Z'$ is at the
same time the open subset of triviality of the log structure and
the largest $G$-stable open subset of $X'$ where $G$ acts freely,
and $G$ acts very tamely on $X'$. However, $f$ is not log smooth
(the fiber of $f$ at $\{0\}$ is the line $Z'$ with the log structure
associated to $\mathbf{N} \to \OO, 1 \mapsto 0$, which is
not log smooth over $\Spec\ \!  k$ with the trivial log structure).

(b) In the above example, let $D_1$, $D_2$ be distinct lines
in $X$ crossing at $\{0\}$, and put the log structure $M(D)$ on
$X$ defined by the divisor with normal crossings $D = D_1 \cup
D_2$. Then $\tilde X = (X,M(D))$ is log smooth over $\Spec\ \! k$
endowed with the trivial log structure, and $G$ acts very tamely
on $(X,M(D))$ (\refext{VI}{6-dcn stricts}).  Moreover, the modification $f$
considered above underlies the log blow up $\tilde f\colon \tilde X'
\to \tilde X$ of $X$ at (the ideal in $M(D)$) of $\{0\}$. While
$f$ depends only on $X$, the log \'etale morphism $\tilde f$
is not canonical, as it depends on the choice of $D$. However,
one can recover $f$ from the \textit{canonical} resolutions of
toric singularities (discussed in the next section). Namely,
as $G$ acts very tamely on $\tilde X$, the quotient $\tilde Y =
\tilde X/G$ is log regular (\refext{VI}{6-3.2}): $\tilde Y = \Spec\
\! k[P]$, where $P$ is the submonoid of $\mathbf{Z}^2$ generated
by $(2,0)$, $(1,1)$ and $(0,2)$, and the projection $p\colon
\tilde X \to \tilde Y$ is a Kummer \'etale cover of group $G$,
in particular is a $G$-\'etale cover of $V = p(U)$, where $U = X -
D$. Let $g\colon \tilde Y' \to \tilde Y$ be the log blow up of $\{0\}
= p(\{0\})$ in $\tilde Y$. We then have a cartesian diagram of log
schemes $$ \xymatrix{\tilde X' \ar[d] \ar[r]& \tilde Y' \ar[d]^g \\
\tilde X \ar[r]^{p} & \tilde Y}, $$ where the horizontal maps are
Kummer \'etale covers of group $G$. Now, as a morphism of schemes,
$\tilde Y' \to \tilde Y$ is the \textit{canonical resolution}
of $\tilde Y$, and the underlying scheme $X'$ of $\tilde X'$
is the \textit{normalization} of $\tilde Y'$ in the $G$-\'etale
cover $p\colon U \to V$. This observation, suitably generalized,
plays a key role in the proof of \ref{8-TH. MODIFICATION}.

\section{Functorial resolutions}\label{8-functorressec} All schemes
considered from now on will be assumed to be noetherian.

\subsection{Towers of blow ups}\label{8-blowtowersec} In this
section we review various known results on the following related
topics: blow ups and their towers, various operations on towers,
such as strict transforms and pushforwards, associated points of
schemes and schematic closure.

\subsubsection{Blow ups}\label{8-blowsec} We start with
recalling basic properties of blow ups; a good reference is
\cite[\S1]{Nagata@Conrad}. Let $X$ be a scheme. By a \textit{blow
up} of $X$ we mean a triple consisting of a morphism $f\colon Y
\to X$, a closed subscheme $V$ of $X$ (the \textit{center}), and
an $X$-isomorphism $\alpha\colon Y\iso \mathrm{Proj}(\oplus_{n
\in \mathbf{N}}\calI^n)$, where $\calI=\calI(V)$ is the ideal
of $V$.  We will write $Y=\Bl_V(X)$. When there is no risk of
confusion we will omit $V$ and $\alpha$ from the notation. A blow
up $(f,V,\alpha)$ is said to be \textit{empty} if $V=\emptyset$. In
this case, $Y=X$ and $f$ is the identity. The only $X$-automorphism
of a blow up is the identity. Also, it is well known that $Y$ is
the universal $X$-scheme such that $V\times_XY$ is a Cartier divisor
(i.e. the ideal $\calI\calO_Y$ is invertible).

\subsubsection{Total and strict transforms}\label{8-transformsec}
Given a blow up $f{\colon}Y=\Bl_V(X)\to X$, there are two natural
ways to pullback closed subschemes $i{\colon}Z\into X$. The {\em
total transform} of $Z$ under $f$ is the scheme-theoretic preimage
$f^\tot(Z)=Z\times_XY$. The {\em strict transform} $f^\st(Z)$
is defined as the schematic closure of $f^{-1}(Z-V)\iso Z-V$.

\begin{remark2}\label{8-strictrem} (i) The strict transform depends
on the centers and not only on $Z$ and the morphism $Y\to X$. For
example, if $D\into X$ is a Cartier divisor then the morphism
$\Bl_D(X)\to X$ is an isomorphism but the strict transform of $D$
is empty.

(ii) While $f^\tot(Z)\to Z$ is just a proper morphism, the morphism
$f^\st(Z)\to Z$ can be provided with the blow up structure because
$f^\st(Z)\iso\Bl_{V\times_XZ}(Z)$ (e.g., if $Z\into V$ then
$f^\st(Z)=\emptyset=\Bl_Z(Z)$). Thus, the strict transform can be
viewed as a genuine blow up pullback of $f$ with respect to $i$.
\end{remark2}

\subsubsection{Towers of blow ups}\label{8-towerssec} Next we
introduce blow up towers and study various operations with them
(see also \cite[\S2.2]{Temkin@qe}). By a \textit{tower of blow ups}
of $X$ we mean a finite sequence of length $n\ge 0$ $$ X_{\bullet}
= (\xymatrix{ X_n \ar[r]^{f_{n-1}} & X_{n-1} \ar[r] &\cdots
\ar[r]^{f_0} & X_0 = X}) $$ of blow ups. In particular, this data
includes the centers $V_i\into X_i$ for $0\le i\le n-1$. Usually,
we will denote the tower as $X_\bullet$ or $(X_\bullet,V_\bullet)$. Also,
we will often use notation $X_n\dashto X_0$ to denote a sequence
of morphisms.

If $n=0$ then we say that the tower is {\em trivial}. Note that the
morphism $X_n\to X$ is projective, and it is a modification if and
only if the centers $V_i$ are nowhere dense. If $X_{\bullet}$ is a
tower of blow ups, we denote by $(X_\bullet)_c$ the {\em contracted}
tower deduced from $X_{\bullet}$ by omitting the empty blow ups.

\subsubsection{Strict transform of a tower}\label{8-strictsec}
Assume that $\calX=(X_\bullet,V_\bullet)$ is a blow up tower
of $X$ and $h{\colon}Y\to X$ is a morphism. We claim that
there exists a unique blow up tower $\calY=(Y_\bullet,W_\bullet)$
of $Y$ such that $Y_i\to Y\to X$ factors through $X_i$ and
$W_i=V_i\times_{X_i}Y_i$. Indeed, this defines $Y_0,W_0$ and $Y_1$
uniquely. Since $V_0\times_XY_1=W_0\times_{Y_0}Y_1$ is a Cartier
divisor, $Y_1\to X$ factors uniquely through $X_1$. The morphism
$Y_1\to X_1$ uniquely defines $W_1$ and $Y_2$, etc. We call $\calY$
the {\em strict transform} of $\calX$ with respect to $h$ and denote
it $h^\st(\calX)$.

\begin{remark2}\label{8-strtowerrem} The following observation
motivates our terminology: if $h{\colon}Y\into X$ is a closed
immersion then $Y_i\into X_i$ is a closed immersion and $Y_{i+1}$
is the strict transform of $Y_i$ under the blow up $X_{i+1}\to X_i$.
\end{remark2}

\subsubsection{Pullbacks}\label{8-pullsec} One can also define
a naive base change of $\calX$ with respect to $h$ simply as
$(Y_\bullet,W_\bullet)=\calX\times_XY$. This produces a sequence
of proper morphisms $Y_n\dashto Y_0$ and closed subschemes
$W_i\into Y_i$ for $0\le i\le n-1$. If this datum is a blow up
sequence, i.e. $Y_{i+1}\iso\Bl_{W_i}(Y_i)$, then we say that
$\calX\times_XY$ is the {\em pullback} of $\calX$ and use the
notation $h^*(\calX)=\calX\times_XY$.

\begin{remark2}\label{8-pullrem} The pullback exists if and only if
$\calX\times_XY\iso h^\st(\calX)$. Indeed, this is obvious for
towers of length one, and the general case follows by induction on
the length.  \end{remark2}

\subsubsection{Flat pullbacks}\label{8-flatpullsec} Blow ups are
compatible with flat base changes $h\colon Y\to X$ in the sense that
$\Bl_{V\times_XY}(Y)\iso\Bl_V(X)\times_XY$ (e.g. just compute
these blow ups in the terms of $\Proj$). By induction on length
of blow up towers it follows that pullbacks of blow up towers with
respect to flat morphisms always exist. One can slightly strengthen
this fact as follows.

\begin{remark2}\label{8-nonflatrem} Assume that $X_\bullet$ is a
blow up tower of $X$ and $h{\colon}Y\to X$ is a morphism. If there
exists a flat morphism $f{\colon}X\to S$ such that the composition
$g{\colon}Y\to S$ is flat and the blow up tower $X_\bullet$ is
the pullback of a blow up tower $S_\bullet$ then the pullback
$\h^*(X_\bullet)$ exists and equals to $g^*(S_\bullet)$.  \end{remark2}

\subsubsection{Equivariant blow ups}\label{8-equivsec} Assume
that $X$ is an $S$-scheme acted on by a flat $S$-group scheme
$G$. We will denote by $p,m\colon X_0=G\times_SX\to X$ the
projection and the action morphisms. Assume that $V\into X$ is
a $G$-equivariant closed subscheme (i.e., $V\times_{X}(X_0;m)$
coincides with $V_0=V\times_{X}(X_0;p)$) then the action
of $G$ lifts to the blow up $Y=\Bl_V(X)$. Indeed, the blow up
$Y_0=\Bl_{V_0}(X_0)\to X_0$ is the pullback of $Y\to X$ with respect
to both $m$ and $p$, i.e. there is a pair of cartesian squares $$
\xymatrix{Y_0 \ar@<0.5ex>[d]^{m'}\ar@<-0.5ex>[d]_{p'} \ar[r]& X_0
\ar@<0.5ex>[d]^m\ar@<-0.5ex>[d]_p \\ Y \ar[r] & X} $$ So, we obtain
an isomorphism $Y_0\iso G\times_SY$ (giving rise to the projection
$p'\colon Y_0\to Y$) and a group action morphism $m'\colon Y_0\to Y$
compatible with $m$. Furthermore, the unit map $e{\colon}X\to X_0$
satisfies the condition of Remark \ref{8-nonflatrem} (with $X=S$),
hence we obtain the base change $e'{\colon}Y\to Y_0$ of $e$.

It is now straightforward to check that $m'$ and $e'$ satisfy
the group action axioms, but let us briefly spell this out
using simplicial nerves. The action of $G$ on $V$ defines a
cartesian sub-simplicial scheme $\Ner(G,V)\into\Ner(G,X)$. By
the flatness of $G$ over $S$ and Remark \ref{8-nonflatrem},
$\Bl_{\Ner(G,V)}(\Ner(G,X))$ is cartesian over $\Ner(G)$, hence
corresponds to an action of $G$ on $\Bl_V(X)$.

\subsubsection{Flat monomorphism}\label{8-flatmonsec} Flat
monomorphisms are studied in \cite{Ray@Raynaud}. In particular, it is
proved in \cite[Prop. 1.1]{Ray@Raynaud} that $i:Y\into X$ is a flat
monomorphism if and only if $i$ is injective and for any $y\in Y$ the
homomorphism $\calO_{X,y}\to\calO_{X,x}$ is an isomorphism. Moreover,
it is proved in \cite[Prop. 1.2]{Ray@Raynaud} that in this case $i$
is a topological embedding and $\calO_Y=\calO_X|_{Y}$. In addition
to open immersions, the main source of flat monomorphisms for us
will be morphisms of the form $\Spec(\calO_{X,x})\into X$ and their
base changes.

\subsubsection{Pushouts of ideals}\label{8-pushidealsec}
Let $i{\colon}Y\to X$ be a flat monomorphism,
e.g. $\Spec(\calO_{X,x})\into X$. By the {\em pushout} $U=i_*(V)$
of a closed subscheme $V\into Y$ we mean its schematic image in
$X$, i.e. $U$ is the minimal closed subscheme such that $V\into X$
factors through $U$. It exists by \ega{I}{9.5.1}.

\begin{lemma2}\label{8-pushideallem} The pushout $U=i_*(V)$ extends
$V$ in the sense that $U\times_XY=V$.  \end{lemma2} \begin{proof} By
minimality of $U$, it suffices to show that $V$ admits some extension
to a closed subscheme of $X$. Furthermore, it suffices to extend $V$
to a neighborhood $X'$ of $i(Y)$ because closed subschemes of $X'$
can be extended to $X$ by the schematic image. Thus we can replace
$X$ with $X'$, in particular, by \cite[1.5(i)]{Ray@Raynaud} we
can assume that $i{\colon}Y\to X$ is affine. Then $i$ is a limit
of affine finite type morphisms $X_i\to X$ by \ega{I}{9.6.6},
and by \egalong{IV}{3}{8.8.2 and 8.10.5(iv)} $V$ is the preimage of a
closed subscheme $V_i\into X_i$ for large enough $i$. We claim that
the morphism $Y\to X_i$ factors through a small enough neighborhood
$X'$ of $i(Y)$. Indeed, by quasi-compactness of $Y$ it suffices to
check that for any $y\in Y$ the morphism $\Spec(\calO_{X,y})\to X_i$
factors through a neighborhood of $y$. Since $\Spec(\calO_{X,y})$ is
the limit of all affine neighborhoods of $y$ the latter follows from
\egalong{IV}{3}{8.13.2}. Having such a factoring $Y\to X'\to X_i$ it
remains to note that $U$ is the preimage of $V'=V_i\times_{X_i}X'$,
hence $U$ extends to $V'$ in $X'$, and as we observed earlier this
implies the lemma.  \end{proof}

\subsubsection{Pushouts of blow up towers}\label{8-pushblowsec}
Given a blow up $f{\colon}\Bl_V(Y)\to Y$ and a flat monomorphism
$i{\colon}Y\into X$ we define the pushforward $i_*(f)$ as the blow
up along $U=i_*(V)$. Using Lemma \ref{8-pushideallem} and flat
pullbacks we see that $i^*i_*(f)=f$ and $\Bl_V(Y)\into\Bl_U(X)$ is
a flat monomorphism. So, we can iterate this procedure to construct
pushforward with respect to $i$ of any blow up tower $Y_\bullet$
of $Y$. It will be denoted  $(X_\bullet,U_\bullet)=i_*(Y_\bullet,V_\bullet)$.

\begin{remark2}\label{8-pushblowrem} (i) Clearly,
$i^*i_*(Y_\bullet)=Y_\bullet$.

(ii) In the opposite direction, a blow up tower $(X_\bullet,U_\bullet)$
of $X$ satisfies  $i_*i^*(X_\bullet)=X_\bullet$ if and only if the
preimage of $Y$ in each center $U_i$ of the tower is schematically
dense.  \end{remark2}

\subsubsection{Associated points of a scheme}\label{8-asschemesec}
Recall that a point $x\in X$ is called associated if $m_x$ is an
associated prime of $\calO_{X,x}$, i.e. $\calO_{X,x}$ contains an
element whose annihilator is $m_x$, see \egalong{IV}{2}{3.1.1}. The
set of all such points will be denoted $\Ass(X)$. The following
result is well known but difficult to find in the literature.

\begin{lemma2}\label{8-asschemelem} Let $i{\colon}Y\into X$ be
a flat monomorphism. Then the schematic image of $i$ coincides
with $X$ if and only if $\Ass(X)\subseteq i(Y)$.  \end{lemma2}
\begin{proof} Note that the schematic image of $i$ can be described
as $\Spec(\calF)$, where $\calF$ is the image of the homomorphism
$\phi{\colon}\calO_X\to i_*(\calO_Y)$. Thus, the schematic image
coincides with $X$ if and only if $\Ker(\phi)=0$.

If $i(Y)$ omits a point $x$ then any $m_x$-torsion element
$s\in\calO_{X,x}$ is in the kernel of $\phi_x{\colon}\calO_{X,x}\to
i_*(\calO_Y)_x$ (we use that $\calO_Y=\calO_X|_Y$ by
\S\ref{8-flatmonsec} and since $\ox$ is the support of $s$, the
restriction $s|_Y$ vanishes). So, if there exists $x\in\Ass(X)$
with $x\notin i(Y)$ then $\Ker(\phi)\neq 0$.

Conversely, if the kernel is non-zero then we take $x$ to
be any maximal point of its support and choose any non-zero
$s\in\Ker(\phi_x)$. In particular, $s|_{Y\cap X_x}=0$ and hence
$x\notin i(Y)$. For any non-trivial generalization $y$ of $x$ the
image of $s$ in $\calO_{X,y}$ vanishes because $\Ker(\calO_{X,y}\to
i_*(\calO_Y)_y)=0$ by maximality of $x$. Thus, $x$ is the support
of $s$, and hence $s$ is annihilated by a power of $m_x$. Since $X$
is noetherian, we can find a multiple of $s$ whose annihilator is
$m_x$, thereby obtaining that $x\in\Ass(X)$.  \end{proof}

\subsubsection{Associated points of blow up
towers}\label{8-astowersec} If $(X_\bullet,V_\bullet)$ is a blow up tower
of $X$ then by the set $\Ass(X_\bullet)$ of its associated points
we mean the union of the images of $\Ass(V_i)$ in $X$. Combining
Remark \ref{8-pushblowrem}(i) and Lemma \ref{8-asschemelem} we
obtain the following:

\begin{lemma2}\label{8-astowerlem} Let $i{\colon}Y\into X$ be a
flat monomorphism and let $X_\bullet$ be a blow up tower of $X$. Then
$i_*i^*(X_\bullet)=X_\bullet$ if and only if $\Ass(X_\bullet)\subset i(Y)$.
\end{lemma2}

\subsection{Normalized blow up towers}\label{8-normaltowersec} For
reduced schemes most of the notions, constructions and results of
\S\ref{8-blowtowersec} have normalized analogs. We develop such a
"normalized" theory in this section.

\subsubsection{Normalization}\label{8-normsec} The normalization of
a reduced noetherian scheme $X$, as defined in \ega{III}{6.3.8},
will be denoted $X^\nor$. Recall that normalization is compatible
with open immersions and for an affine $X=\Spec(A)$ its normalization
is $X^\nor=\Spec(B)$ where $B$ is the integral closure of $A$
in its total ring of fractions (which is a finite product of
fields). The normalization morphism $X^\nor\to X$ is integral but
not necessarily finite (since we do not impose the quasi-excellence
assumption on schemes).

\subsubsection{Functoriality}\label{8-functornormsec} Recall
(see expos\'e II) that a morphism $f{\colon}Y\to X$ is called
{\em maximally dominating} if it takes generic points of $Y$ to
generic points of $X$. Normalization is a functor on the category of
reduced schemes with maximally dominating morphisms. Furthermore,
it possesses the following universal property: any maximally
dominating morphism $Y\to X$ with normal $Y$ factors uniquely
through $X^\nor$. (Both claims are local on $Y$ and $X$ and are
obvious for affine schemes.)

\subsubsection{Normalized blow ups}\label{8-normblowsec} By the {\em
normalized blow up} of a reduced scheme $X$ along a closed subscheme
$V$ we mean the morphism $f{\colon}\Bl_V(X)^\nor\to X$. Note that $f$
is universally closed but does not have to be of finite type. As
in the case of usual blow ups, $V$ is a part of the structure. In
particular, $\Bl_V(X)$ has no $X$-automorphisms and we can talk
about equality of normalized blow ups (as opposed to an isomorphism).

\begin{proposition2}\label{8-normblowprop} (i) Keep the above
notation. Then $\Bl_V(X)^\nor\to X$ is the universal maximally
dominating morphism $Y\to X$ such that $Y$ is normal and $V\times_XY$
is a Cartier divisor.

(ii) For any blow up $f{\colon}Y=\Bl_V(X)\to X$ its normalization
$f^\nor{\colon}Y^\nor\to X^\nor$ is the normalized blow up along
$V\times_XX^\nor$.  \end{proposition2} \begin{proof} Combining the
universal properties of blow ups and normalizations we obtain (i),
and (ii) is its immediate corollary.  \end{proof}

Towers of normalized blow ups and their transforms can now be
defined similarly to their non-normalized analogs.

\subsubsection{Towers of normalized blow ups}\label{8-normtowersec}
A {\em tower of normalized blow ups} is a finite sequence $X_n\dashto
X_0$ of normalized blow ups with centers $V_i\into X_i$ for $0\le
i\le n-1$. The centers are part of the datum.

\subsubsection{Normalization of a blow up tower}\label{8-norblowsec}
Using induction on length and Proposition \ref{8-normblowprop}(ii),
we can associate to a blow up tower $\calX=(X_\bullet,V_\bullet)$
of a reduced scheme $X$ a normalized blow up tower
$\calX^\nor=(Y_\bullet,W_\bullet)$, where $Y_0=X$, $Y_i=X_i^\nor$ for
$i\ge 1$, and $W_i=V_i\times_{X_i}X^\nor_i$. We call $\calX^\nor$
the {\em normalization} of $\calX$.

\subsubsection{Strict transforms}\label{8-strnorsec}
If $\calX=(X_\bullet,V_\bullet)$ is a normalized blow up
tower of $X=X_0$ and $f{\colon}Y\to X$ is a morphism
then we define the strict transform $f^\st(\calX)$ as
the normalized blow up tower $(Y_\bullet,W_\bullet)$ such that
$Y_0=Y$ and $W_i=V_i\times_{X_i}Y_i$. Using induction on the
length and the universal property of normalized blow ups, see
\ref{8-normblowprop}(i), one shows that such a tower exists and is
the universal normalized blow up tower of $Y$ such that $f=f_0$
extends to a compatible sequence of morphisms $f_i{\colon}Y_i\to
X_i$.

\subsubsection{Pullbacks}\label{8-normpullsec} The strict transform
$f^\st(\calX)$ as above will be called the {\em pullback} and denoted
$f^*(\calX)$ if $Y_i\iso X_i\times_XY$ for any $0\le i\le n$.

\begin{lemma2}\label{8-smoothpulllem} If $f{\colon}Y\to X$ is
regular then any normalized blow up tower $\calX$ of $X$ admits
a pullback $f^*(\calX)$.  \end{lemma2} \begin{proof} Blow ups
are compatible with flat morphisms hence we should only show that
normalization is compatible with regular morphisms: if $f{\colon}Y\to
X$ is a regular morphism of reduced schemes then the morphism
$h{\colon}Y^\nor\to X^\nor\times_XY$ is an isomorphism. Note that
$h$ is an integral morphism which is generically an isomorphism and
the target is normal because it is regular over a normal scheme,
see \cite[21.E(iii)]{CA@Matsumura}. Hence $h$ is an isomorphism.
\end{proof}

\begin{remark2} Using \egalong{IV}{2}{6.14.1} instead of the
reference in the proof, one obtains that the claim of the lemma
holds, more generally, whenever $f$ is a normal morphism (i.e. flat
with geometrically normal fibres). We will use in the sequel only
the case when $f$ is regular.  \end{remark2}

\subsubsection{Fpqc descent of blow up towers}\label{8-descsubsec}
The classical fpqc descent of ideals (and modules) implies that
there is also an fpqc descent for blow up towers. Namely, if $Y\to
X$ is an fpqc covering and $Y_\bullet$ is a blow up tower of $Y$
whose both pullbacks to $Y\times_XY$ are equal then $Y_\bullet$
canonically descends to a blow up tower of $X$ because the centers
descend. In the same way, normalized blow up towers descend with
respect to quasi-compact surjective regular morphisms.

\subsubsection{Associated points}\label{8-assnorsec} The
material of \S\S\ref{8-pushblowsec}--\ref{8-astowersec} extends
to normalized blow up towers almost verbatim. In particular,
if $\calX=(X_\bullet,V_\bullet)$ is such a tower then $\Ass(\calX)$
is the union of the images of $\Ass(V_j)$ and for any flat
monomorphism $i{\colon}Y\into X$ (which is a regular morphism by
\S\ref{8-flatmonsec}) with a blow up tower $\calY$ of $Y$ we always
have that $i^*i_*\calY=\calY$, and we have that $i_*i^*\calX=\calX$
if and only if $\Ass(\calX)\subset i(Y)$.

\subsection{Functorial desingularization}\label{8-desingsec} In
this section we will formulate the desingularization result about
toric varieties that will be used later in the proof of Theorem
\ref{8-TH. MODIFICATION}. Then we will show how it is obtained from
known desingularization results.

\subsubsection{Desingularization of a scheme}\label{8-desschsec}
By a {\em resolution} (or {\em desingularization}) {\em tower} of
a scheme $X$ we mean a tower of blow ups with nowhere dense centers
$X_{\bullet}$ such that $X=X_0$, $X_n$ is regular and no $f_i$ is an
empty blow up. For example, the trivial tower is a desingularization
if and only if $X$ itself is regular.

\subsubsection{Normalized desingularization}\label{8-normdessec} We
will also consider normalized blow up towers such that each center
is non-empty and nowhere dense, $X=X_0$ and $X_n$ is regular. Such
a tower will be called a {\em normalized desingularization tower}
of $X$.

\begin{remark2}\label{8-normdesrem} (i) For any desingularization
tower $\calX$ of $X$ its normalization $\calX^\nor$ is a normalized
desingularization tower of $X$.

(ii) Usually one works with non-normalized towers; they are subtler
objects that possess more good properties. All known constructions of
functorial desingularization (see below) produce blow up towers by
an inductive procedure, and one cannot work with normalized towers
instead. However, it will be easier for us to deal with normalized
towers in log geometry because in this case one may work only with
fs log schemes.  \end{remark2}

\subsubsection{Functoriality of
desingularization}\label{8-functorsec} For concreteness, we will
consider desingularizations in \S\ref{8-functorsec}, but all what
we say holds for normalized desingularizations too. Assume that a
family $\calS^0$ of schemes is provided with desingularizations
$\calF(X)=X_\bullet$ for any $X\in\calS^0$. We say that the
desingularization (family) $\calF$ is {\em functorial} with respect
to a family $\calS^1$ of morphisms between the elements of $\calS^0$
if for any $f{\colon}Y\to X$ from $\calS^1$ the desingularization of
$X$ induces that of $Y$ in the sense that $f^*\calF(X)$ is defined
and its contraction coincides with $\calF(Y)$ (so, $\mathscr{F}(Y)
= (Y \times_X \mathscr{F}(X))_c$). Note that we put the $=$ sign
instead of an isomorphism sign, which causes no ambiguity by the
fact that any automorphism of a blow up is the identity as we
observed above.

\begin{remark2}\label{8-functorrem} (i) Contractions in the pulled
back tower appear when some centers of $\calF(X)$ are mapped to
the complement of $f(Y)$ in $X$. In particular, if $f\in\calS^1$
is surjective then the precise equality $\mathscr{F}(Y) = Y \times_X
\mathscr{F}(X)$ holds.

(ii) Assume that $X=\cup_{i=1}^nX_i$ is a Zariski covering and
the morphisms $X_i\into X$ and $\coprod_{i=1}^nX_i\to X$ are
in $\calS^1$. In general, one cannot reconstruct $\calF(X)$
from $\calF(X_i)$'s because the latter are contracted
pullbacks and it is not clear how to glue them with correct
synchronization. However, all information about $\calF(X)$ is
kept in $\calF(\coprod_{i=1}^nX_i)$. The latter is the pullback
of $\calF(X)$ hence we can reconstruct $\calF(X)$ by gluing the
restricted blow up towers $\calF(\coprod_{i=1}^nX_i)|_{X_i}$. Note
that $\calF(\coprod_{i=1}^nX_i)|_{X_i}$ can be obtained from
$\calF(X_i)$ by inserting empty blow ups, and these empty blow ups
make the gluing possible. This trick with synchronization of the
towers $\calF(X_i)$ by desingularizing disjoint unions is often
used in the modern desingularization theory, and one can formally
show (see \cite[Rem. 2.3.4(iv)]{Temkin@qe}) that such approach is
equivalent to the classical synchronization of the algorithm with
an invariant.

(iii) Assume that $\calS^1$ contains all identities $\Id_X$
with $X\in\calS^0$ and a morphism $Y\coprod Z\to X$ is in
$\calS^1$ whenever its restrictions $Y\to X$ and $Z\to X$ are in
$\calS^1$. As an illustration of the above trick, let us show that
even if $f,g{\colon}Y\to X$ are in $\calS^1$ but not surjective,
we have an equality $\calF(X)\times_X(Y,f)=\calF(X)\times_X(Y,g)$
of non-contracted towers. Indeed, set $Y'=Y\coprod X$ and consider
the morphisms $f',g'{\colon}Y'\to X$ that agree with $f$ and
$g$ and map $X$ by identity. Then $\calF(X)\times_X(Y',f')$
and $\calF(X)\times_X(Y',g')$ are equal because $f'$ and $g'$
are surjective, hence their restrictions onto $Y$ are also
equal, but these are precisely $\calF(X)\times_X(Y,f)$ and
$\calF(X)\times_X(Y,g)$.  \end{remark2}

\subsubsection{Gabber's magic box}\label{8-magicboxsec} Now we have
tools to formulate the aforementioned desingularization result.

\begin{theorem2}\label{8-magicboxth} Let $\calS^0$
denote the family of disjoint unions of affine toric
varieties over $\QQ$, i.e. $X\in\calS^0$ if it is of the form
$\coprod_{i=1}^n\Spec(\QQ[P_i])$, where $P_1\. P_n$ are fs torsion
free monoids. Let $\calS^1$ denote the family of smooth morphisms
$$f\colon\coprod_{j=1}^m\Spec(\QQ[Q_j])\to\coprod_{i=1}^n\Spec(\QQ[P_i])$$
such that for each $1\le j\le m$ there exists $1\le i=i(j)\le n$ and a
homomorphism of monoids $\phi_j\colon P_i\to Q_j$ so that the restriction
of $f$ onto $\Spec(\QQ[Q_j])$ factors through the toric morphism
$\Spec(\QQ[\phi_j])$. Then there exists a desingularization $\calF$
on $\calS^0$ which is functorial with respect to $\calS^1$ and,
in addition, satisfies the following compatibility condition:
if $\calO_1\.\calO_l$ are complete noetherian rings containing
$\QQ$, $Z=\coprod_{i=1}^l\Spec(\calO_i)$, and $g,h{\colon}Z\to X$
are two regular morphisms with $X\in\calS^0$ then \begin{equation}
\label{8-equalizer2a} (Z,g) \times_X \mathscr{F}(X) = (Z,h) \times_X
\mathscr{F}(X) \end{equation} \end{theorem2}

Before showing how this theorem follows from known desingularization
results, let us make few comments.

\begin{remark2} (i) Gabber's original magic box also requires that
the centers are smooth schemes. This (and much more) can also be
achieved as will be explained later, but we prefer to emphasize
the minimal list of properties that will be used in the proof of
Theorem \ref{8-TH. MODIFICATION}.

(ii) It is very important to allow disjoint unions in the theorem
in order to deal with synchronization issues, as explained in
Remark \ref{8-functorrem}(ii). This theme will show up repeatedly
throughout the paper.  \end{remark2}

\subsubsection{Desingularization of qe schemes
over $\mathbf{Q}$}\label{8-desingQsec} Gabber's magic box
\ref{8-magicboxth} is a particular case of the following
theorem, see \cite[Th. 1.2.1]{Temkin@qe}. Indeed, due to Remark
\ref{8-functorrem}(iii), functoriality with respect to regular
morphisms implies (\ref{8-equalizer2a}).

\begin{theorem2}\label{8-qeth} There exists a desingularization
algorithm $\calF$ defined for all reduced quasi-excellent schemes
over $\QQ$ and functorial with respect to all {\em regular}
morphisms. In addition, $\calF$ blows up only regular centers.
\end{theorem2}

\begin{remark2}
Although this is not stated in \cite{Temkin@qe}, one can strengthen
Theorem \ref{8-qeth} by requiring that $\calF$ blows up only regular centers
contained in the singular locus. An algorithm $\calF$ is constructed in
\cite{Temkin@qe} from an algorithm $\calF_\Var$ that desingularizes varieties
of characteristic zero, and one can check that if the centers of $\calF_\Var$ lie
in the singular loci (of the intermediate varieties) then the same is true
for $\calF$. Let us explain how one can choose an appropriate $\calF_\Var$.
In \cite{Temkin@qe}, one uses the algorithm of Bierstone-Milman to construct $\calF$,
see Theorem 6.1 and its Addendum in \cite{BMT@Bierstone-Milman-Temkin}
for a description of this algorithm and its properties.
It follows from the Addendum that the algorithm blows up centers lying in the
singular loci until $X$ becomes smooth, and then it performs some additional blow ups
to make the exceptional divisor snc. Eliminating the latter blow ups we obtain
a desingularization algorithm $\calF_\Var$ which only blows up regular
centers lying in the singular locus.
\end{remark2}

It will be convenient for us to use the algorithm $\calF$
from Theorem \ref{8-qeth} in the sequel. Also, to simplify the
exposition we will freely use all properties of $\calF$ but the
careful reader will notice that only the properties of Gabber's
magic box will be crucial in the end. Also, instead of working with
$\calF$ itself we will work with its normalization $\calF^\nor$
which assigns to a reduced qe scheme over $\QQ$ the normalized
blow up tower $\calF(X)^\nor$. It will be convenient to use the
notation $\tilcalF=\calF^\nor$ in the sequel.

\begin{remark2}\label{8-qenorth} (i) Since normalized blow ups
are compatible with regular morphisms, it follows from Theorem
\ref{8-qeth} that the normalized desingularization $\tilcalF$
is functorial with respect to all regular morphisms.

(ii) The feature which is lost under normalization (and which is
not needed for our purposes) is some control on the centers. The
centers $\tilV_i$ of $\tilcalF(X)$ are preimages of the centers
$V_i\into X_i$ of $\calF(X)$ under the normalization morphisms
$X_i^\nor\to X_i$, so they do not have to be even reduced. It will
only be important that $\tilV_i$'s are equivariant when a smooth
group acts on $X$. In the original Gabber's argument it was important
to blow up only regular centers because they were not part of the
blow up data, and one used that a regular center without codimension one
components intersecting the regular locus is determined already by the
underlying morphism of the blow up.  \end{remark2}

\subsubsection{Alternative desingularization
inputs}\label{8-altersec} For the sake of completeness, we discuss
how other algorithms could be used instead of $\calF$. Some
desingularization algorithms for reduced varieties over $\QQ$
are constructed in \cite{BM@resolution}, \cite{Wlodarczyk} and
\cite{Kollar@resolution}. They all are functorial with respect to
equidimensional smooth morphisms (though usually one "forgets"
to mention the equidimensionality restriction). It is shown in
\cite[\S6.3]{BMT@Bierstone-Milman-Temkin} how to make the algorithm
of \cite{BM@resolution} fully functorial by a slight adjusting of
the synchronization of its blow ups. There are also a few other canonical
constructions (e.g. due to Villamayor), which are probably functorial
with respect to equidimensional smooth morphisms too. All these
algorithms can be used to produce a desingularization of log regular
schemes (see \S\ref{8-logregsec}), so the only difficulty is in
establishing the compatibility (\ref{8-equalizer2a}).

For the algorithm of \cite{BMT@Bierstone-Milman-Temkin}
it was shown by Bierstone-Milman (unpublished, see
\cite[Rem. 7.1(2)]{BMT@Bierstone-Milman-Temkin}) that the induced
desingularization of a formal completion at a point depends only on
the formal completion as a scheme. This is precisely what we need in
(\ref{8-equalizer2a}).

Finally, there is a much more general result by Gabber, see
Theorem \ref{8-invarth}, whose proof uses Popescu's theorem and the
cotangent complex. It implies that, actually, any desingularization
of reduced varieties over $\QQ$ which is functorial with respect to
smooth morphisms automatically satisfies (\ref{8-equalizer2a}). So,
in principle, any functorial desingularization of varieties over
$\QQ$ could be used for our purposes. Since Gabber's result and
its proof are powerful and novel for the desingularization theory
(and were missed in \cite{BMT@Bierstone-Milman-Temkin}, mainly
due to a not so trivial involvement of the cotangent complex),
we include them in \S\ref{8-invarsec}.

\subsubsection{Invariance of the regular locus}\label{8-consfunsec}
Until the end of \S\ref{8-desingsec} we consider only qe schemes
of characteristic zero, and our aim is to establish a few useful
properties of $\calF$ (and $\tilcalF$) that are consequences of
the functoriality property $\calF$ satisfies. First, we claim that
$\calF$ does not modify the regular locus of $X$, and even slightly
more than that:

\begin{corollary2}\label{8-reglocuscor} All centers of $\calF(X)$
and $\tilcalF(X)$ sit over the singular locus of $X$. In particular,
$X$ is regular if and only if $\calF(X)$ is the trivial tower.
\end{corollary2} \begin{proof} It suffices to study $\calF$. The
claim is obvious for $S=\Spec(\QQ)$ because $S$ does not contain
non-dense non-empty subschemes. By functoriality, $\calF(T)$ is
trivial for any regular $T$ of characteristic zero, because it is
regular over $S$. Finally, if $T$ is the regular locus of $X$ then
$\calF(T)=(\calF(X)\times_XT)_c$ and hence any center $V_i\into X_i$
of $\calF(X)$ does not intersect the preimage of $T$.  \end{proof}

\subsubsection{Equivariance of the
desingularization}\label{8-equivdessec} It is well known that
functorial desingularization is equivariant with respect to
any smooth group action (and, moreover, extends to functorial
desingularization of stacks). For the reader's convenience we
provide an elementary argument.

\begin{corollary2}\label{8-cor BMT group action} Let $S$ be
a qe scheme over $\QQ$, $G$ be a smooth $S$-group and $X$
be a reduced $S$-scheme of finite type acted on by $G$. Then
the action of $G$ on $X$ extends naturally to an action of
$G$ on $\mathscr{F}(X)$ and $\tilcalF(X)$.  \end{corollary2}
\begin{proof} Again, it suffices to study $\calF$. Let $\calF(X)$
be given by $X_n\dashto X_0=X$ and $V_i\into X_i$ for $0\le i\le
n-1$. By $p,m\colon Y=G\times_S X\to X$ we denote the projection
and the action morphisms. Note that $m$ is smooth (e.g. $m$ is the
composition of the automorphism $(g,x)\to(g,gx)$ of $G\times_S X$ and
$p$). Therefore, $\calF(X)\times_{X}(Y;m)=\calF(Y)=\calF(X)\times_SG$
by Theorem \ref{8-qeth} and Remark \ref{8-functorrem}(i). In
particular, $V_0\times_X (Y; p) = V_0\times_X(Y; m)$, i.e. $V_0$
is $G$-equivariant. By \S\ref{8-equivsec}, $X_1$ inherits
a $G$-action. Then the same argument implies that $V_1$ is
$G$-equivariant and $X_2$ inherits a $G$-action, etc.  \end{proof}

\subsection{Complements on functorial
desingularizations}\label{8-invarsec} This section is devoted to
Gabber's result on a certain non-trivial compatibility property
that any functorial desingularization satisfies. It will not be
used in the sequel, so an uninterested reader may safely skip it.

\begin{theorem2}\label{8-invarth} Assume that $S$ is a scheme,
$\calS^0$ is a family of reduced $S$-schemes of finite type and
$\calS^1$ is a family of morphisms between elements of $\calS^0$
such that if $f{\colon}Y\to X$ is smooth and $X\in\calS^0$
then $Y\in\calS^0$ and $f\in\calS^1$. Let $\mathscr{F}$
be a desingularization on $\calS^0$ which is functorial with
respect to all morphisms of $\calS^1$. Then any pair of regular
morphisms $g{\colon}Z\to X$ and $h{\colon}Z\to Y$ with targets
in $\calS^0$ induces the same desingularization of $Z$; namely,
$\calF(X)\times_XZ=\calF(Y)\times_YZ$.  \end{theorem2}

Note that the theorem has no restrictions on the characteristic
(because no such restriction appears in Popescu's theorem). Before
proving the theorem let us formulate its important corollary,
whose main case is when $S=\Spec(k)$ for a field $k$ and $\calS^0$
is the family of all reduced $k$-schemes of finite type.

\begin{corollary2}\label{8-invarcor} Keep the notation of Theorem
\ref{8-invarth}. Then $\calF$ canonically extends to the family
$\widehat{\calS}^0$ of all schemes that admit a regular morphism to
a scheme from $\calS^0$ and the extension is functorial with respect
to all regular morphisms between schemes of $\widehat{\calS}^0$.
\end{corollary2}

The main ingredient of the proof will be the following result that
we are going to establish first.

\begin{proposition2}\label{8-invarprop} Consider a commutative
diagram of schemes $$ \xymatrix{{} & Z \ar[dl]_g \ar[dr]^h  \ar[d]^f
& {} \\ X \ar[dr]_a & Z' \ar[l]_{g'} \ar[r]^{h' }& Y \ar[dl]^b \\
{} & S & {}} $$ such that $a$ and $b$ are of finite type, $g$ and
$h$ are regular and $g'$ is smooth. Then $h'$ is smooth around the
image of $f$.  \end{proposition2}

For the proof we will need the following three lemmas. In the first
one we recall the Jacobian criterion of smoothness, rephrased in
terms of the cotangent complex.

\begin{lemma2}\label{8-invarlem1} Let $f{\colon} X \ra S$ be a
morphism which is locally of finite presentation, and let $x \in
X$. Then the following conditions are equivalent:

(i) $f$ is smooth at $x$;

(ii) $H_1(L_{X/S} \otimes^L k(x)) = 0$.  \end{lemma2} In the lemma
we use the convention $H_i = H^{-i}$, and $L_{X/S}$ denotes the
cotangent complex of $X/S$.  \begin{proof} (i) $\Rightarrow$
(ii) is trivial: as $f$ is smooth at $x$, up to shrinking $X$
we may assume $f$ smooth, then $L_{X/S}$ is cohomologically
concentrated in degree zero and locally free \cite[III
3.1.2]{Illusie@cotangent}. Let us prove (ii) $\Rightarrow$ (i).
We may assume that we have a factorization $$ \xymatrix{X \ar[r]^i
\ar[d]^f & Z \ar[dl]^g \\ S}, $$ where $i$ is a closed immersion of
ideal $I$ and $g$ is smooth. Consider the standard exact sequence
$$ I/I^2 \ra i^*\Omega^1_{Z/S} \ra \Omega^1_{X/S} \ra 0. \leqno
(*) $$ By the Jacobian criterion \egalong{IV}{4}{17.12.1} and
\egazero{IV}{19.1.12}, the smoothness of $f$ at $x$
is equivalent to the fact that the morphism $$ (I/I^2) \otimes k(x)
\ra \Omega^1_{Z/S} \otimes k(x) \leqno (**) $$ deduced from the left
one in (*) is injective.  Now, $(I/I^2) \otimes k(x) = H_1(L_{X/Z}
\otimes^L k(x))$ \cite[III 3.1.3]{Illusie@cotangent}, and (**)
is a morphism in the exact sequence associated with the triangle
deduced from the transitivity triangle $Li^*L_{Z/S} \ra L_{X/S}
\ra L_{X/Z} \ra Li^*L_{Z/S}[1]$ by applying $\otimes^L k(x)$ :
$$ H_1(L_{X/S} \otimes^L k(x)) \ra H_1(L_{X/Z} \otimes^L k(x))
(= (I/I^2) \otimes k(x)) \ra \Omega^1_{Z/S} \otimes k(x).  $$ By
(ii), $H_1(L_{X/S} \otimes^L k(x)) = 0$, hence (**) is injective,
which completes the proof.  \end{proof}

\begin{lemma2}\label{8-invarlem2} Consider morphisms $f{\colon}
X \ra Y$, $g{\colon} Y \ra S$, $h = gf{\colon} X \ra S$, and let
$x \in X$, $y=f(x)\in Y$. Assume that

(i) $H_1(L_{X/S} \otimes ^L k(x)) = 0$

(ii) $H_2(L_{X/Y} \otimes^L k(x)) = 0$.

Then $H_1(L_{Y/S} \otimes^L k(y)) = 0$. In particular, if $g$
is locally of finite presentation then $g$ is smooth at $y$.
\end{lemma2} \begin{proof}It is equivalent to show that $H_1(L_{Y/S}
\otimes^L k(x)) = 0$, and this follows trivially from the exact
sequence $$ H_2(L_{X/Y} \otimes^L k(x)) \ra H_1(L_{Y/S} \otimes^L
k(x)) \ra H_1(L_{X/S} \otimes^L k(x)).  $$ \end{proof}

\begin{lemma2}\label{8-invarlem3} Let $f {\colon} X \ra S$ be a
regular morphism between locally noetherian schemes. Then $L_{X/S}$
is cohomologically concentrated in degree zero and $H_0(L_{X/S})
= \Omega^1_{X/S}$ is flat.  \end{lemma2} \begin{proof} We may
assume $X = \SP B$ and $S = \SP A$ affine. Then, by Popescu's
theorem \cite[1.1]{NPD@Swan}, $X$ is a filtering projective limit
of smooth affine $S$-schemes $X_{\alpha} = \SP B_{\alpha}$. By
\cite[II (1.2.3.4)]{Illusie@cotangent}, we have $$ L_{B/A} =
\colim_{\alpha} L_{B_{\alpha}/A}.  $$ By \cite[III 3.1.2 and II
2.3.6.3]{Illusie@cotangent}, $L_{B_{\alpha}/A}$ is cohomologically
concentrated in degree zero and $H_0(L_{B_{\alpha}/A}) =
\Omega^1_{B_{\alpha}/A}$ is projective of finite type over
$B_{\alpha}$, so the conclusion follows.  \end{proof}

\begin{proof}[Proof of Proposition \ref{8-invarprop}] The question is
local around a point $y = f(x) \in Z'$, $x \in Z$. In view of Lemma
\ref{8-invarlem1}, by Lemma \ref{8-invarlem2} applied to $Z \ra Z'
\ra Y$ it suffices to show that $H_1(L_{Z/Y} \otimes^L k(x)) = 0$
and $H_2(L_{Z/Z'} \otimes^L k(x)) = 0$. As $Z$ is regular over $Y$,
the first vanishing follows from Lemma \ref{8-invarlem3}. For the
second one, consider the exact sequence
 $$
H_2(L_{Z/X} \otimes^L k(x)) \ra H_2(L_{Z/Z'} \otimes^L k(x)) \ra
H_1(L_{Z'/X} \otimes^L k(x)).  $$ By the regularity of $Z/X$ and
Lemma \ref{8-invarlem3}, $H_2(L_{Z/X} \otimes^L k(x)) = 0$. As $Z'$
is smooth over $X$, $H_1(L_{Z'/X} \otimes^L k(x)) = 0$ by Lemma
\ref{8-invarlem1}, which proves the desired vanishing and finishes
the proof.  \end{proof}

\begin{proof}[Proof of Theorem \ref{8-invarth}] Find finite
affine coverings $X=\cup_i X_i$, $Y=\cup_i Y_i$ and $Z=\cup_i
Z_i$ such that $g(Z_i)\subseteq X_i$ and $h(Z_i)\subseteq
Y_i$. Set $X'=\coprod_i X_i$, $Y'=\coprod_i Y_i$ and
$Z'=\coprod_i Z_i$ and let $Z'\to X'$ and $Z'\to Y'$ be the
induced morphisms. It suffices to check that $\calF(X)\times_XZ$
and $\calF(Y)\times_YZ$ become equal after pulling them back to
$Z'$. So, we should check that $(\calF(X)\times_XX')\times_{X'}Z'$
coincides with $(\calF(Y)\times_YY')\times_{Y'}Z'$. The
morphisms $X'\to X$ and $Y'\to Y$ are smooth and hence
contained in $\calS^1$. So, $\calF(X)\times_XX'=\calF(X')$
and similarly for $Y$. In particular, it suffices to prove that
$\calF(X')\times_{X'}Z'=\calF(Y')\times_{Y'}Z'$. This reduces the
problem to the case when all schemes are affine, so in the sequel
we assume that $X,Y$ and $Z$ are affine.

Next, note that it suffices to find factorizations $g = g_0f$ and $h
= h_0f$, where $f{\colon} Z \ra Z_0$ is a morphism with target in
$\calS^0$ and $g_0{\colon} Z_0 \ra X$, $h_0{\colon} Z_0 \ra X$ are
smooth. By Popescu's theorem, one can write $ g {\colon} Z \ra X$ as
a filtering projective limit of affine smooth morphisms $g_{\alpha}
{\colon} Z_{\alpha} \ra X$, $\alpha \in A$. As $Y$ is of finite
type over $S$, $h$ will factor through one of the $Z_{\alpha}$'s
(\egalong{IV}{3}{8.8.2.3}): there exists $\alpha \in A$, $f_{\alpha}
{\colon}Z \ra Z_{\alpha} $, $h_{\alpha} {\colon} Z_{\alpha} \ra Y$
such that $g = g_{\alpha}f_{\alpha}$, $h = h_{\alpha}f_{\alpha}$. By
Proposition \ref{8-invarprop}, $h_{\alpha}$ is smooth around the
image of $f_{\alpha}$, so we can take $Z_0$ to be a sufficiently
small neighborhood of the image of $f_\alpha$.  \end{proof}

\section{Resolution of log regular log schemes}\label{8-logregsec}
Unless said to the contrary, by log structure we mean
a log structure with respect to the \'etale topology. We
will say that a log structure $M_X$ on a scheme $X$ is {\em
Zariski} if $\varepsilon^*\varepsilon_*M_X\iso M_X$, where
$\varepsilon{\colon}X_\et\to X$ is the morphism between the \'etale
and Zariski sites. In this case, we can safely view the log structure
as Zariski log structure $\varepsilon_*M_X$. A similar convention
will hold also for log schemes.

\subsection{Fans}\label{8-fansec} Many definitions/constructions on
log schemes are of "combinatorial nature". Roughly speaking, these
constructions use only multiplication and ignore addition. Naturally,
there exists a category of geometric spaces whose structure
sheaves are monoids, and most of combinatorial constructions
can be described as "pullbacks" of analogous "monoidal"
operations. The first definition of such a category was done by
Kato in \cite{Toric@Kato}. Kato called his spaces {\em fans} to
stress their relation to the classical combinatorial fans obtained
by gluing polyhedral cones. For example, to any combinatorial fan
$C$ one can naturally associate a fan $F(C)$ whose set of points
is the set of faces of $C$. The main motivation for the definition
is that fans can be naturally associated to various log schemes.

It took some time to discover that fans are sort of "piecewise
schemes" rather than a monoidal version of schemes. A more geometric
version of combinatorial schemes was introduced by Deitmar in
\cite{Deitmar@F1}. He called them $\FF_1$-schemes, but we prefer
the terminology of monoschemes introduced by Ogus in his book
in preparation \cite{Ogus@book}. Note that when working with a
log scheme $X$, we use the sheaf $M_X$ in some constructions and
we use its sharpening $\oM_X$ (see \S\ref{8-notmonsec}) in other
constructions. Roughly speaking, monoschemes naturally arise when we
work with $M_X$ while fans naturally arise when we work with $\oM_X$.

In \S\ref{8-logregsec}, we will show that: (a) a functorial
desingularization of toric varieties over $\QQ$ descends to
a desingularization of monoschemes, (b) to give the latter is
more or less equivalent to give a desingularization of fans,
(c) desingularization of fans can be used to induce a monoidal
desingularization of log schemes, (d) the latter induces a
desingularization of log regular schemes, which (at least in some
cases) depends only on the underlying scheme.

In principle, we could work locally, using desingularization of
disjoint unions of all charts for synchronization. In this case, we
could almost ignore the intermediate categories by working only with
fine monoids and blow up towers of their spectra. However, we decided
to emphasize the actual geometric objects beyond the constructions,
and, especially, stress the difference between fans and monoschemes.

\subsubsection{Sharpening}\label{8-notmonsec} For a monoid $M$,
by $M^\times$ (or $M^*$) we denote the group of its invertible
elements, and its {\em sharpening} $\oM$ is defined as $M/M^\times$.

\subsubsection{Localization}\label{8-localmonsec} By {\em
localization} of a monoid $M$ along a subset $S$ we mean the
universal $M$-monoid $M_S$ such that the image of $S$ in $M_S$
is contained in $M_S^\times$. If $M$ is integral then $M_S$ is
simply the submonoid $M[S^{-1}]\subseteq M^\gp$ generated by $M$
and $S^{-1}$. If $M$ is a fine then any localization is isomorphic
to a localization at a single element $f$, and will be denoted $M_f$.

\subsubsection{Spectra of fine monoids}\label{8-specmonoidsec} All
our combinatorial objects will be glued from finitely many spectra of
fine monoids. Recall that with any fine monoid $P$ one can associate
the set $\Spec(P)$ of prime ideals (with the convention that
$\emptyset$ is also a prime ideal) equipped with the Zariski topology
whose basis is formed by the sets $D(f)=\{p\in\Spec(P)|\ f\notin
p\}$ for $f\in P$, see, for example, \cite[\S9]{Toric@Kato}. The
structure sheaf ${M}_P$ is defined by ${M}_P(D(f))=P_f$, and the
sharp structure sheaf $\oM_P={M}_P/{M}_P^\times$ is the sharpening
of ${M}_P$ (we will see in Remark \ref{8-specrem}(iii) that actually
$\oM_P(D(f))=\overline{P_f}=P_f/(P_f^\times)$, i.e. no sheafification
is needed).

\begin{remark2}\label{8-specrem} (i) Since $P\setminus P^\times$
and $\emptyset$ are the maximal and the minimal prime ideals of
$P$, $\Spec(P)$ possesses unique closed and generic points $s$ and
$\eta$. The latter is the only point whose stalk $M_{P,\eta}=P^\gp$
is a group.

(ii) The set $\Spec(P)$ is finite and its topology is the
specialization topology, i.e. $U$ is open if and only if it is
closed under generalizations. (More generally, this is true for any
finite sober topological space, such as a scheme that has finitely
many points.)

(iii) A subset $U\subseteq\Spec(P)$ is affine (and even of the form
$D(f)$) if and only if it is the localization of $\Spec(P)$ at a
point $x$ (i.e. the set of all generalizations of $x$). Any open
covering $U=\cup_i U_i$ of an affine set is trivial (i.e. $U$
is equal to some $U_i$), therefore any functor $\calF(U)$
on affine sets uniquely extends to a sheaf on $\Spec(P)$. In
particular, this explains why no sheafification is needed when
defining $\oM_P$. Furthermore, we see that, roughly speaking, any
notion/construction that is "defined in terms of" localizations
$X_x$ and stalks $M_x$ or $\oM_x$ is Zariski local. This is very
different from the situation with schemes.  \end{remark2}

\subsubsection{Local homomorphisms of monoids}\label{8-localhomsec}
Any monoid $M$ is local because $M\setminus M^\times$ is its unique
maximal ideal. A homomorphism $f{\colon}M\to N$ of monoids is {\em
local} if it takes the maximal ideal of $M$ to the maximal ideal
of $N$. This happens if and only if $f^{-1}(N^\times)=M^\times$.

\subsubsection{Monoidal spaces}\label{8-monoidspsec} A {\em
monoidal space} is a topological space $X$ provided with
a sheaf of monoids $M_X$. A morphism of monoidal spaces
$(f,f^\#){\colon}(Y,M_Y)\to(X,M_X)$ is a continuous map
$f{\colon}Y\to X$ and a homomorphism $f^\#{\colon}f^{-1}(M_X)\to
M_Y$ such that for any $y\in Y$ the homomorphism of monoids
$f_y^\#{\colon}{M}_{X,f(y)}\to{M}_{Y,y}$ is local.

\begin{remark2} Strictly speaking one should have called the above
category the category of locally monoidal spaces and allow non-local
homomorphisms in the general category of monoidal spaces. However,
we will not use the larger category, so we prefer to abuse the
notation slightly.  \end{remark2}

Spectra of monoids possess the usual universal property, namely:

\begin{lemma2}\label{8-monoidsplem} Let $(X,M_X)$ be a monoidal
space and $P$ be a monoid.

(i) The global sections functor $\Gamma$ induces
a bijection between morphisms of monoidal spaces
$(f,f^\#){\colon}(X,M_X)\to(\Spec(P),M_P)$ and homomorphisms
$\phi{\colon}P\to\Gamma(M_X)$.

(ii) If $M_X$ has sharp stalks then $\Gamma$ induces
a bijection between morphisms of monoidal spaces
$(f,f^\#){\colon}(X,M_X)\to(\Spec(P),\oM_P)$ and homomorphisms
$\phi{\colon}P\to\Gamma(M_X)$.  \end{lemma2} \begin{proof} (i)
Let us construct the opposite map. Given a homomorphism $\phi$,
for any $x\in X$ we obtain a homomorphism $\phi_x{\colon}P\to
M_{X,x}$. Clearly, $m=P\setminus\phi_x^{-1}(M^\times_{X,x})$ is a
prime ideal and hence $\phi_x$ factors through a uniquely defined
local homomorphism $P_m\to M_{X,x}$. Setting $f(x)=m$ we obtain
a map $f{\colon}X\to\Spec(P)$, and the rest of the proof of (i)
is straightforward.

If the stalks of $M_X$ are sharp then any morphism
$(X,M_X)\to(\Spec(P),M_P)$ factors uniquely through
$(\Spec(P),\oM_P)$. Also, $\Gamma(M_X)$ is sharp, hence
any homomorphism to it from $P$ factors uniquely through
$\oP$. Therefore, (ii) follows from (i).  \end{proof}

\subsubsection{Fine fans and monoschemes}\label{8-finfansec} A
fine {\em monoscheme} (resp. a fine {\em fan}) is a monoidal space
$(X,{M}_X)$ that is locally isomorphic to $A_P=(\Spec(P),{M}_P)$
(resp. $\oA_P=(\Spec(P),\oM_P)$), where $P$ is a fine monoid. We
say that $(X,M_X)$ is {\em affine} if it is isomorphic to $A_P$
(resp.f $\oA_P$). A morphism of monoschemes (resp. fans) is a
morphism of monoidal spaces. A monoscheme (resp. a fan) is called
{\em torsion free} if it is covered by spectra of $P$'s with torsion
free $P^\gp$'s. It follows from Remark \ref{8-specrem}(iii) that this
happens if and only if all groups $M^\gp_{X,x}$ are torsion free.

\begin{remark2}\label{8-fanrem} (i) Any fs fan is torsion
free because if an fs monoid is torsion free then so is any its
localization. This is not true for general fine fans. For example,
if $\mu_2=\{\pm 1\}$ then $P=\NN\oplus\mu_2\setminus\{(0,-1)\}$
is a sharp monoid with $P^\gp=\ZZ\oplus\mu_2$.

(ii) For any point $x$ of a fine monoscheme (resp. fan) $X$ the
localization $X_x$ that consists of all generalizations of $x$ is
affine. In particular, by Remark \ref{8-specrem}(i) there exists
a unique maximal point generalizing $x$, and hence $X$ is a finite
disjoint union of irreducible components.  \end{remark2}

\subsubsection{Comparison of monoschemes and
fans}\label{8-comparesec} There is an obvious sharpening functor
$(X,{M}_X)\mapsto(X,\oM_X)$ from monoschemes to fans, and there is
a natural morphism of monoidal spaces $(X,\oM_X)\to(X,M_X)$. The
sharpening functor loses information, and one needs to know
${M}_X^\gp$ to reconstruct ${M}_X$ from $\oM_X$ as a fibred product
(see \cite[2.1.8.3]{Ogus@book}). Actually, there are much more fans
than monoschemes. For example, the generic point $\eta\in\Spec(P)$
is open and $\oM_{P,\eta}$ is trivial hence for any pair of fine
monoids $P$ and $Q$ we can glue their sharpened spectra along the
generic points. What one gets is sort of "piecewise scheme" and,
in general, it does not correspond to standard geometric objects,
such as schemes or monoschemes. We conclude that, in general,
fans can be lifted to monoschemes only locally.

\begin{remark2} As a side remark we note that sharpened monoids
naturally appear as structure sheaves of piecewise linear spaces
(a work in progress of the second author on skeletons of Berkovich
spaces). In particular, PL functions can be naturally interpreted
as sections of the sharpened sheaf of linear functions on polytopes.
\end{remark2}

\subsubsection{Local smoothness}\label{8-localsmoothfansec}
A local homomorphism of fine monoids $\phi{\colon}P\to Q$
is called {\em smooth} if it can be extended to an isomorphism
$P\oplus\NN^r\oplus\ZZ^s\iso Q$. The following lemma checks
that this property is stable under localizations.

\begin{lemma2}\label{8-smoothmonoidlem} Assume that $\phi{\colon}P\to
Q$ is smooth and $P',Q'$ are localizations of $P,Q$ such that $\phi$
extends to a local homomorphism $\phi'{\colon}P'\to Q'$. Then $\phi'$
is smooth.  \end{lemma2} \begin{proof} Recall that $P'=P_a$ for $a\in
P$ (notation of \S\ref{8-localmonsec}), and $\phi'$ factors through
the homomorphism $\phi_a{\colon}P'\to Q_a$, which is obviously
smooth. Therefore, replacing $\phi$ with $\phi_a$ we can assume
that $P=P'$. Let $b=(p,n,z)\in Q$ be such that $Q'=Q_b$. Then $p\in
P^\times$ because $P\to Q'$ is local, and hence $Q'$ is isomorphic
to $P\oplus(\NN^r)_n\oplus\ZZ^s$. It remains to note that any
localization of $\NN^r$ is of the form $\NN^{r-t}\oplus\ZZ^t$.
\end{proof}

\subsubsection{Smoothness}\label{8-smoothfansec} The lemma
allows to globalize smoothness: a morphism $f{\colon}Y\to X$
of monoschemes is called {\em smooth} if the homomorphisms of
stalks ${M}_{X,f(y)}\to{M}_{Y,y}$ are smooth. In particular, $X$
is {\em smooth} if its morphism to $\Spec(1)$ is smooth, that is,
the stalks ${M}_{X,x}$ are of the form $\NN^r\oplus\ZZ^s$. In
particular, a smooth monoscheme is torsion free.

Analogous smoothness definitions are given for fans. Moreover,
in this case we can consider only sharp monoids, and then the
group component $\ZZ^s$ is automatically trivial. It follows
that we can rewrite the above paragraph almost verbatim but with
$s=0$. Obviously, a morphism of monoschemes is smooth if and only
if its sharpening is a smooth morphism of fans.

\begin{remark2}\label{8-chartsmoothrem} (i) Recall that
any fine monoscheme $X$ admits an open affine covering
$X=\cup_{x\in X}\Spec(M_{X,x})$. It follows that a morphism
of fine monoschemes $f{\colon}Y\to X$ is smooth if and
only if it is covered by open affine charts of the form
$\Spec(P\oplus\NN^r\oplus\ZZ^s)\to\Spec(P)$.

(ii) Smooth morphisms of fine fans admit a similar local description,
and we leave the details to the reader.  \end{remark2}

\subsubsection{Saturation}\label{8-satursec} As usually, for a
fine monoid $P$ we denote its saturation by $P^\sat$ (it consists
of all $x\in P^\gp$ with $x^n\in P$ for some $n>0$). Saturation is
compatible with localizations and sharpening and hence extends to
a saturation functor $F\mapsto F^\sat$ on the categories of fine
monoschemes (resp. fine fans). We also have a natural morphism
$F^\sat\to F$, which is easily seen to be bijective. So, actually,
$(F,M_F)^\sat=(F,M_F^\sat)$.

\subsubsection{Ideals}\label{8-idealsec} A subsheaf of ideals
$\calI\subseteq M_X$ on a monoscheme $(X,{M}_X)$ is called a
{\em coherent ideal} if for any point $x\in X$ the restriction of
$\calI$ on $X_x$ coincides with $\calI_xM_{X_x}$. (Due to Remark
\ref{8-specrem}(iii) this means that $\calI$ is coherent in the usual
sense, i.e. its restriction on an open affine submonoscheme $U$ is
generated by the global sections over $U$.) We will consider only
coherent ideals, so we will omit the word "coherent" as a rule. An
ideal $\calI\subseteq{M}_X$ is {\em invertible} if it is locally
generated by a single element.

\subsubsection{Blow ups}\label{8-blowmonosec} Similarly to schemes,
for any non-empty ideal $\calI\subseteq\calO_X$ there exists a
universal morphism of monoschemes $h{\colon}X'\to X$ such that the
{\em pullback ideal} $h^{-1}\calI=\calI{M}_{X'}$ is invertible. We
call $\calI$ the {\em center} of the blow up. (One does not have
an adequate notion of closed submonoscheme, so unlike blow ups
of the scheme it would not make sense that "$V(\calI)$" is the
center.) An explicit construction of $X'$ copies its scheme analog:
it is local on the base and for an affine $X=\Spec(P)$ with an ideal
$I\subseteq P$ corresponding to $\calI$ one glues $X'$ from the
charts $\Spec(P[a^{-1}I])$, where $a\in I$ and $P[a^{-1}I]$ is the
submonoid of $P^\gp$ generated by the fractions $b/a$ for $b\in I$
(see \cite[Ch. 2, \S1.6]{Ogus@book} for details).

\begin{remark2}\label{8-blowmonorem} Blow ups induce isomorphisms on
the stalks of $M^\gp$; this is an analog of the fact that blow ups
of schemes along nowhere dense subschemes are birational morphisms.
\end{remark2}

\subsubsection{Saturated blow ups}\label{8-satblowsec} Analogously
to normalized blow ups, one defines {\em saturated blow up} of a
monoscheme $X$ along an ideal $\calI\subset M_X$ as the saturation
of $\Bl_\calI(X)$. The same argument as for schemes shows that
$\Bl_\calI(X)^\sat$ is the universal saturated $X$-monoscheme such
that the pullback of $\calI$ is invertible.

\subsubsection{Towers and pullbacks}\label{8-towpullsec} Towers
of (saturated) blow ups of a monoscheme $X$ are defined in the
obvious way. Given such a tower $X_\bullet$ with $X=X_0$ and a
morphism $f{\colon}Y\to X$ we define the {\em pullback tower}
$Y_\bullet=f^*(X_\bullet)$ as follows: $Y_0=Y$ and $Y_{i+1}$ is the
(saturated) blow up of $Y_i$ along the pullback of the center
$\calI_i$ of $X_{i+1}\to X_i$. Due to the universal property of
(saturated) blow ups this definition makes sense and $Y_\bullet$
is the universal (saturated) blow up tower of $Y$ that admits a
morphism to $X_\bullet$ extending $f$.

\begin{remark2}\label{8-towpullrem} Unlike pullbacks of
(normalized) blow up towers of schemes, see \S\S\ref{8-pullsec}
and \ref{8-normpullsec}, we do not distinguish strict transforms
and pullbacks. The above definition of pullback covers our needs,
and we do not have to study the base change of monoschemes. For the
sake of completeness, we note that fibred products of monoschemes
exist and in the affine case are defined by amalgamated sums of
monoids, see \cite{Deitmar@F1}. Also, it is easy to check that for
a smooth $f$ (which is the only case we will use) one indeed has
that $f^*(X_\bullet)=X_\bullet\times_XY$ for any (saturated) blow up
tower $X_\bullet$. For blow ups one checks this with charts and in
the saturated case one also uses that saturation is compatible
with a smooth morphism $f{\colon}Y\to X$, i.e. $Y^\sat\iso
X^\sat\times_XY$.  \end{remark2}

\subsubsection{Compatibility with sharpening}\label{8-compsharpsec}
Ideals and blow ups of fans are defined in the same way, but with
$\oM_X$ used instead of ${M}_X$ (Kato defines their saturated version
in \cite[9.7]{Toric@Kato}). Towers of blow ups of fans are defined
in the obvious way.

\begin{lemma2}\label{8-idealfanlem} Let $X=(X,{M}_X)$ be a
monoscheme, let $(F,M_F)=(X,\oM_X)$ be the corresponding fan and
let $\lambda{\colon}M_X\to M_F$ denote the sharpening homomorphism.

(i) $\calI\mapsto\lambda(\calI)$ induces a natural bijection between
the ideals on $X$ and on $F$.

(ii) Blow ups are compatible with sharpening, that is,
the sharpening of $\Bl_\calI(X))$ is naturally isomorphic to
$\Bl_{\lambda(\calI)}(F)$. The same statement holds for saturated
blow ups.

(iii) Sharpening induces a natural bijection between the (saturated)
blow up towers of $X$ and $F$.  \end{lemma2} \begin{proof} (i) is
obvious. (ii) is shown by comparing the blow up charts. Combining
(i) and (ii), we obtain (iii).  \end{proof}

\subsubsection{Desingularization}\label{8-fandesingsec} Using
the above notions of smoothness and blow ups, one can copy
other definitions of the desingularization theory. By a {\em
desingularization} (resp. {\em saturated desingularization}) of
a fine monoscheme $X$ we mean a blow up tower (resp. saturated
blow up tower) $X_n\dashto X_0=X$ with smooth $X_n$. By Remark
\ref{8-blowmonorem}, if $X$ admits a desingularization then it is
torsion free, and we will later see that the converse is also true.

For concreteness, we consider below non-saturated desingularizations,
but everything extends to the saturated case verbatim. A family
$\calF^\mono(X)$ of desingularizations of torsion free monoschemes
is called {\em functorial} (with respect to smooth morphisms) if for
any smooth $f{\colon}Y\to X$ the desingularization $\calF^\mono(Y)$
is the contracted pullback of $\calF^\mono(X)$. The same argument
as for schemes (see Remark \ref{8-functorrem}(ii)) shows that
$\calF^\mono$ is already determined by its restriction onto the
family of finite disjoint unions of affine monoschemes.

The definition of a functorial desingularization $\calF^\fan$ of fine
torsion free fans is similar. Since blow up towers and smoothness
are compatible with the sharpening functor, it follows that a
desingularization of a monoscheme $X$ induces a desingularization of
its sharpening. Moreover, any affine fan can be lifted to an affine
monoscheme and $\calF^\fan$ is determined by its restriction onto
disjoint unions of affine fans, hence we obtain the following result.

\begin{theorem2}\label{8-fandesingth} The sharpening functor
induces a natural bijection between functorial desingularizations
of fine torsion free monoschemes and functorial desingularizations
of fine torsion free fans. A similar statement holds for saturated
desingularizations.  \end{theorem2}

\begin{remark2}\label{8-satdesrem} Similarly to the normalization
of a desingularization tower, to any desingularization $\calF$ of
monoschemes or fans one can associate a saturated desingularization
$\calF^\sat$: one replaces all levels of the towers, except the zero
level, with their saturations. In this case blow ups are replaced
with saturated blow ups along the pulled back ideals. If $\calF$
is functorial with respect to all smooth morphisms then the same is
true for $\calF^\sat$. Indeed, for any smooth $Y\to X$ the centers
of $\calF(Y)$, $\calF^\sat(X)$ and $\calF^\sat(Y)$ are the pullbacks
of those of $\calF(X)$. In addition, the saturation construction
is compatible with the bijections from Theorem \ref{8-fandesingth}
in the obvious way.  \end{remark2}

\begin{remark2}\label{8-fandesrem} In principle, (saturated)
desingularization of fans or monoschemes can be described in
purely combinatorial terms of fans and their subdivisions (e.g. see
\cite[9.6]{Toric@Kato} or \cite[\S4]{Toric@Niziol}). However, it is
not easy to construct a functorial one directly. We will instead
use a relation between monoschemes and toric varieties to descend
desingularization of toric varieties to monoschemes and fans.
\end{remark2}

\subsection{Monoschemes and toric varieties}\label{8-monotoricsec}

\subsubsection{Base change from monoschemes to
schemes}\label{8-monoliftsec} Let $S$ be a scheme. The following
proposition introduces a functor from monoschemes to $S$-schemes
that can be intuitively viewed as a base change with respect to a
"morphism" $S\to\Spec(1)$.

\begin{proposition2}\label{8-monoliftprop} Let $S$ be a scheme
and $F$ be a monoscheme. Then there exists an $S$-scheme $X=S[F]$
with a morphism of monoidal spaces $f{\colon}(X,\calO_X)\to(F,M_F)$
such that any morphism $(Y,\calO_Y)\to(F,M_F)$, where $Y$ is an
$S$-scheme, factors uniquely through $f$.  \end{proposition2}
\begin{proof} Assume, first, that $F=\Spec(P)$ is affine. By Lemma
\ref{8-monoidsplem}(i), to give a morphism $(Y,\calO_Y)\to(F,M_F)$
is equivalent to give a homomorphism of monoids
$\phi{\colon}P\to\Gamma(\calO_Y)$, and the latter factors uniquely
through a homomorphism of sheaves of rings $\calO_S[P]\to\calO_Y$. It
follows that $S[F]=\mathbf{\Spec}(\calO_S[P])$ in this case. Since
the above formula is compatible with localizations by elements $a\in
P$, i.e. $S[F_a]\iso S[F]_{\phi(a)}$, it globalizes to the case of
an arbitrary monoscheme. Thus, for a general monoscheme $F$ covered
by $F_i=\Spec(P_i)$, the scheme $S[F]$ is glued from $S[F_i]$.
\end{proof}

\begin{remark2}\label{8-monoliftremrem} Note that if $S=\Spec(R)$
and $F=\Spec(P)$ then $S[F]=\Spec(R[P])$. However, we will often
consider an "intermediate" situation where $S=\Spec(R)$ is affine
and $F$ is a general monoscheme. To simplify notation, we will
abuse them by writing $R[F]$ instead of $\Spec(R[F])$. Such "mixed"
notation will always refer to a scheme.  \end{remark2}

\subsubsection{Toric schemes}\label{8-toricfansec} If $F$ is torsion
free and connected then we call $S[F]$ a {\em toric scheme} over
$S$. Recall that by Remark \ref{8-fanrem}(ii), $F$ possesses a unique
maximal point $\eta=\Spec(P^\gp)$, where $\Spec(P)$ is any affine
open submonoscheme. Hence $X=S[F]$ possesses a dense open subscheme
$T=S[\eta]$, which is a split torus over $S$, and the action of $T$
on itself naturally extends to the action of $T$ on $X$.

\begin{remark2} Assume that $k$ is a field. Classically, a toric
$k$-variety is defined as a normal finite type separated $k$-scheme
$X$ that contains a split torus $T$ as a dense open subscheme such
that the action of $T$ on itself extends to the whole $X$. If
$F$ is saturated then our definition above is equivalent to the
classical one. However, we also consider non-normal toric varieties
corresponding to non-saturated monoids.  \end{remark2}

\subsubsection{Canonical log structures}\label{8-canonsec}
For any monoscheme $F$, the $S$-scheme $X=S[F]$ possesses
a natural log structure induced by the universal morphism
$f{\colon}(X,\calO_X)\to(F,M_F)$. Namely, $M_X$ is the log structure
associated with the pre-log structure $g^{-1}M_F\to\calO_{X_\et}$,
where $g{\colon}(X_\et,\calO_{X_\et})\to(X,\calO_X)\to(F,M_F)$ is
the composition. We call $M_X$ the {\em canonical} log structure
of $X=S[F]$.

\begin{remark2}\label{8-canonrem} (i) The canonical log
structure is Zariski, as $\varepsilon_* M_X$ coincides with
the Zariski log structure associated with the pre-log structure
$f^{-1}M_F\to\calO_X$.

(ii) The log scheme $(X=S[F],M_X)$ is log smooth over the scheme
$S$ provided with the trivial log structure. In particular,
if $S$ is regular and $F$ is saturated then $(X,M_X)$ is fs and
log regular. Without the saturation assumption we still have that
$X^\sat\iso S[F^\sat]$ is log regular, hence $X$ is log regular
in the sense of Gabber (see \S\ref{8-nonsatursec}).

(iii) If $\eta$ is the set of generic points of $F$ then $T=S[\eta]$
is the open subset of $X$ which is the triviality set of its
log structure. However, the map $M_{X} \to \OO_X
\cap j_* \OO_{T}^\times$ is not an isomorphism in
general, as the case where $T = \Spec \mathbf{C}[t,t^{-1}]
\subset X = \Spec \mathbf{C}[t^2,t^3]$ already shows: the
image of $t^2 + t^3$ in $\OO_{X,\{0\}}$ belongs to $(
j_*\OO_{T}^\times)_{\{0\}}$, but does not belong to
$M_{\{0\}} = t^{P}\OO_{X,\{0\}}^\times$, where $P$ is the
(fine, but not saturated) submonoid of $\mathbf{N}$ generated by 2
and 3. (We use that $\OO_{X,\{0\}}$ is strictly smaller than
its normalization $\mathbf{C}[t]_{(t)}=\OO_{X,\{0\}}[t]$
and hence $\frac{t^2+t^3}{t^2}=1+t$ is not contained in
$\OO_{X,\{0\}}$.)  \end{remark2}

\subsubsection{Toric saturation}\label{8-toricsatsec} Saturation of
monoschemes corresponds to normalization of schemes. This will play
an essential role later, since we get a combinatorial description
of the normalization.

\begin{lemma2}\label{8-toricsatlem} If $S$ is a normal scheme
and $F$ is a fine monoscheme then there is a natural isomorphism
$S[F]^\nor\iso S[F^\sat]$.  \end{lemma2} \begin{proof} Note
that $f{\colon}\ZZ[F^\sat]\to\Spec(\ZZ)$ is a flat morphism and
its fibers are normal because  they are classical toric varieties
$\FF_p[F^\sat]$. So, $f\times S{\colon}S[F^\sat]\to S$ is a flat
morphism with normal fibers and normal target, and we obtain that
its source is normal by \cite[21.E(iii)]{CA@Matsumura}. It remains
to note that $S[F^\sat]\to S[F]$ is a finite morphism inducing
isomorphism of dense open subschemes $S[F^\gp]$, hence $S[F^\sat]$
is the normalization of $S[F]$.  \end{proof}

\begin{remark2} The same argument shows that if $S$ is Cohen-Macaulay
then so is $S[F^\sat]$.  \end{remark2}

\subsubsection{Toric smoothness}\label{8-toricsmoothsec} Next,
let us compare smoothness of morphisms of monoschemes as defined
in \S\ref{8-smoothfansec} and classical smoothness of toric
morphisms. The following lemma slightly extends the classical result
(e.g. see \cite[\S2.1]{Toric@Fulton}) that if $P$ is fs and $\CC[P]$
is regular then $P\iso\NN^r\oplus\ZZ^s$.

\begin{lemma2}\label{8-smoothfanlem} Let $f{\colon}F\to F'$ be a
morphism of fine monoschemes and let $S$ be a scheme.

(i) If $f$ is smooth then $S[f]$ is smooth (as a morphism of
schemes).

(ii) If $F$ is torsion free and the morphism $S[F]\to S$ is
smooth then $F$ is smooth.  \end{lemma2} \begin{proof} Part
(i) is obvious, so let us check (ii). We can also assume that
$F=\Spec(P)$ is affine. Also, we can replace $S$ with any of its
points achieving that $S=\Spec(k)$. Then $P$ is a fine torsion
free monoid and $k[P]\subseteq k[P^\gp]\iso k[\ZZ^n]$. It
follows that $X=\Spec(k[P])$ is an integral smooth $k$-variety
of dimension $n$. Note that $\Spec(k[P^\sat])$ is a finite
modification of $X$ which is generically an isomorphism. Since $X$
is normal we have that $\Spec(k[P^\sat])\iso X$, and it follows
that $P$ is saturated. Now, $P\iso\oP\oplus\ZZ^l$ and hence
$X\iso\Spec(k[\oP])\times_k\GG_m^l$. Obviously, $\Spec(k[\oP])$
is smooth of dimension $r=n-l$ and our task reduces to showing
that $\oP\iso\NN^{r}$.

Let $m=\oP\setminus\{1\}$ be the maximal ideal of $\oP$. Then
$I=k[m]$ is a maximal ideal of $k[\oP]$ with residue field
$k$. In particular, by $k$-smoothness of $k[\oP]$ we have that
$\dim_k(I/I^2)=r$. On the other hand, $I=\oplus_{x\in m}xk$
and $I^2=\oplus_{x\in m^2}xk$, hence $I/I^2\iso\oplus_{x\in
m\setminus m^2}xk$ and we obtain that $m\setminus m^2$ consists of
$r$ elements $t_1\.t_r$. Note that these elements generate $\oP$
as a monoid and hence they generate $\oP^\gp$ as a group. Since
$\oP^\gp\iso\ZZ^r$, the elements $t_1\.t_r$ are linearly
independent in $\oP^\gp$, and we obtain that the surjection
$\oplus_{i=1}^rt_i^\NN\to\oP$ is an isomorphism.  \end{proof}

\begin{remark2} It seems very probable that, much more
generally, $f$ is smooth whenever $S[f]$ is smooth as a
morphism of schemes and one of the following conditions holds:
(a) $S$ has points in all characteristics, (b) the homomorphisms
${M}^\gp_{F',x'}\to{M}^\gp_{F,x}$ induced by $f$ have torsion free
kernels and cokernels. We could prove this either in the saturated
case or under some milder but unnatural restrictions. The main ideas
are similar but the proof becomes more technical. We do not develop
this direction here since the lemma covers our needs.  \end{remark2}

\subsubsection{Toric ideals}\label{8-toricidealsec} Let $F$ be
a torsion free monoscheme, $k$ be a field and $X=k[F]$. For any
ideal $\calI$ on $F$ one naturally defines an ideal $k[\calI]$
on $X$: in local charts, an ideal $I\subseteq P$ goes to the
ideal $Ik[P]=k[I]=\oplus_{a\in I}ak$ in $k[P]$. We say that
$\calJ=k[\calI]$ is a {\em monoidal ideal} on $k[F]$. Note that
$\calJ$ determines $\calI$ uniquely and, actually, $\calI=\calJ\cap
M_X$ (this is obvious in local charts).

\begin{lemma2}\label{8-toricideallem} Assume that $F$ is connected,
$\eta$ is its maximal point, $X=k[F]$, and $T=k[\eta]$ is the torus
of the toric scheme $X$. A coherent ideal $\calJ\subseteq{M}_X$
is $T$-equivariant if and only if it is monoidal.  \end{lemma2}
\begin{proof} Any monoidal ideal is obviously $T$-equivariant,
so let us prove the inverse implication. The claim is local, so
we should prove that any $T$-equivariant ideal $J\subset A=k[P]$
is of the form $k[I]$ for an ideal $I$ of $P$. Consider the
coaction homomorphism $\mu{\colon}A\to A\otimes_k k[P^\gp]=B$. The
equivariance of $J$ means that $JB$ (with respect to the embedding
$A\into A\otimes_k k[P^\gp]$) is equal to $\mu(J)B$. In particular,
$\mu|_J{\colon}J\to JB=J\otimes_k k[P^\gp]$ induces a $P^\gp$-grading
on $J$ compatible with the $P^\gp$-grading $A=\oplus_{\gamma\in P}
A_\gamma$. Thus, $J$ is a homogeneous ideal in $A$ and, since $A_1=k$
is a field and each $k$-module $A_\gamma$ is of rank one, we obtain
that $J=\oplus_{\gamma\in I} A_\gamma$ for a subset $I\subseteq
P$. Thus, $J=k[I]$, and clearly $I$ is an ideal.  \end{proof}

\subsubsection{Toric blow ups}\label{8-toricblowsec} We will also
need the well known fact that toric blow ups are of combinatorial
origin, i.e. they are induced from blow ups of monoschemes.

\begin{lemma2}\label{8-blowfanlem} Assume that $F$ is a torsion
free monoscheme, $\calI\subseteq{M}_F$ is an ideal, $X=k[F]$,
and $\calJ=k[\calI]$. Then there is a canonical isomorphism
$\Bl_\calJ(X)\iso k[\Bl_\calI(F)]$.  \end{lemma2} \begin{proof}
Assume first that $F=\Spec(P)$. Then $\calI$ corresponds to an ideal
$I\subseteq P$ and we can simply compare charts: $\Bl_\calI(F)$
is covered by the charts $\Spec(P[a^{-1}I])$ for $a\in I$, and,
since $I$ generates $J$, the charts $k[a^{-1}J]=k[P[a^{-1}I]]$
cover $\Bl_J(X)$. This construction is compatible with localizations
$(P,I)\mapsto (P_b,b^{-1}I)$ hence it globalizes to the case of a
general fine monoscheme with an ideal.  \end{proof}

Using Lemma \ref{8-toricsatlem} we obtain a similar relation between
saturated blow ups and normalized toric blow ups.

\begin{corollary2}\label{8-blowfancor} Keep notation of
Lemma \ref{8-blowfanlem}. Then $\Bl_\calJ(X)^\nor\iso
k[\Bl_\calI(F)^\sat]$.  \end{corollary2}

\subsubsection{Desingularization of
monoschemes}\label{8-desingmonosec} Let $F$ be a torsion free
monoscheme and $X=k[F]$ for a field $k$ of characteristic zero
(e.g. $k=\QQ$). Recall that the normalized desingularization functor
$\tilcalF$ from \S\ref{8-desingQsec} is compatible with the action
of any smooth $k$-group, hence the centers of $\tilcalF(X)\colon
X_n\dashto X_0=X$ are $T$-equivariant ideals. By Lemma
\ref{8-toricideallem}, the blown up ideal of $X_0$ is of the form
$k[\calI]$ for an ideal $\calI\subseteq{M}_F$, hence $X_1=k[F_1]$
for $F_1=\Bl_\calI(F)^\sat$ by Lemma \ref{8-blowfancor}. Applying
this argument inductively we obtain that the entire normalized
blow up tower $\tilcalF(X)$ descends to a saturated blow up tower
of $F$, which we denote as $\tilcalF^\mono(F)$ (in other words,
$\tilcalF(X)=k[\tilcalF^\mono(F)]$). Since $F_n$ is smooth by
Lemma \ref{8-smoothfanlem}(ii), the tower $\tilcalF^\mono(F)$
is a desingularization of $F$. Moreover, part (i) of the same
lemma implies that $\tilcalF^\mono$ is functorial with respect to
smooth morphisms of monoschemes. Namely, for any smooth morphism
$F'\to F$, $\tilcalF^\mono(F')$ is the contracted pullback of
$\tilcalF^\mono(F)$ (see \S\ref{8-towpullsec}). Summarizing, we
have obtained:

\begin{theorem2}\label{8-monodesinth} Let $k$ be a field and let
$\tilcalF$ be a normalized desingularization of disjoint unions
of toric $k$-varieties which is functorial with respect to smooth
morphisms. Then each normalized blow up tower $\tilcalF(k[F])$
is the pullback of a uniquely defined saturated blow up tower
of monoschemes $\tilcalF^\mono(F)$. This construction produces a
saturated desingularization of monoschemes which is functorial with
respect to smooth morphisms.  \end{theorem2}

Combining theorems \ref{8-fandesingth} and \ref{8-monodesinth}
we obtain a functorial saturated desingularization of fans that
will be denoted $\tilcalF^\fan$.

\begin{remark2} We will work with normalized and saturated
desingularizations, so we formulated the theorem for $\tilcalF$. The
same argument shows that $\calF$ induces desingularizations
$\calF^\mono$ and $\calF^\fan$ that are functorial with respect
to all smooth morphisms. Moreover, the descent from toric
desingularization is compatible with normalization/saturation,
i.e. $\tilcalF^\mono=(\calF^\mono)^\sat$ and similarly for fans.
\end{remark2}

\subsection{Monoidal desingularization}\label{8-logfansec} In this
section we will establish, what we call, monoidal desingularization
of fine log schemes $(X,M_X)$. This operation "resolves" the sheaf
$\oM_X$ but does not "improve" the log strata of $X$.

\subsubsection{Log stratification}\label{8-logstratsec}
Using charts one immediately checks that for any fine log scheme $(X,M_X)$
the function $x\mapsto {\rm rank}(\oM^\gp_\ox)$ is constructible.
The corresponding stratification of $X$, whose strata are the maximal locally
closed subsets on which this function is constant, will be called
the {\em log stratification} (the analogous stratification in \refext{VI}{6-1.5} was called
canonical or stratification by rank).

\subsubsection{Monoscheme charts of log
schemes}\label{8-monochartsec} A (global) {\em monoscheme
chart} of a Zariski log scheme $(X,M_X)$ is a morphism of
monoidal spaces $c{\colon}(X,\varepsilon_*M_X)\to(F,{M}_F)$
such that the target is a monoscheme and $\varepsilon_*M_X$
is isomorphic to the Zariski log structure associated with the
pre-log structure $c^{-1}M_F\to\calO_X$ (obtained as $c^{-1}M_F\to
M_X\to\calO_X$). In particular, $M_X$ is the log structure associated
with $(c\circ\varepsilon)^{-1}M_F\to\calO_{X_\et}$. We say that
the chart is {\em fine} if $(F,M_F)$ is so. For example, any toric
scheme $R[F]$, where $F$ is a monoscheme, possesses a canonical
chart $R[F]\to F$.

\begin{lemma2}\label{8-monochartlem} Let $(X,M_X)$ be a log scheme
and let $(F,M_F)$ be a monoscheme. Then any morphism of monoidal
spaces $f{\colon}(X,M_X)\to(F,M_F)$ factors uniquely into the
composition of a morphism of log schemes $(X,M_X)\to\ZZ[F]$ and
the canonical chart $\ZZ[F]\to F$.  \end{lemma2} \begin{proof}
Note that $(Id_X,\alpha){\colon}(X,\calO_X)\to(X,M_X)$ is
a morphism of monoidal spaces, hence so is the composition
$h{\colon}(X,\calO_X)\to(F,M_F)$. If $F=\Spec(P)$ then $h$ is
determined by the homomorphism $P\to\Gamma(\calO_X)$ by Lemma
\ref{8-monoidsplem}. Since the latter factors uniquely into the
composition of the homomorphism of monoids $P\to\ZZ[P]$ and the
homomorphism of rings $\ZZ[P]\to\Gamma(\calO_X)$, we obtain a
canonical factoring $X\to\Spec(\ZZ[P])\to\Spec(P)$. Furthermore,
this affine construction is compatible with localizations of $P$,
hence it globalizes to the case when the monoscheme $F$ is arbitrary.
\end{proof}

\begin{remark2}\label{8-usualchartrem} (i) Usually, one works with
log schemes using local charts $(X,M_X)\to\Spec(\ZZ[P])$. By Lemma
\ref{8-monochartlem} this is equivalent to working with affine
monoscheme charts.

(ii) In particular, any fine log scheme $(X,M_X)$ admits a fine
monoscheme chart \'etale-locally, i.e., there exists a strict (in
the log sense) \'etale covering  $(Y,M_Y)\to(X,M_X)$ whose source
possesses a fine monoscheme chart. Similarly, any Zariski fine log scheme
admits a fine monoscheme chart Zariski locally. \end{remark2}

\subsubsection{Chart base change}\label{8-chartchangesec}
Given a fine monoscheme chart $(X,M_X)\to F$ and a morphism of
monoschemes $F'\to F$ we will write $(X,M_X)\times_FF'$ instead
of $(X,M_X)\times_{\ZZ[F]}\ZZ[F']$, where the second product
is taken in the category of fine log schemes. This notation is
partially justified by the following result.

\begin{lemma2}\label{8-chartbaselem} Keep the above notation and
let $(X',M_{X'})=(X,M_X)\times_FF'$.

(i) The morphism $c'{\colon}(X',M_{X'})\to F'$ is a monoscheme chart.

(ii) If $(Y,M_Y)$ is a log scheme over $(X,M_X)$ and
$d{\colon}(Y,M_Y)\to F$ is the induced morphism of monoidal spaces,
then any lifting of $d$ to a morphism $(Y,M_Y)\to F'$ factors
uniquely through $c'$.  \end{lemma2} \begin{proof} Strictness is
stable under base changes, hence $(X',M_{X'})\to\ZZ[F']$ is strict
and we obtain (i). The assertion of (ii) is a consequence of Lemma
\ref{8-monochartlem}.  \end{proof}

\subsubsection{Log ideals}\label{8-logidealsec} By {\em log
ideal} on a fine log scheme $(X,M_X)$ we mean any ideal
$\calI\subseteq M_X$ that \'etale-locally on $X$ admits
a coherent chart as follows: there exists a strict \'etale
covering $f{\colon}(Y,M_Y)\to(X,M_X)$, a fine monoscheme chart
$c{\colon}(Y,M_Y)\to F$ and a coherent ideal $\calI_F\subseteq M_F$
such that $f^{-1}(\calI)M_Y=c^{-1}(\calI_F)M_Y$.

\subsubsection{Log blow ups}\label{8-logblowsec} It is proved
in \cite[4.2]{Toric@Niziol} that there exists a universal
morphism $f{\colon}(X',M_{X'})\to(X,M_X)$ such that the ideal
$f^{-1}(\calI)M_{X'}$ is invertible, i.e. locally (in the \'etale
topology) generated by one element. (We use here that, unlike rings,
any principal ideal $aM$ of an integral monoid $M$ is invertible
in the usual sense, i.e. $M\iso aM$ as $M$-sets.) Actually,
the formulation in \cite{Toric@Niziol} refers only to saturated
blow ups, but the proof deals also with the non-saturated ones.

The construction of log blow us is standard and it also shows
that they are compatible with arbitrary strict morphisms. If
$(X,M_X)$ and $\calI$ admit a chart $F,\calJ\subseteq M_F$ then
$(X',M_{X'})=(X,M_X)\times_F\Bl_\calJ(F')$ is as required. If
$(Y,M_Y)\to(X,M_X)$ is strict then $F$ is also a chart of $(Y,M_Y)$,
hence the local construction is compatible with strict morphisms. The
general case now follows by descent because any fine log scheme
admits a chart  \'etale locally. We call $f$ the {\em log blow up}
of $(X,M_X)$ along $\calI$ and denote it $\LogBl_{\calI}(X,M_X)$
(it is called unsaturated log blow up in \cite{Toric@Niziol}). Log
blow up towers are defined obviously. As usually, contraction
of such a tower is obtained by removing all empty log blow ups
(i.e. blow ups along $\calI=M_X$).

\subsubsection{Saturated log blow ups}\label{8-satlogblowsec}
Saturated log blow up along a log ideal $\calI$ is defined
as $(\LogBl_{\calI}(X,M_X))^\sat$. It satisfies an obvious
universal property too. (It is called log blow up in
\cite{Toric@Niziol}). Towers of saturated log blow ups, their
pullbacks, and saturation of a tower of log blow ups are defined
in the obvious way.

\subsubsection{Pullbacks}\label{8-logpullsec} Let
$f{\colon}(Y,M_Y\to(X,M_X)$ be a morphism of log schemes. By
{\em pullback} of the log blow up $\LogBl_\calI(X,M_X)$
along a log ideal $\calI\subseteq M_X$ we mean the log blow up
$\LogBl_{\calJ}(Y,M_Y)$, where $\calJ=f^{-1}(\calI)M_Y$. This is
the universal log scheme over $(Y,M_Y)$ whose morphism to $(X,M_X)$
factors through $\LogBl_\calI(X,M_X)$. The pullback of saturated blow
ups is defined similarly, and these definitions extend inductively
to pullbacks of towers of (saturated) log blow ups.

\subsubsection{Basic properties}\label{8-basiclogblowsec}
Despite the similarity with usual blow ups of schemes,
log blow ups (resp. saturated log blow ups) have nice
properties that are not satisfied by usual blow ups. First,
it is proved in \cite[4.8]{Toric@Niziol} that log blow ups
are compatible with any log base change $f{\colon}Y\to X$,
i.e. $\LogBl_{f^{-1}\calI}(Y)\iso\LogBl_\calI(X)\times_XY$
for a monoidal ideal $\calI$ on $X$. In particular, saturated
blow ups are compatible with saturated base changes. Second, log
blow ups (resp. saturated log blow ups) are log \'etale morphisms
because so are both saturation morphisms and charts of the form
$\ZZ[\Bl_\calI(F)]\to\ZZ[F]$.

\subsubsection{Fan charts}\label{8-fanchart} A {\em
fan chart} of a Zariski log scheme $(X,M_X)$ is a
morphism $d{\colon}(X,\varepsilon_*\oM_X)\to(F,M_F)$
of monoidal spaces such that the target is a fan and
$d^{-1}(M_F)\iso\varepsilon_*\oM_X$. For example, for any
monoscheme chart $c{\colon}(X,\varepsilon_*M_X)\to(F,M_F)$,
its sharpening $\oc{\colon}(X,\varepsilon_*\oM_X)\to(F,\oM_F)$
is a fan chart. Fan charts were considered by Kato (e.g.,
in \cite[9.9]{Toric@Kato}). They contain less information than
monoscheme charts, but "remember everything about ideals and blow
ups" because there is a one-to-one correspondence between ideals and
blow up towers of $M_F$ and $M_\oF$. Let us make this observation
rigorous. For concreteness, we discuss only non-saturated (log)
blow ups, but everything easily extends to the saturated case.

\begin{remark2}\label{8-fanchartrem} (i) Assume that
$\oc{\colon}(X,\oM_X)\to(F,M_F)$ is a fan chart. Any ideal
$\calI_F\subseteq M_F$ induces a log ideal $\calI\subset
M_X$, which is the preimage of $\oc^{-1}(\calI_F)\oM_X$
under $M_X\to\oM_X$. We say that the blow up $F'=\Bl(\calI_F)$
induces the log blow up $(X',\oM_{X'})=\LogBl_\calI(X,\oM_X)$
or that the latter log blow up is the {\em pullback} of
$\Bl(\calI_F)$. Furthermore, $(X',\oM_{X'})\to F'$ is also a fan
chart (see \cite[9.9]{Toric@Kato}, where the fs case is treated),
hence this definition iterates to a tower $F_\bullet$ of blow ups
of $F$. We will denote the pullback tower of log blow ups as
$\oc^*(F_\bullet)$.

(ii) By a slight abuse of language, Kato and Niziol denote
$\oc^*(F_\bullet)$ as $(X,M_X)\times_FF_\bullet$. One should be very
careful with this notation because, in general, there is no morphism
$(X,M_X)\to F$ that lifts $\oc$. Also, one cannot define analogous
"base change" for an arbitrary morphism of fans $F'\to F$. The reason
is that there are many "unnatural" gluings in the category of fans
(e.g. along generic points), and such gluings cannot be lifted to
log schemes (and even to monoschemes).

(iii) For blow up towers, however, the base change
notation is safe and agrees with the base change from
the monoschemes. Namely, if $\oc$ is the sharpening of a
monoscheme chart $c{\colon}(X,M_X)\to(F,M_F)$ then there exists
a one-to-one correspondence between blow up towers of the
monoscheme $F=(F,M_F)$ and the fan $\oF=(F,\oM_F)$, see Lemma
\ref{8-idealfanlem}. Clearly, the matching towers induce the same
log blow up tower of $(X,M_X)$. In particular, $\calF^\mono(F,M_F)$
and $\calF^\fan(F,\oM_F)$ (see Theorem \S\ref{8-fandesingth})
induce the same log blow up tower of $(X,M_X)$.  \end{remark2}

\subsubsection{Monoidal desingularization of log
schemes}\label{8-monoidaldessec} Let $(X,M)$ be a fine log
scheme and assume that $(X,M)$ is {\em monoidally torsion free}
in the sense that the groups $\oM^\gp_\ox$ are torsion free. By
{\em monoidal desingularization} (resp. {\em saturated monoidal
desingularization}) of a fine log scheme $(X,M)$ we mean a
tower of log blow ups (resp. tower of saturated log blow ups)
$(X_n,M_n)\dashto (X_0,M_0)=(X,M)$ such that for any geometric
point $\ox\to X_n$ the stalk of $\oM_n$ at $\ox$ is a free monoid. A
morphism $(Y,N)\to(X,M)$ is called {\em monoidally smooth} if each
induced homomorphism of stalks of monoids $\oM_\ox\to\oN_\oy$ can
be extended to an isomorphism $\oM_\ox\oplus\NN^r\iso\oN_\oy$.

\begin{theorem2}\label{8-monoidaldesth} Let $\tilcalF^\fan$ be a
saturated desingularization of fine torsion free fans which is
functorial with respect to smooth morphisms. Then there exists
unique saturated monoidal desingularization $\tilcalF^\rmlog(X,M)$
of monoidally torsion free fine log schemes $(X,M)$, such that
$\tilcalF^\rmlog$ is functorial with respect to all monoidally
smooth morphisms and $\tilcalF^\rmlog(X,M)$ is the contraction of
$c^*(\tilcalF^\fan(F))$ for any log scheme $(X,M)$ that admits a
fan chart $c{\colon}(X,\oM_X)\to F$. In the same way, a functorial
desingularization $\calF^\fan$ induces a monoidal desingularization
$\calF^\rmlog$.  \end{theorem2}

\begin{remark2}\label{8-monoidaldesrem} Since any monoscheme
chart induces a fan chart, it then follows from Remark
\ref{8-fanchartrem}(iii) that $\tilcalF^\rmlog(X,M)$ is the
contraction of $d^*(\tilcalF^\mono(F))$ for any log scheme
$(X,M)$ that admits a monoscheme chart $d{\colon}(X,\oM_X)\to F$.
\end{remark2}

\begin{proof}[Proof of Theorem \ref{8-monoidaldesth}] Both
cases are established similarly, so we prefer to deal with
$\calF^\fan$ (to avoid mentioning saturations at any step of
the proof). By descent, it suffices to show that the pullback
from fans induces a functorial monoidal desingularization of
those fine log schemes that admit a global fan chart. Thus,
if $\calF^\rmlog$ exists then it is unique, and our aim is to
establish existence and functoriality. Both are consequences of
the following claim: assume that $f{\colon}(Y,M_Y)\to(X,M_X)$
is a monoidally smooth morphism whose source and target admit fan
charts $d{\colon}(Y,\oM_Y)\to G$ and $d'{\colon}(X,\oM_X)\to F$,
then contractions of $d^*(\calF^\fan(G))$ and $c^*(\calF^\fan(F))$
are equal, where $c=d'\circ\of{\colon}(Y,\oM_Y)\to(X,\oM_X)\to
F$. Note that $d^{-1}(\oM_G)\iso\oM_Y$ and the homomorphism
$c^{-1}(\oM_F)\iso\oM_Y$ is smooth.

Choose a point $y\in Y$ and consider the localizations
$Y'=\Spec(\calO_{Y,y})$, $F'=\Spec(M_{F,c(y)})$ and
$G'=\Spec(M_{G,d(y)})$ at $y$ and its images in the fans. Since
$\phi{\colon}M_{G,d(y)}\to\oM_{Y,y}$ is an isomorphism and the
homomorphism $\psi{\colon}M_{F,c(y)}\to\oM_{Y,y}$ is smooth,
we obtain a factorization of $\psi$ into a composition of
a homomorphism $\lambda{\colon}M_{F,c(y)}\to M_{G,d(y)}$
and $\phi$, where $\lambda$ is smooth. Set $U=c^{-1}(F')\cap
d^{-1}(G')$. Then $U$ is a neighborhood of $y$, and $c$ and $d$
induce homomorphisms $\phi_U{\colon}M_{G,d(y)}\to\oM_Y(U)$ and
$\psi_U{\colon}M_{F,c(y)}\to\oM_Y(U)$. Sine the monoids are fine,
we can shrink $U$ so that the equality $\psi_U=\phi_U\circ\lambda$
holds. It then follows from Lemma \ref{8-monoidsplem}(ii) that
$c|_U$ factors into a composition of $d|_U$ and the smooth morphism
$\Spec(\lambda){\colon}G'\to F'$.

By quasi-compactness of $Y$ we can now find finite coverings
$Y=\cup_{i=1}^nY_i$, $F=\cup_{i=1}^nF_i$ and $G=\cup_{i=1}^nG_i$,
and smooth morphisms $\lambda_i{\colon}G_i\to F_i$ such that
$Y_i$ is mapped to $F_i$ and $G_i$ by $c$ and $d$, respectively,
and the induced maps of monoidal spaces $c_i{\colon}Y_i\to
F_i$ and $d_i{\colon}Y_i\to G_i$ satisfy $c_i=\lambda_i\circ
d_i$. Set $Y'=\coprod_{i=1}^nY_i$, $F'=\coprod_{i=1}^n F_i$,
$G'=\coprod_{i=1}^nG_i$, $c'{\colon}Y'\to F$ and $d'{\colon}Y'\to
G$. By descent, it suffices to check that contractions of
$c'^*(\calF^\fan(F))$ and $d'^*(\calF^\fan(G))$ are equal. Since,
the morphism $Y'\to F$ factors through the surjective smooth morphism
$F'\to F$, and similarly for $G$, these two pullbacks are equal to
the contracted pullbacks of $\calF^\fan(F')$ and $\calF^\fan(G')$,
respectively. It remains to note that $Y'\to F'$ factors through the
smooth morphism $\coprod_{i=1}^n\lambda_i{\colon}G'\to F'$. Hence
$\calF^\fan(G')$ is the contracted pullback of $\calF^\fan(F')$,
and their contracted pullbacks to $Y'$ coincide.  \end{proof}

\subsection{Desingularization of log regular log
schemes}\label{8-resolution of log schemes1} In this section we will
see how saturated monoidal desingularization leads to normalized
desingularization of log regular log schemes. Up to now we freely
considered saturated and unsaturated cases simultaneously, and did
not feel any essential difference. This will not be the case in the
present section because the notion of log regularity was developed
by Kato and Niziol in the saturated case. Gabber generalized the
definition for non-saturated case and extended to that case all
main results about log regular log schemes. This was necessary for
his original approach, but can be by passed by use of saturated
monoidal desingularization. So, we prefer to stick to the saturated
case and simply refer to all foundational results about log regular
fs log schemes to \cite{Toric@Kato} and \cite{Toric@Niziol}. For
the sake of completeness, we will outline Gabber's results about
the general case in \S\ref{8-nonsatursec}.

\subsubsection{Conventions}\label{8-logregconvsec} Recall
that Kato's notion of log regular fs log schemes was already
used in expos\'e VI, \S1.2. Throughout \S\ref{8-resolution of
log schemes1} we assume that $(X,M_X)$ is a log regular fs log
scheme. Note that the homomorphism $\alpha_X{\colon}M_X\to\calO_X$
of $X_\et$-sheaves is injective by \cite[2.6]{Toric@Niziol}, and
actually $M_X=\calO^\times_U\cap\calO_X$, where $U\subseteq X$ is
the triviality locus of $M_X$. So, we will freely identify $M_X$
with a multiplicative submonoid of $\calO_X$.

\subsubsection{Monoidal ideals}\label{8-logregidealsec}
For any log ideal $\calI\subseteq M_X$ consider the ideal
$\calJ=\alpha(\calI)\calO_X$ it generates. We call $\calJ$ a {\em
monoidal ideal} and by a slight abuse of language, we will write
$\calJ=\calI\calO_X$.

\begin{lemma2}\label{8-logregideallem} Let $X=(X,M_X)$ be as in
\S\ref{8-logregconvsec}. The rules $\calI\mapsto\calI\calO_X$
and $\calK\mapsto\calK\cap M_X$ give rise to a one-to-one
correspondence between log ideals $\calI\subseteq M_X$
and monoidal ideals $\calK\subseteq\calO_X$.  \end{lemma2}
\begin{proof} It suffices to show that any log ideal $\calI$
coincides with $\calJ=\calI\calO_X\cap M_X$. The question easily
reduces to the case when $X$ admits a chart $X\to\Spec(\ZZ[P])$
and $\calI=IM_X$, $\calJ=JM_X$ for ideals $I\subseteq J$
of $P$. Assume that the lemma fails, i.e. $I\subsetneq
J$, and consider the exact sequence \begin{equation}
\label{8-logregeq} I\ZZ[P]\otimes_{\ZZ[P]}\calO_X\to
J\ZZ[P]\otimes_{\ZZ[P]}\calO_X\to
J\ZZ[P]/I\ZZ[P]\otimes_{\ZZ[P]}\calO_X\to 0 \end{equation}
Since $\Tor_1^{\ZZ[P]}(\ZZ[P]/I\ZZ[P],\calO_X)=0$
by \cite[6.1(ii)]{Toric@Kato},
$I\ZZ[P]\otimes_{\ZZ[P]}\calO_X=I\calO_X$ and similarly
for $J$, and we obtain that the first morphism in the
sequence (\ref{8-logregeq}) is $I\calO_X\to J\calO_X$. To
obtain a contradiction, it suffices to show that
$J\ZZ[P]/I\ZZ[P]\otimes_{\ZZ[P]}\calO_X\neq 0$. Note that
$\ZZ[P]/m_P\ZZ[P]$ is a quotient of $J\ZZ[P]/\ZZ[I]$,
so it remains to note that $m_P\calO_X\neq \calO_X$ and hence
$\ZZ[P]/m_p\ZZ[P]\otimes_{\ZZ[P]}\calO_X\neq 0$.  \end{proof}

\subsubsection{Interpretation of monoidal smoothness}\label{8-regsec}
Note that by \refext{VI}{6-1.7} $(X,M_X)$ is monoidally smooth if and only if
$X$ is regular and in this case the non-triviality locus of $M_X$
is a normal crossings divisor $D$.

\subsubsection{Saturated log blow ups of log regular
log schemes}\label{8-logregsatblowsec} Using Kato's
$\Tor$-independence result \cite[6.1(ii)]{Toric@Kato} Niziol
proved in \cite[4.3]{Toric@Niziol} that saturated log blow
ups of $(X,M_X)$ are compatible with normalized blow ups along
monoidal ideals. Namely, if $(Y,M_Y)=\LogBl_\calI(X,M_X)^\sat$
then $Y\iso\Bl_{\calI\calO_X}(X)^\nor$. We will also need more
specific results that showed up in the proof of loc.cit., so we
collect them altogether in the following lemma.

\begin{lemma2}\label{8-logregblowlem} Let
$f{\colon}(X,M_X)\to(Y,M_Y)$ be a strict morphism of fs log
regular log schemes, let $\calI\subseteq M_Y$ be a log ideal and
$\calJ=f^{-1}\calI M_X$. Set $(X',M_{X'})=\LogBl_\calJ(X,M_X)$,
$(X'',M_{X''})=(X',M_{X'})^\sat$, $(Y',M_{Y'})=\LogBl_\calI(Y,M_Y)$
and $(Y'',M_{Y''})=(Y',M_{Y'})^\sat$. Then

(i) The (saturated) log blow up of $(X,M_X)$ is compatible with
(normalized) blow up of $X$: $X'\iso\Bl_{\calI\calO_X}(X)$
and $X''\iso\Bl_{\calI\calO_X}(X)^\nor$.

(ii) The (normalized) blow up of $X$ along $\calJ$ is the pullback
of the (normalized) blow up of $Y$ along $\calI$. In particular,
$X'\iso X\times_YY'$ and $X''\iso X\times_YY''$.  \end{lemma2}
\begin{proof} All claims can be checked \'etale locally, hence we can
assume that there exists a chart $g{\colon}(Y,M_Y)\to(\ZZ[P],P)$
and $\calI=g^{-1}(I_0)M_X$ for an ideal $I_0\subseteq P$. Then it
suffices to prove (ii) for $g$ and the induced chart $g\circ f$ of
$(X,M_X)$. In particular, this reduces the lemma to the particular
case when $X=\Spec(A)$ and $f$ is a chart $(X,M_X)\to(\ZZ[P],P)$. It
is shown in the first part of the proof of \cite[4.3]{Toric@Niziol}
that $$X'\iso\Proj(A\otimes_{\ZZ[P]}(\oplus_{n=0}^\infty
I_0^n))\iso\Proj(\oplus_{n=0}^\infty\calJ^n)$$ The first
isomorphism implies that $X'\to X$ is the base change of
$\Proj(\oplus_{n=0}^\infty I_0^n)=Y'\to Y$, and the second
isomorphism means that $X'\iso\Bl_{\calI\calO_X}(X)$. This
establishes the unsaturated and unnormalized part of the Lemma,
and the second part follows in the same way from the second part
of the proof of \cite[4.3]{Toric@Niziol}.  \end{proof}

\begin{remark2} It follows from the lemma that the unsaturated log
blow up $(X',M_{X'})$ is log regular in the sense of Gabber, see
\S\ref{8-nonsatursec}. Thus, once log regularity is correctly defined
in full generality, it becomes a property preserved by log blow ups
(as it should be, since log blow ups are log smooth).  \end{remark2}

\subsubsection{Desingularization of log regular log
schemes}\label{8-desingloregsec} By Lemmas \ref{8-logregideallem}
and \ref{8-logregblowlem}, any saturated log blow up tower
$(X_n,M_n)\dashto (X,M_X)$ induces a normalized blow up tower
$X_n\dashto X$, which completely determines it. Furthermore,
by \S\ref{8-regsec}, the first tower is a saturated monoidal
desingularization if and only if the second tower is a normalized
desingularization. In particular, the saturated monoidal
desingularization $\tilcalF^\rmlog(X,M_X)$ induces a normalized
desingularization of the scheme $X$ that depends on $(X,M_X)$
and will be denoted $\tilcalF(X,M_X)$.

\begin{theorem2}\label{8-deslogregth} The saturated
monoidal desingularization $\tilcalF^\rmlog$ (see Theorem
\ref{8-monoidaldesth}) gives rise to desingularization
$\tilcalF$ of log regular log schemes that possesses the same
functoriality properties: if $\phi{\colon}(Y,M_Y)\to(X,M_X)$
is a monoidally smooth morphism of log regular log
schemes then $\tilcalF(Y,M_Y)$ is the contraction of
$\phi^\st(\tilcalF(X,M_X))$. Furthermore, if $\phi$ is strict then
$\phi^\st(\tilcalF(X,M_X))=\tilcalF(X,M_X)\times_XY$.  \end{theorem2}
Strictness of $\phi$ in the last claim is not needed. To remove it
one should work out the assertion of Remark \ref{8-towpullrem}.
\begin{proof} We should only establish functoriality. Let
$(X',M_{X'})=\LogBl_\calI(X,M_X)^\sat$ be the first saturated blow
up of $\tilcalF^\rmlog(X,M_X)$. Set $\calJ=\phi^{-1}\calI M_Y$, then
$(Y',M_{Y'})=\LogBl_\calJ(Y,M_Y)^\sat$ is the first saturated blow up
of $\tilcalF^\rmlog(Y,M_Y)$ by functoriality of $\tilcalF^\rmlog$. By
Lemma \ref{8-logregblowlem}, $X'=\Bl_{\calI\calO_X}(X)^\nor$
and $Y'=\Bl_{\calJ\calO_Y}(Y)^\nor$, and using that
$\phi^{-1}(\calI\calO_X)\calO_Y=\calJ\calO_Y$ we obtain that $Y'$
is the strict transform of $X'$. It remains to inductively apply
the same argument to the other levels of the towers. The last claim
follows from Lemma \ref{8-logregblowlem}(ii).  \end{proof}

\begin{remark2} The same results hold for (non-saturated) monoidal
desingularization, which induces a (usual) desingularization of
log regular log schemes. For non-saturated log regular log schemes
(see \S\ref{8-nonsatursec}) one should first establish analogs of
Lemmas \ref{8-logregideallem} and \ref{8-logregblowlem} (where the
input in the second one does not have to be saturated). Then the
same proof as above applies.  \end{remark2}

\subsubsection{Canonical fans and associated
points}\label{8-canonfansec} By the {\em canonical fan} $\Fan(X)$ of $(X,M_X)$
we mean the set of maximal points of the log stratification (i.e. the maximal
points of the log strata). Alternatively, $\Fan(X)$ can be described as in
\cite[\S2.2]{Toric@Niziol} as the set of points $x\in X$ such that
$m_\ox$ coincides with the ideal $I_\ox\subseteq\calO_\ox$
generated by $\alpha(M_\ox-M_\ox^\times)$.

We provide $F=\Fan(X)$ with the induced topology and define $M_F$ to be the
restriction of $\oM_X$ onto $F$. For example, for a toric $k$-variety
$X=\Spec(k[Z])$, where $Z$ is a monoscheme, $(F,M_F)$ is isomorphic
to the sharpening of $Z$. More generally, if log scheme $(X,M_X)$
is Zariski then $(F,M_F)$ is a fan and the map $c{\colon}X\to F$ sending any
point to the maximal point of its log stratum is a fan chart of $X$.
This follows easily from the fact that such $(X,M_X)$ admits monoscheme
charts Zariski locally.

\begin{remark2}
In general, $(F,M_F)$ does not have to be a fan, but it seems probable that
it can be extended to a meaningful object playing the role of algebraic
spaces in the category of fans. We will not investigate this direction here.
\end{remark2}

\begin{lemma2}\label{8-assoclem} Let $X=(X,M_X))$ be an fs log
regular log scheme with a monoidal ideal $\calI\subseteq\calO_X$. Then:

(i) The set of associated points of $\calO_X/\calI$ is contained
in $\Fan(X)$.

(ii) The fans of $\Bl_\calI(X)$ and $\Bl_\calI(X)^\nor$ are contained
in the preimage of $\Fan(X)$.

(iii) For any tower of monoidal blow ups (resp. normalized
monoidal blow ups) $X_n\dashto X$, its set of associated points
is contained in $\Fan(X)$.  \end{lemma2} \begin{proof} (i) It
suffices to check that any $x\in X-\Fan(X)$ is not an associated
point of $\calO_X/\calI$. Since associated points are compatible
with flat morphisms, we can pass to the formal completion
$\hatX_x=\Spec(\hatcalO_{X,x})$. Let us consider first the more
difficult case when $A=\hatcalO_{X,x}$ is of mixed characteristic
$(0,p)$. By \refext{VI}{6-str loc log reg}, $A\iso B/(f)$ where $B=C(k)[[Q]][[\ut]]$,
$Q$ is a sharp monoid defining the log structure, $\ut=(t_1\. t_n)$,
and $f\in B$ equals to $p$ modulo $(Q\setminus\{1\},\ut)$. Note that
$n\ge 1$ as otherwise $Q\setminus\{1\}$ would generate the maximal
ideal of $A$. In this case, $x$ is the log stratum of $\hatX_x$
and hence $x$ is the maximal point of its log stratum in $X_x$,
which contradicts that $x\notin\Fan(X)$. The completion $\hatcalI$
of $\calI$ is of the form $IA$ for an ideal $I\subset Q$. We should
prove that $x\in\hatX_x$ is not an associated point of $A/\hatcalI$,
or, equivalently, $\depth(A/\hatcalI)\ge 1$. Since $B/(fB+\hatcalI
B)=A/\hatcalI$, it suffices to show that $\depth(B/\hatcalI B)\ge 2$
and $f$ is a regular element of $B/\hatcalI B$.

Regularity of $f$ follows from the following easy claim by
taking $C=C(k)$, $J=I$ and $R=Qt_1^\NN\dots t_n^\NN$: if $C$
is a domain, $R$ is a sharp monoid with an ideal $J$, $f$ is an
element of $C[[R]]$ with a non-zero free term, and $g\in C[[R]]$
satisfies $fg\in JC[[R]]$ then $g\in JC[[R]]$. Next, let us bound
the depth of $B/\hatcalI B$. In view of \cite[21.C]{CA@Matsumura},
depth is preserved by completions of local rings hence it suffices
to show that $\depth(D_r/ID_r)\ge 2$, where $D=C(k)[Q][\ut]$ and
$r=(m_Q,p,\ut)$ is the ideal generated by $m_Q=Q\setminus\{1\}$,
$p$ and $\ut$. Note that $D/ID\iso C(k)[\ut][Q\setminus I]$
as a $C(k)[\ut]$-module, hence it is a flat $C(k)[\ut]$-module
and the local homomorphism $C(k)[\ut]_{(p,\ut)}\to
D_r/ID_r$ is flat. By \cite[21.C]{CA@Matsumura},
$\depth(D_r/ID_r)\ge\depth(C(k)[\ut]_{(p,\ut)})=n+1\ge 2$, so
we have established the mixed characteristic case. In the equal
characteristic case we have that $A=k[[Q]][[\ut]]$, and the same
argument shows that $\depth(A/\hatcalI)\ge\depth(k[\ut]_{(\ut)})=n\ge
1$.

To prove (ii) we should check that if $x\in X-\Fan(X)$ then no
point of $\Fan(\Bl_\calI(X))$ sits over $x$. We will only check the
mixed characteristic case since it is more difficult. As earlier,
$\hatX_x=\Spec(A)$ where $A=B/(f)$ and $B=C(k)[[Q]][[\ut]]$. Note
that $\psi{\colon}\hatX_x\to X$ is a flat strict morphism of
log schemes. Hence $\psi$ is compatible with blow ups and it
maps the fans of $\hatX_x$ and $\Bl_\hatcalI(\hatX_x)$ to the
fans of $X_x$ and $\Bl_\calI(X)$, respectively. Therefore,
we should only check that $\Fan(\Bl_\hatcalI(\hatX_x))$ is
disjoint from the preimage of $x$. The latter blow up is covered
by the charts $V_a=\Spec(A[a^{-1}\hatcalI])$ with $a\in I$. Set
$P'=Q[a^{-1}I]$, $B'=C(k)[[P']][[t_1\. t_n]]/(f)$ and $A'=B'/(f)$
(where $f$ is as above). Then the $m_x$-adic completion of $V_a$
is $\hatV_a=\Spec(A')$. The ideal defining the closed point of
$\Fan(\hatV_a)$ is generated by the maximal ideal $m'$ of $P'$. This
ideal does not contain any $t_i$. Indeed, $t_i\notin m'B'+fB'$
because $t_i\notin fC(k)[[t_1\.t_n]]$ as $f-p\in(t_1\.t_n)$. Thus,
$\Fan(\hatV_a)$ is disjoint from the preimage of $x$, and hence the
same is true for $\Fan(V_a)$. This proves (ii) in the non-saturated
case, and the saturated case is dealt with similarly but with $P'$
replaced by its saturation. Finally, (iii) follows immediately from
(i) and (ii).  \end{proof}

\subsubsection{Independence of the log structure}\label{8-indepsec}
Dependence of $\tilcalF(X,M_X)$ on $M_X$ is a subtle question. In
this section we will use functoriality of $\tilcalF$ to prove that
$\tilcalF(X,M_X)$ is independent of $M_X$ in characteristic zero
(i.e., it is compatible with any automorphism of $X$). This will
cover our needs, but the restriction on the characteristic will
complicate our arguments later. Conjecturally, $\tilcalF(X,M_X)$
does not depend on $M_X$ at all and the following result of
Gabber supports this conjecture: if $P$ and $Q$ are two fine
sharp monoids and $k[[P]]\iso k[[Q]]$ for a field $k$ (of any
characteristic!) then $P\iso Q$. For completeness, we will give
a proof of this in \S\ref{8-reconstrsec}.

\begin{theorem2}\label{8-Morel 2.4.2} Let $\tilcalF$ be
a functorial normalized desingularization of regular qe
schemes of characteristic zero, and let $\tilcalF(X,M_X)$ be
the normalized desingularization of log regular log schemes
it induces (see Theorem \ref{8-deslogregth}). Assume that
$(X,M_X)$ and $(Y,M_Y)$ are saturated log regular log schemes
such that there exists an isomorphism $\phi{\colon}Y\iso X$
of the underlying schemes. Assume also that the maximal points
of the strata of the stratifications of $X$ and $Y$ by the
rank of $\overline{M}^{gp}$ are of characteristic zero. Then
$\tilcalF(X,M_X)$ and $\tilcalF(Y,M_Y)$ are compatible with $\phi$,
that is, $\tilcalF(Y,M_Y)=\phi^*(\tilcalF(X,M_X))$.  \end{theorem2}
\begin{proof} We can check the assertion of the theorem \'etale
locally. Namely, we can replace $X$ with a strict \'etale covering
$X'$ and replace $Y$ with $Y'=Y\times_XX'$ with the log structure
induced from $Y$. In particular, we can assume that the log
structures are Zariski, and so the canonical fans $\Fan(X)$ and
$\Fan(Y)$ are defined. Our assumption  on the maximal points actually
means that $\Fan(X)$ is contained in $X_\QQ=X\otimes_\ZZ\QQ$. By
Lemma \ref{8-assoclem}(iii) and \S\ref{8-assnorsec}, $\tilcalF(X,M)$
is the pushforward of its restriction onto $X_\QQ$, and
similarly for $Y$. So, it suffices to prove that the normalized
desingularizations of $X_\QQ$ and $Y_\QQ$ are compatible with
respect to $\phi\otimes_\ZZ\QQ{\colon}X_\QQ\iso Y_\QQ$. Thus,
it suffices to prove the theorem for $X$ and $Y$ of characteristic
zero, and, in the sequel, we assume that this is the case.

To simplify notation we identify $X$ and $Y$, and set $N_X=M_Y$. It
suffices to check that the blow up towers $\tilcalF(X,M_X)$ and
$\tilcalF(X,N_X)$ coincide after the base change to each completion
$\hatX_\ox=\Spec(\hatcalO_{X,\ox})$ at a geometric point $\ox$. By
\refext{VI}{6-str loc log reg}, we have that $\hatX_\ox\iso\Spec(k[[P]][[t_1\. t_n]])$,
where $P\iso\oM_{X,\ox}$, and the morphism of fs log
schemes $(\hatX_\ox,P)\to(X,M_X)$ is strict. By Theorem
\ref{8-monoidaldesth} the contracted pullback of $\tilcalF(X,M_X)$
to $\hatX_\ox$ coincides with $\tilcalF^\rmlog(\hatX_\ox,P)$. In
the same way, the contracted pullback of $\tilcalF(X,N_X)$
coincides with $\tilcalF^\rmlog(\hatX_\ox,Q)$, where
$\hatcalO_{X,\ox}\iso\Spec(k[[Q]][[t_1\. t_m]])$ (we use that $k$
is isomorphic to the residue field of this ring and hence depends
only on the ring $\hatcalO_{X,\ox}$).

Let us now recall how $\tilcalF^\rmlog(\hatX_\ox,P)$ is
constructed (Theorems \ref{8-fandesingth}, \ref{8-monodesinth}
and \ref{8-monoidaldesth}). We have the obvious strict
morphism $\hatX_\ox\to Z:=\Spec(\QQ[P][t_1\. t_n])$,
hence $\tilcalF^\rmlog(\hatX_\ox,P)$ is the pullback
of $\tilcalF^\rmlog(Z,P)$. The latter is the pullback of
$\tilcalF^\mono(P)=\tilcalF^\fan(P)$, which, in its turn, is
induced from $\tilcalF(Z)$. Therefore, $\tilcalF(Z,P)=\tilcalF(Z)$
and, by the functoriality of $\tilcalF$, its pullback to
$\hatX_\ox$ is $\tilcalF(\hatX_\ox)$. The same argument shows
that $\tilcalF(\hatX_\ox)$ is the contracted pullback of
$\tilcalF(X,N_X)$.

In order to prove that $\tilcalF(X,M_X)=\tilcalF(X,N_X)$ it only
remains to resolve the synchronization issues, i.e. to prove
equality without contractions. For this one should take the
union $S$ of the fans of $(X,M_X)$ and $(X,N_X)$, and consider
the morphism $\hatX_S:=\coprod_{x\in S}\hatX_\ox$ rather than
the individual completions. The pullbacks of $\tilcalF(X,M_X)$
and $\tilcalF(X,N_X)$ to $\hatX_S$ have no empty blow ups because
the fans contain all associated points of the towers by Lemma
\ref{8-assoclem}(iii). Hence the same argument as above shows that
they both coincide with $\tilcalF(\hatX_S)$.  \end{proof}

\begin{remark2} (i) Without taking the completion, $\tilcalF(X)$
does not even have to be defined as $X$ may be non-qe. In order
to pass to the completion we used functoriality of the monoidal
desingularization with respect to strict morphisms, which may be very
bad (e.g. non-flat) on the level of usual morphisms of schemes. Even
when $(X,M)$ is log regular, the formal completion morphism
$\hatX_x\to X$ can be non-regular in the non-qe case. However,
one can show that it is regular over the fan, and this is enough
to relate the (log) desingularization of $X$ and $\hatX_x$.

(ii) We used the very strong desingularization $\tilcalF$ from
Theorem \ref{8-qeth}. However, it is easy to see that only the
properties listed in Theorem \ref{8-magicboxth} were essential.
\end{remark2}

\subsection{Complements on non-saturated log regular log
schemes}\label{8-nonsatursec} For the sake of completeness, we
mention Gabber's results on non-saturated log regular log schemes
that will not be used. We only formulate results but do not give
proofs. Gabber defines a fine log scheme $(X,M_X)$ to be {\em log
regular} if its saturation is log regular in the usual sense. Assume
now that $(X,M_X)$ is fine and log regular. The key result that
lifts the theory off the ground is that Kato's Tor independence
extends to non-saturated log regular log schemes. Namely, if
$(X,M_X)$ admits a chart $\ZZ[P]$ and $I\subseteq P$ is an
ideal then ${\rm Tor}_1^{\ZZ[P]}(\calO_X,\ZZ[P]/I)=0$. For fs
log schemes this is due to Kato, and Gabber deduces the general
case using a non-flat descent. It then follows similarly to
the saturated case that if $(Y,M_Y)=\LogBl_\calI(X,M_X)$ then
$Y\iso\Bl_{\calI\calO_X}(X)$ and $(Y,M_Y)$ is log regular. In
addition, one shows that $(X,M_X)$ is saturated if and only if $X$ is
normal, and if $(Y,M_Y)=(X,M_X)^\sat$ then $Y\iso X^\nor$. Using
these foundational results on log regular fine log schemes one can
imitate the method of \S\ref{8-resolution of log schemes1} to extend
Lemma \ref{8-assoclem} to the non-saturated case. As a corollary,
one obtains an analog of Theorem \ref{8-Morel 2.4.2} for $\calF$
and $\calF^\rmlog$.

\subsection{Reconstruction of the monoid}\label{8-reconstrsec}
This section will not be used in the sequel. Its aim is to prove
that a fine torsion free monoid $P$ can be reconstructed from a ring
$A=k[[P]]$ where $k$ is a field. The main idea of the proof is that
an isomorphism $k[[P]]\iso A$ defines an action of the torus
$\Spec(k[P^\gp])$ on $A$, and any two maximal tori in $\Aut_k(A)$
are conjugate.

\subsubsection{Automorphism groups of complete
rings}\label{8-automsec} Let $k$ be a field and $A$ be a complete
local noetherian $k$-algebra with residue field $k$. Let $m$ denote
the maximal ideal and set $A_n=A/m^{n+1}$. Consider the algebraic
$k$-groups $G_n=\Aut_k(A_n)$, i.e. $G_n(B)=\Aut_B(A_n\otimes_kB)$
for any $k$-algebra $B$. They form a filtering projective family
$\dots G_2\to G_1\to G_0$, which we call algebraic pro-group
$G_\bullet$. Note that $G_\bullet$ induces a functor $G(B)=\lim_n G_n(B)$
and $G(k)=\Aut_k(A)$.

\begin{remark2} Gabber also considered more complicated filtering
families, but we stick to the simplest case we need.  \end{remark2}

\subsubsection{Stabilization}\label{8-stabsec} We say that an
algebraic pro-group $G_\bullet$ is {\em stable} if the homomorphisms
$G_{n+1}\to G_n$ are surjective for large enough $n$. Any algebraic
group can be stabilized as follows. For each $G_n$ and $i\ge
0$ let $G_{n,i}$ denote the image of $G_{n+i}$ in $G_n$. Then
$G_{n,0}\supseteq G_{n,1}\supseteq\dots$ is a decreasing sequence
of algebraic subgroups of $G$, hence it stabilizes on a subgroup
$G^\st_n\subseteq G_n$. The family $G^\st_\bullet$ with obvious
transition morphisms is an algebraic pro-subgroup of $G_\bullet$, and
it is clear from the definition that $G^\st_\bullet$ is stable. Note
also that $G_\bullet^\st$ is isomorphic to $G_\bullet$ at least in the
sense that $G(B)\iso\lim_n G^\st_n(B)$ for any $k$-algebra $B$.

\subsubsection{Maximal pro-tori}\label{8-maxtorsec} By a {\em
pro-torus} $T_\bullet$ in $G_\bullet$ we mean a {\em compatible}
family of tori $T_n\into G_n$  for $n\gg 0$, in the sense that
$\pi_n(T_{n+1})=T_n$. It is  called a {\em torus} if $\pi_n$'s
are isomorphisms for $n\gg 0$. A pro-torus is {\em split} or {\em
maximal split} if so are $T_n$ for $n\gg 0$. Pro-tori $T_\bullet$
and $T'_\bullet$ are conjugate if for $n\gg 0$ there exists a {\em
compatible} family of conjugations $G_n\to G_n$ taking $T_n$
to $T'_n$.

\begin{proposition2}\label{8-maxtorprop} Assume that $k$
is a field, $G_\bullet$ is a stable pro-algebraic $k$-group,
and $T_\bullet,T'_\bullet\into G_\bullet$ are split pro-tori. If
$T_\bullet$ is maximal then $T'_\bullet$ is conjugate to a
sub-pro-torus of $T_\bullet$. In particular, any two maximal
split pro-tori are conjugate.  \end{proposition2} \begin{proof}
It is a classical result that maximal split tori in algebraic
$k$-groups are conjugate. In particular, for each $n$ we can
move $T'_n$ into $T_n$ by a conjugation, and the only issue
is compatibility of the conjugations. Naturally, to achieve
compatibility we should lift conjugations inductively from $G_n$
to $G_{n+1}$. It suffices to prove that if $\pi_n$ is surjective and
$c_n{\colon}G_n\to G_n$ conjugates $T'_n$ into $T_n$ then it lifts
to $c_{n+1}{\colon}G_{n+1}\to G_{n+1}$ that conjugates $T'_{n+1}$
into $T_{n+1}$. By the stability assumption, we can lift $c_n$
to a conjugation $c'$ of $G_{n+1}$. It takes $T'_{n+1}$ to the
subgroup $H=KT_{n+1}$, where $K$ is the kernel of $\pi_n$. Since
maximal split tori in $H$ are conjugate and conjugation by elements
of $T_{n+1}$ preserves $T_{n+1}$, we can find a conjugation $c''$
by an element of $K$ that takes $c'(T'_{n+1})$ to $T_{n+1}$. Then
$c_{n+1}=c''c'$ is a lifting of $c_n$ as required \end{proof}

\begin{corollary2}\label{8-maxtorcor} Let $k$ be a
field and let $G_\bullet$ be an pro-algebraic $k$-group. If
$T_\bullet,T'_\bullet\subseteq G_\bullet$ are stable pro-tori and $T_\bullet$
is maximal as a split pro-torus of $G_\bullet$, then $T'_\bullet$ is
conjugate to a subtorus of $T_\bullet$.  \end{corollary2} \begin{proof}
Obviously, $T_n,T'_n\subseteq G^\st_n$ for $n\gg 0$. So, $T'_\bullet$
is conjugate to a subtorus of $T_\bullet$ already inside of $G^\st$
by Proposition \ref{8-maxtorprop}.  \end{proof}

\subsubsection{Certain tori in $\Aut_k(A)$}\label{8-torisec}
Any isomorphism $C[[P]]\iso A$, where $C$ is a complete
local $k$-algebra and $P$ is a sharp fine monoid, induces an
algebraic action of the split torus $T_P=\Spec(k[P^\gp])$ on $A$:
a character $\chi{\colon}P\to k$ acts on $C$ trivially and acts
on $p\in P$ by $p\mapsto \chi(p)p$ (and the action of $B$-points
$\chi{\colon}P\to B$ is analogous). Thus we obtain homomorphisms
$\psi{\colon}T_P\to G'_n$ which are injective for $n>0$. In
particular, the image is a split torus of $G$. Furthermore, we claim
that if $C=k$ then the torus is maximal (as a pro-torus). Indeed, if
$\phi\in\Aut_k(A)$ commutes with $T_P$ then its action on $k[[P]]$
preserves the $T_P$-eigenspaces $pk$ for $p\in P$ and on each
$pk$ it acts by multiplication by a number $\lambda(p)$. Clearly,
$\lambda{\colon}P\to k$ is a homomorphism and we obtain that $\phi$
belongs to $\psi(T_P(k))$ and corresponds to $\lambda\in T_P(k)$.

\begin{theorem2}\label{8-toritheor} Assume that $k$ is a field,
$P$ is a sharp fine monoid and $A=k[[P]]$. If $C$ is a complete
local $k$-algebra, $Q$ is a sharp fine monoid and $C[[Q]]\iso A$
then $C\iso k[[R]]$ and $P\iso Q\times R$ for a sharp fine
monoid $R$. In particular, $P$ is uniquely determined by $A$.
\end{theorem2} \begin{proof} Consider $G=\Aut_k(A)$ with split tori
$T_P,T_Q\into G$ corresponding to these isomorphisms. By maximality
of $T_P$ and Corollary \ref{8-maxtorcor} there exists a conjugation
of $G$ that maps $T_Q$ into $T_P$. This produces a new isomorphism
$C[[Q]]\iso A=k[[P]]$ that respects the grading, i.e. each $pk$
lies in some $qC$, and we obtain a surjective map $f{\colon}P\to Q$,
which is clearly a homomorphism. If $C=k$ then $f$ is an isomorphism
and we obtain that $P$ is determined by $A$.

Set $R_q=\prod_{p\in f^{-1}(q)}pk$. We have natural embeddings
$\prod_{p\in f^{-1}(q)}pk\into qC$ which are all isomorphisms
because $A=\prod_{q\in Q} R_q$ is isomorphic to $C[[Q]]=\prod_{q\in
Q}qC$. In particular, for $R=R_1$ we have that $C=\prod_{q\in R}
qk=k[[R]]$. Therefore,\linebreak[4] $A\iso k[[R]][[Q]]\iso k[[R\times Q]]$,
and since the monoid is determined by $A$ we obtain that $P\iso
Q\times R$.  \end{proof}

\section{Proof of Theorem \ref{8-TH. MODIFICATION} -- preliminary
steps}\label{8-prelimsec} The goal of \S\ref{8-prelimsec} is to
reduce the proof of Theorem \ref{8-TH. MODIFICATION} to the case
when the following conditions are satisfied: (1) $X$ is regular,
(2) the log structure is given by an snc divisor $Z$ which is {\em
$G$-strict} in the sense that for any $g\in G$ and a component
$Z_i$ either $gZ_i=Z_i$ or $gZ_i\cap Z_i=\emptyset$, (3) $G$ acts
freely on $X\setminus Z$ and for any geometric point $\ox\to X$
the inertia group $G_\ox$ is abelian.

\subsection{Plan} A general method for constructing a $G$-equivariant
morphism $f$ as in Theorem \ref{8-TH. MODIFICATION} is to construct
a tower of $G$-equivariant morphisms of log schemes $X'=X_n\dashto
X_0=X$, where the underlying morphisms of schemes are normalized
blow ups along $G$-stable centers sitting over $Z\cup T$, such that
various properties of the log scheme $X_i$ with the action of $G$
gradually improve to match all assertions of the Theorem. To simplify
notation, we will, as a rule, replace $X$ with the new log scheme
after each step. The three conditions above will be achieved in
three steps as follows.

\subsubsection{Step 1. Making $X$ regular.}\label{8-step1}
This is achieved by the saturated log blow up
tower $\tilcalF(X,Z){\colon}X'\dashto X$ from Theorem
\ref{8-deslogregth}. In particular, the morphism $X'\to X$ is
even log smooth. In the sequel, we assume that $X$ is regular, in
particular, $Z$ is a normal crossings divisor by \refext{VI}{6-1.7}. We will
call $Z$ the {\em boundary} of $X$. In the sequel, these conditions
will be preserved, so let us describe an appropriate restriction
on further modifications.

\subsubsection{Permissible blow ups}\label{8-permissec} After
Step 1 any modification in the remaining tower will be of the form
$f{\colon}(X',Z')\to(X,Z)$ where $X'=\Bl_V(X)$, $Z'=f^{-1}(Z\cup
V)$ and $V$ has {\em normal crossings} with $Z$, i.e. \'etale
locally on $X$ there exist regular parameters $t_1\. t_d$ such that
$Z=V(\prod_{i=1}^lt_i)$ and $V=V(t_{i_1}\. t_{i_m})$. We call such
modification {\em permissible} and a blow up $f{\colon}X'\to X$
is called {\em permissible} (with respect to the boundary $Z$)
if it underlies a permissible modification. Since there is an
obvious bijective correspondence between permissible modifications
and blow ups we will freely pass from one to another. Note that
$Z'=f^\st(Z)\cup E'$, where $E'=f^{-1}(V)$ is the exceptional
divisor.

\subsubsection{Permissible towers}\label{8-permistowersec} A {\em
permissible modification tower} $(X_d,Z_d)\dashto(X_0,Z_0)=(X,Z)$
is defined in the obvious way and we say that a blow up tower
$X_d\dashto X$ is permissible if it underlies such a modification
tower. Again, we will freely pass between permissible towers of
these types. Note that $Z_i=Z^\st_i\cup E_i$, where $Z_i^\st$
is the strict transform of $Z$ under $h_i{\colon}X_i\dashto X$
and $E_i$ is the exceptional divisor of $h_i$ (i.e. the union of
the preimages of the centers of $h_i$).

\begin{remark2}\label{8-permisrem} (i) Consider a permissible tower
as above. It is well known that for any $i$ one has that $X_i$
is regular, $Z_i$ is normal crossings and $E_i$ is even snc. For
completeness, let us outline the proof. Both claims follow from
the following: if $Z$ is snc then $Z_i$ is snc. Indeed, the claim
about $E_i$ follows by taking $Z=\emptyset$ and the claim about
$Z_i$ can be checked \'etale locally, so we can assume that
$Z$ is snc. Finally, if $Z$ is snc then $Z_i$ is snc by Lemma
\ref{8-regblowlem} below.

(ii) Permissible towers are very common in embedded desingularization
because they do not destroy regularity of the ambient scheme and the
normal crossings (or snc) property of the boundary (or accumulated
exceptional divisor). Even when one starts with an empty boundary,
a non-trivial boundary appears after the first step, and this
restricts the choice of further centers. Actually, any known
self-contained proof of embedded desingularization constructs a
permissible resolution tower.  \end{remark2}

\subsubsection{$G$-permissible towers}\label{8-Gpermissec} In
addition, we will only blow up $G$-equivariant centers $V$. So,
$f{\colon}X'=\Bl_V(X)\to X$ is $G$-equivariant and the exceptional
divisor $E=f^{-1}(V)$ is regular and $G$-equivariant and hence
$G$-strict. Such a blow up (or their tower) will be called {\em
$G$-permissible}. It follows by induction that the exceptional
divisor of such a tower is $G$-strict.

\subsubsection{Step 2. Making $Z$ snc and $G$-strict}\label{8-step2}
Consider the stratification of $Z$ by multiplicity: a point $z\in Z$
is in $Z^n$ if it has exactly $n$ preimages in the normalization
of $Z$. Note that $\{Z^n\}$ is precisely the log stratification
as defined in \S\ref{8-logstratsec}. By
{\em depth} of the stratification we mean the maximal $d$ such
that $Z^d\neq\emptyset$. In particular, $Z^d$ is the only closed
stratum. Step 2 proceeds as follows: $X_{i+1}\to X_i$ is the blow
up along the closed stratum of $Z_i^\st$.

\begin{remark2}\label{8-step2rem} What we use above is
the standard algorithm that achieves the following two
things: $Z'$ is snc and $Z^\st=\emptyset$ (see, for example,
\cite[7.2]{Alterations@de_Jong}). Even when $Z$ is snc, the second
property is often used in the embedded desingularization algorithms
to get rid of the old components of the boundary.  \end{remark2}

\subsubsection{Justification of Step 2}\label{8-step2proof}
Since the construction is well known, we just sketch the
argument. First, observe that $Z^d$ has normal crossings with
$Z$, that is, $X'=\Bl_{Z^d}(X)\to X$ is permissible. Thus, $Z'$
is normal crossings and hence $Z^\st$ is also normal crossings. A
simple computation with blow up charts shows that the depth of
$Z^\st$ is $d-1$ (for example, one can work \'etale-locally, and
then this follows from Lemma \ref{8-weaksnclem} below). It follows
by induction that the tower produced by Step 2 is permissible,
of length $d$ and with $Z_d^\st=\emptyset$. So, $Z_d=E_d$ is snc
by Remark \ref{8-permisrem} and $G$-strict by \S\ref{8-Gpermissec}.

\subsubsection{Step 3. Making the inertia groups abelian
and the action of $G$ on $X\setminus Z$ free}\label{8-step3}
Recall (\refext{VI}{6-stratification}) that the inertia strata are of the
form $X_H=X^H\setminus\cup_{H\subsetneq H'}X^{H'}$. Step 3 runs
analogously to Step 2, but this time we will work with the inertia
stratification of $X$ instead of the log stratification, and will
have to apply the same operation a few times. Let us first describe
the blow up algorithm used in this step; its justification will be
given in \S\ref{8-step3proof}.

Let $f_{\{X^H\}}{\colon}X'\dashto X$ denote the following blow
up tower. First we blow up the union of all closed strata
$X_H$. In other words, $V_0$ is the union of all non-empty
minimal $X^H$, i.e. non-empty $X^H$ that do not contain $X^K$
with $\emptyset\subsetneq X^K\subsetneq X^H$. Next, we consider
the family of all strict transforms of $X^H$ and blow up the union
of the non-empty minimal ones, etc. Obviously, the construction is
$G$-equivariant. We will prove in Proposition \ref{8-towerprop} that
$f_{\{X^H\}}$ is permissible of length bounded by the length of the
maximal chains of subgroups. Also, we will show in \S\ref{8-step3sec}
that $f_{\{X^H\}}$ decreases all non-abelian inertia groups,
so the algorithm of Step 3 goes as follows: until all inertia
groups become abelian, apply $f_{\{X^H\}}{\colon}X'\dashto X$
(i.e. replace $(X,Z)$ with $(X',Z')$).

\subsection{Justification of Step 3}\label{8-step3proof}

\subsubsection{Weakly snc families}\label{8-weaksncsec} Assume that
$X$ is regular, $Z\into X$ is an snc divisor, and $\{X_i\}_{i\in
I}$ is a finite collection of closed subschemes of $X$. For any
$J\subseteq I$ we denote by $X_J$ the scheme-theoretic intersection
$\cap_{j\in J}X_j$. The family $\{X_i\}$ is called {\em weakly
snc} if each $X_i$ is nowhere dense and $X_J$ is regular. The
family $\{X_i\}$ is called {\em weakly $Z$-snc} if it is weakly
snc and each $X_J$ has normal crossings with $Z$. In particular,
$\{X_i\}_{i\in I}$ is weakly snc (resp. weakly $Z$-snc) if and only
if the family $\{X_J\}_{\emptyset\neq J\subseteq I}$ is weakly snc
(resp. weakly $Z$-snc).

\begin{remark2}\label{8-weaksncrem} (i) Here is a standard criterion
of being an snc divisor, which is often taken as a definition. Let
$D\into X$ be a divisor with irreducible components $\{D_i\}_{i\in
I}$. Then $D$ is snc if and only if it is weakly snc and each
irreducible component of $D_J$ is of codimension $|J|$ in $X$.

(ii) The condition on the codimension is essential. For example,
$xy(x+y)=0$ defines a weakly snc but not snc divisor in
$\AA^2_k=\Spec(k[x,y])$.

(iii) The criterion from (i) implies that if $X$ is qe then the snc
locus of $D$ is open -- it is the complement of the union of singular
loci of $D_J$'s. (Note that this makes sense for all points of $X$
because $D$ is snc at a point $x\in X\setminus D$ if and only if
$X=D_\emptyset$ is regular at $x$.)
\end{remark2}

\begin{lemma2}\label{8-Zsnclem} Let $(X,Z)$ and $G$ be as achieved
in Step 2. Then the family $\{X^H\}_{1\neq H\subseteq G}$ is
weakly $Z$-snc.
\end{lemma2}

\begin{proof} Recall that for any
subgroup $H\subseteq G$ the fixed point subscheme $X^H$ is regular
by Proposition \refext{VI}{6-pts fixes reguliers}, and $X^H$ is nowhere dense for $H\neq 1$
by generic freeness of the action of $G$. Since for any pair of
subgroups $K,H\subseteq G$ we have that $X^H\times_XX^K=X^{KH}$,
the family is weakly snc. It remains to show that $Y=X^H$ has normal
crossings with $Z$. Note that it is enough to consider the case when
$X$ is local with closed point $x$ and $G=H$. The cotangent spaces
at $x$ will be denoted $T^*X=m_{X,x}/m^2_{X,x}$, $T^*Y$, etc. Their
dual spaces will be called the tangent spaces, and denoted $TX$,
$TY$, etc. Let $\phi^*:T^*X\onto T^*Y$ denote the natural map
and let $\phi:TY\into TX$ denote its dual. We will systematically
use without mention that $|G|$ is coprime to ${\rm char} k(x)$,
in particular, the action of $G$ on $T^*X$ is semi-simple.

The proof of \refext{VI}{6-pts fixes reguliers} also shows that for any point $x\in X^H$, the
tangent space $T_x(X^H)$ is isomorphic to $(T_xX)^H$. In particular,
$TY\iso(TX)^G$ and hence $\phi^*$ maps $(T^*X)^G\subset
T^*X$ isomorphically onto $T^*Y$. Therefore, $U=\Ker(\phi^*)$
is the $G$-orthogonal complement to $(T^*X)^G$, i.e. the
only $G$-invariant subspace such that $(T^*X)^G\oplus U\iso
T^*X$. Let $Z_i=V(t_i)$, $1\le i\le n$ be the components of $Z$ and
let $dt_i\in T^*X$ denote the image of $t_i$. By $Z$-strictness
of $G$, each line $L_i=\Span(dt_i)$ is $G$-invariant, so $G$
acts on $dt_i$ by a character $\chi_i$. Without restriction of
generality, $\chi_1\.\chi_l$ for some $0\le l\le n$ are the only
trivial characters. In particular, $L=\Span(dt_1\. dt_n)$ is the
direct sum of $L^G=\Span(dt_1\. dt_l)$ and its $G$-orthogonal
complement $L\cap U$, which (by uniqueness of the complement)
coincides with $\Span(dt_{l+1}\.dt_n)$. Complete the basis of
$L$ to a basis $\{dt_1\.dt_n,e_1\. e_m\}$ of $T^*X$ such that
$\{dt_{l+1}\.dt_n,e_1\. e_r\}$ for some $r\le m$ is a basis of
$U$ and choose functions $s_1\.s_m$ on $X$ so that $ds_j=e_j$
and $s_1\. s_r$ vanish on $Y$. Clearly, $t_1\. t_n,s_1\.s_m$ is a
regular family of parameters of $\calO_{X,x}$, so the lemma would
follow if we prove that $Y=V(t_{l+1}\.t_n,s_1\.s_r)$.

Since $Y$ is regular and $\Ker(\phi^*)$ is spanned by the images of
$t_{l+1}\. t_n,s_1\.s_r$, we should only check that these functions
vanish on $Y$. The $s_j$'s vanish on $Y$ by the construction, so we
should check that $t_i$ vanishes on $Y$ whenever $l<i\le n$. Using
the functorial definition from \refext{VI}{6-stratification} of the subscheme of fixed
points, we obtain that $Z_i^G=Z_i\times_XX^G$, hence $Z_i\times_XY$
is regular by \refext{VI}{6-pts fixes reguliers}. However, $TY$ is contained $TZ_i$, which
is the vanishing space of $dt_i$, hence we necessarily have that
$Y\into Z_i$.  \end{proof}

\subsubsection{Snc families}\label{8-sncsec} A family of nowhere
dense closed subschemes $\{X_i\}_{i\in I}$ is called {\em snc}
(resp. {\em $Z$-snc}) at a point $x$ if $X$ is regular at $x$ and
there exists a regular family of parameters $t_j\in\calO_{X,x}$ such
that in a neighborhood of $x$ each $X_i$ (resp. and each irreducible
component $Z_k$ of $Z$) passing through $x$ is given by the vanishing of a subfamily
$t_{j_1},\dots,t_{j_l}$. Note that the family $\{X_i\}_{i\in I}$
is $Z$-snc if and only if the union $\{X_i\}\cup\{Z_k\}$ is snc. A
family is {\em snc} if it is so at any point of $X$ (in particular,
$X$ is regular).

\begin{remark2}\label{8-sncrem} (i) It is easy to see that the family
$\{X^H\}_{H\subseteq G}$ is snc whenever $G$ is abelian. Indeed,
it suffices to show that for any point $x$ there exists a basis of
$T_xX$ such that each $T_x(X^H)$ is given by vanishing of some of the
coordinates. But this is so because the action of $G$ on $T_xX$ is
(geometrically) diagonalizable and $T_x(X^H)=(T_xX)^H$. In general,
the family $\{X^H\}_{H\subseteq G}$ does not have to be snc, as
the example of a dihedral group $D_n$ with $n\ge 3$ acting on the
plane shows.

(ii) If $Z\into X$ is an snc divisor with components $Z_i$ and
$V\into X$ is a closed subscheme then the family $\{Z_i,V\}$ is
snc if and only if $V$ has normal crossings with $Z$.  \end{remark2}

The transversal case of the following lemma can be deduced from
\egalong{IV}{4}{\S19.1}, but we could not find the general case in
the literature (although it seems very probable that it should have
appeared somewhere).  \begin{lemma2}\label{8-snclem} Any weakly snc
family with $|I|=2$ is snc.  \end{lemma2} \begin{proof} We should
prove that if $X=\Spec(A)$ is a regular local scheme and $Y$, $Z$
are regular closed subschemes such that $T=Y\times_XZ$ is regular
then there exists a regular family of parameters $t_1\.t_n\in A$
such that $Y$ and $Z$ are given by vanishing of some set of these
parameters. Let $m$ be the maximal ideal of $A$, and let $I$, $J$
and $K=I+J$ be the ideals defining $Y$, $Z$ and $T$, respectively. By
$T^*X=m/m^2$, $T^*Y$, etc., we denote the cotangent spaces at the
closed point of $X$. Note that $I/mI\iso\Ker(T^*X\to T^*Y)$,
and similar formulas hold for $J/mJ$ and $K/mK$. Indeed, we can
choose the parameters so that $Y=V(t_1\.t_l)$ and then the images
of $t_1\.t_l$ form a basis both of $I/mI$ and $\Ker(T^*X\to T^*Y)$.

Now, let us prove the lemma. Assume first that $Y$ and $Z$ are
transversal, i.e. $T^*X\into T^*Y\oplus T^*Z$. Choose elements
$t_1\. t_{l+k}$ such that $Y=V(t_1\.t_l)$, $Z=V(t_{l+1}\.t_{l+k})$,
$l=\codim(Y)$ and $k=\codim(Z)$. Then the images $dt_i\in T^*X$
of $t_i$ are linearly independent because $dt_1\.dt_l$ span
$\Ker(T^*X\to T^*Y)$ and $dt_{l+1}\.dt_{l+k}$ span $\Ker(T^*X\to
T^*Z)$. Hence we can complete $t_i$'s to a regular family of
parameters by choosing $t_{l+k+1}\.t_n$ such that $dt_1\.dt_n$
is a basis of $T^*X$. This proves the transversal case, and to
establish the general case it now suffices to show that if $Y$
and $Z$ are not transversal then there exists an element $t_1\in
m\setminus m^2$ which vanishes both on $Y$ and $Z$. (The we can
replace $X$ with $X_1=V(t_1)$ and repeat this process until $Y$
and $Z$ are transversal in $X_a=V(t_1\.t_a)$.) Tensoring the
exact sequence $0\to I\cap J\to I\oplus J\to K\to 0$ with $A/m$
we obtain an exact sequence $$(I\cap J)/m(I\cap J)\to I/mI\oplus
J/mJ\stackrel\phi\to K/mK\to 0$$ The failure of transversality
is equivalent to non-injectivity of $\phi$, hence there exists
an element $f\in I\cap J$ with a non-zero image in $I/mI\oplus
J/mJ$. Thus, $f\in m\setminus m^2$ and we are done.  \end{proof}

\subsubsection{Blowing up the minimal strata of a weakly
snc family}\label{8-weakblowsec} Given a weakly snc family
$\{X_i\}_{i\in I}$, we say that a scheme $X_J$ with $J\subseteq I$
is {\em minimal} if it is non-empty and any $X_{J'}\subsetneq X_J$
is empty. Also, we will need the following notation: if $Z\into X$
is a closed subscheme and $D\into X$ is a Cartier divisor with the
corresponding ideals $\calI_Z,\calI_D\subset\calO_X$, then $Z+D$
is the closed subscheme defined by the ideal $\calI_Z\calI_D$.

\begin{proposition2}\label{8-weaksnclem} Assume that $X$ is regular,
$Z\into X$ is an snc divisor with irreducible components $Z_1\. Z_l$,
and $\{X_i\}_{i\in I}$ is a $Z$-snc (resp. weakly $Z$-snc) family
of subschemes. Let $V$ be the union of all non-empty minimal
subschemes $X_J$, $f{\colon}X'=\Bl_V(X)\to X$, $X'_i=f^\st(X_i)$
and $Z'=f^{-1}(Z\cup V)$. Then

(i) $X'$ is regular and $Z'$ is snc.

(ii) The family $\{X'_i\}_{i\in I}$ is $Z'$-snc (resp. weakly
$Z'$-snc).

(iii) For any $J\subseteq I$, the scheme-theoretical intersection
$X'_J=\cap_{j\in J}X'_j$ coincides with $f^\st(X_J)$.

(iv) For any $J\subseteq I$ the total transform $X_J\times_XX'$ is of
the form $X'_J+D'_J$ where $D'_J$ is the divisor consisting of all
connected components of $E'=f^{-1}(V)$ contained in $f^{-1}(X_J)$.
\end{proposition2} \begin{proof} We start with the following lemma.

\begin{lemma2}\label{8-regblowlem} Assume that $X$ is regular,
$Z$ is an snc divisor and $V\into Y$ are closed subschemes
having normal crossings with $Z$. Let $f{\colon}X'\to X$ be
the blow up along $V$, $Y'=f^\st(Y)$, $Z'=f^{-1}(Z\cup V)$, and
$E'=f^{-1}(V)$. Then $Z'$ is snc, $Y'$ has normal crossings with
$Z'$ and $Y\times_XX'=Y'+E'$.  \end{lemma2} \begin{proof} The proof
is a usual local computation with charts. Take any point $u\in V$
and choose a regular family of parameters $t_1\.t_n$ at $u$ such
that $Y$ (resp. $V$, resp. $Z$) are given by the vanishing of
$t_1\.t_m$ (resp. $t_1\.t_l$, resp. $\prod_{i\in I}t_i$), where
$0\le m\le l\le n$ and $I\subseteq\{1\. n\}$. Locally over $u$
the blow up is covered by $l$ charts, and the local coordinates on
the $i$-th chart are $t'_j$ such that $t'_j=t_j$ for $j>l$ or $j=i$
and $t'_j=\frac{t_j}{t_i}$ otherwise. On this chart, $Y\times_XX'$
(resp. $Y'$, resp. $E'$, resp. $Z'$) is given by the vanishing of
$t_1\.t_m$ (resp. $t'_1\.t'_m$, resp. $t'_i$, resp. $\prod_{j\in
I\cup\{i\}}t'_j$), hence the lemma follows.  \end{proof}

The lemma implies (i). In addition, it follows from the
lemma that $f^\st(X_J)$ has normal crossings with $Z'$ and
$X_J\times_XX'=f^\st(X_J)+D'_J$. Thus, (iii) implies (iv), and
(iii) implies (ii) in the case when the family $\{X_i\}_{i\in I}$
is weakly $Z$-snc. Note also that if this family is $Z$-snc then
locally at any point $x\in X$ there exists a family of regular
parameters $\underline{t}=\{t_1\. t_n\}$ such that each $X_J$ and
each component of $Z$ is given by the vanishing of a subfamily of
$\underline{t}$ locally in a neighborhood of $x$. Then the same local
computation as was used in the proof of Lemma \ref{8-regblowlem}
proves also claims (ii) and (iii) of the proposition. So, it remains
to prove (iii) when the family is weakly snc. It suffices to prove
that if (iii) holds for $X_J$ and $X_K$ then it holds for $X_{J\cup
K}$. Moreover, (iii) does not involve the boundary so we can assume
that $Z=\emptyset$. It remains to note that $\{X_J,X_S\}$ is an
snc family by Lemma \ref{8-snclem}, hence our claim follows from
the snc case.  \end{proof}

\subsubsection{Blow up tower of a weakly $Z$-snc
family}\label{8-weaktowersec} Let $X$ be a regular scheme, $Z$
be an snc divisor and $\{X_i\}_{i\in I}$ be a weakly $Z$-snc
family. By the blow up tower $f_{\{X_i\}}$ of $\{X_i\}$ we mean
the following tower: the first blow up $h_1:X_1\to X_0=X$ is
along the union of all non-empty minimal schemes of the form $X_J$
for $\emptyset\neq J\subseteq I$, the second blow up is along the
union of all non-empty minimal schemes of the form $h_1^\st(X_J)$
for $\emptyset\neq J\subseteq I$, etc.

\begin{proposition2}\label{8-towerprop} Keep the above notation. Then
the tower $f_{\{X_i\}}$ is permissible with respect to $Z$ and
its length equals to the maximal length of chains $\emptyset\neq
X_{J_1}\subsetneq\dots\subsetneq X_{J_d}$ with $\emptyset\neq
J_d\subsetneq\dots\subsetneq J_1\subseteq I$. Furthermore, the strict
transform of any scheme $X_J$ is empty and the total transform of
$X_J$ is a Cartier divisor.  \end{proposition2} \begin{proof} The
claim about the length is obvious. Let $f_{\{X_i\}}:X_d\dashto X_0=X$
and let $h_n:X_n\dashto X_0=X$ be its $n$-th truncation. For each
$i\in I$ set $X_{n,i}=h_n^\st(X_i)$ and for each $J\subseteq I$ set
$X_{n,J}=\cap_{i\in J}X_{n,i}$. Using Proposition \ref{8-weaksnclem}
and straightforward induction on length, we obtain that the family
$\{X_{n,i}\}_{i\in I}$ is weakly $Z_n$-snc, $X_{n,J}=h_n^\st(X_J)$,
the blow up $X_{n+1}\to X_n$ is along the union of non-empty
minimal $X_{n,J}$'s, and $X_J\times_XX_n=X_{n,J}+D_{J,n}$, where
$D_{J,n}$ is a divisor. So, the tower is permissible, and since
$X_{d,J}=\emptyset$ we also have that $X_J\times_XX_d$ is a divisor.
\end{proof}

\begin{remark2} We will not need this, but it is easy to deduce from
the proposition that on the level of morphisms the modifications
$X_d\to X$ is isomorphic to the blow up along $\prod_{\emptyset\neq
J\subseteq I}\calI_{X_J}$.  \end{remark2}

\subsubsection{Justification of Step 3}\label{8-step3sec} The
blow up tower $f_{\{X^H\}}{\colon}X'\dashto X$ from Step 3 is
$G$-equivariant in an obvious way, and it is permissible by
Proposition \ref{8-towerprop}. In addition, $Z'=f^{-1}(Z)\cup
f^{-1}(\cup_{1\neq H\subseteq G} X^H)$ and, since $G$ acts freely
on $X'\setminus f^{-1}(\cup_{1\neq H\subseteq G} X^H)\iso
X\setminus\cup_{1\neq H\subseteq G}X^H$, it also acts freely on
$X'\setminus Z'$. It remains to show that applying $f_{\{X^H\}}$ we
decrease all non-abelian inertia groups. Namely, for any $x'\in X'$
mapped to $x\in X$ we want to show that either $G_\ox$ is abelian
or the inclusion $G_{\ox'}\subsetneq G_\ox$ is strict.

Let $H\subseteq G$ be any non-abelian subgroup with commutator
$K=[H,H]$. Since $X^K\times_XX'$ is a divisor by Proposition
\ref{8-towerprop}, the universal property of blow ups implies that
$X'\to X$ factors through $Y=\Bl_{X^K}(X)$. On the other hand, it
is proved in \refext{VI}{6-eclatements} that $Y\to X$ is an $H$-equivariant
blow up, and if a geometric point $\ox\to X$ with $G_\ox=H$ lifts to
a geometric point $\oy\to Y$ then $G_\oy\subsetneq H$. Therefore,
the same is true for the $G$-equivariant modification $X'\to X$,
and we are done.

\section{Proof of Theorem \ref{8-TH. MODIFICATION} -- abelian
inertia}\label{8-endsec}

\subsection{Conventions}\label{8-notationendsec} Throughout
\S\ref{8-endsec} we assume that $(X,Z)$ and $G$ satisfy
all conditions achieved at Steps 1, 2, 3, and our aim is to
construct a modification $f_{(G,X,Z)}{\colon}X'\to X$ as in Theorem
\ref{8-TH. MODIFICATION}. Unless specially mentioned, we do not assume 
that $X$ is qe. This is done in order to isolate the only place where
this assumption is needed (existence of rigidifications).

\subsection{Outline of our method and other
approaches}\label{8-outlinesec}

\subsubsection{Combinatorial nature of the problem}\label{8-combsec}
On the intuitive level it is natural to expect that "everything
relevant to our problem" should be determined by the following
"combinatorial" data: the log structure of $X$, the inertia
stratification of $X$ by $X_H:=X^H\setminus\cup_{H'\subsetneq
H}X^{H'}$ and the representations of the inertia groups on the
tangent spaces (which are essentially constant along $X_H$). This
combinatorial nature is manifested in both approaches to the problem
that we describe below.

\subsubsection{Combinatorial algorithm}\label{8-algorsec}
The most natural approach is to seek for a "combinatorial
algorithm" that iteratively blows up the closures of {\em
log-inertia strata} (i.e. intersections of a log stratum with
an inertia stratum) or some generalization thereof. Actually, our (very simple) algorithms in
\S\ref{8-prelimsec} were of this type. The algorithm should be
governed by a combinatorial data, such as the number of components
of $Z$ through a point $x$ and the history of their appearance
(similarly to the desingularization algorithms), the representation
of $G_\ox$ on $T_x$ plus some history (e.g. representation of $G_\ox$
on the tangent spaces to the components of $Z$), etc.

It is natural to expect that building such an algorithm would lead to
a relatively simple proof of Theorem \ref{8-TH. MODIFICATION}. In
particular, it would be non-sensitive to quasi-excellence
issues. Unfortunately, despite partial positive results, we could
not construct such an algorithm.
Thus, the question whether such an algorithm exists remains
open.\footnote{F. Pop told to the second author that he has a plan
of constructing such a combinatorial algorithm.}

\subsubsection{Our method}\label{8-ourmethodsec} A general plan
of our method is as follows. In \S\ref{8-localsec} we will show
that such a modification $f$ exists \'etale-locally on the base
if $X$ is qe. A priori, our construction will be canonical up to
an auxiliary choice, but then we will prove in \S\ref{8-globalsec}
that actually it is independent of the choice and hence descends to a
modification $f$ as required. To prove independence we will show in
\S\ref{8-functsec} that the construction is functorial with respect
to strict {\em inert} morphisms (i.e. morphisms that preserve both
the log and the inertia structures, see \S\ref{8-inertsec}). The
latter is a manifestation of the "combinatorial nature" of our
algorithm.

\subsection{Local construction}\label{8-localsec}

\subsubsection{Very tame action and Zariski
topology}\label{8-verytamesec} Note that the log structure on
$(X,Z)$ is Zariski since $Z$ is snc. Thus, it will be convenient
to describe very tame action in terms of Zariski topology. We say
that the action of $G$ is {\em very tame at} a point $x\in X$ if
for any geometric point $\ox$ over $x$ the action is very tame at
$\ox$. Since the log structure is Zariski, the $G_\ox$-equivariant
log scheme $\Spec(\calO_{X,\ox})$ is independent of the choice of
$\ox$ up to an isomorphism. In particular, the action is very tame
at $x$ if and only if it is very tame at a single geometric point
$\ox$ above $x$.

\begin{lemma2}\label{8-tamelocallem} Assume that $(G,X,Z)$ is as
in \S\ref{8-notationendsec}, $x\in X$ is a point, and $T$ is the
set of points of $X$ at which the action is very tame. Then

(i) $T$ is open,

(ii) $x\in T$ if and only if locally at $x$ the log stratification of
$X$ is finer than the inertia stratification (i.e. any log stratum
is contained in an inertia stratum).  \end{lemma2} \begin{proof}
Choose a geometric point $\ox$ above $x$. If $x\in T$ then by
Corollary \refext{VI}{6-tres modere dans voisinage} the action is very tame on a suitable \'etale
neighborhood $X'$ of $\ox$. The image of $X'$ in $X$ is an open
neighborhood of $x$ which is contained in $T$. Thus, $T$ is open. To
prove the second claim we note that the conditions (i) and (ii) from
\refext{VI}{6-3.1} are automatically satisfied at $\ox$. Indeed, (i) is satisfied
because the action is tame by assumption of \ref{8-TH. MODIFICATION}
and (ii) is satisfied because $Z$ is snc and $G$-strict. Since
log and inertia stratifications are compatible with the strict
henselization morphism, condition (iii) from \refext{VI}{6-3.1} is satisfied
at $\ox$ if and only if the log stratification is finer than the
inertia stratification at $x$. This proves (ii).  \end{proof}

\subsubsection{Admissibility}\label{8-admissec} In the sequel,
by saying that $X$ is {\em admissible} we mean that the action
of $G$ on $X$ is admissible in the sense of \sga{1}{V}{1.7}
(e.g. $X$ is affine). This is needed to ensure that $X/G$ exists as
a scheme. An alternative would be to allow $X/G$ to be an algebraic
space (and $(X/G,Z/G)$ to be a log algebraic space).

\subsubsection{Rigidification}\label{8-rigidsec} By a {\em
rigidification} of $X$ we mean a $G$-equivariant normal crossings
divisor $\oZ$ that contains $Z$ and such that the action of $G$ on
the log regular log scheme $\oX=(X,\oZ)$ is very tame. If $\oZ$ is
snc then we say that the rigidification is {\em strict}. Sometimes,
by a rigidification of $X=(X,Z)$ we will mean the log scheme $\oX$
itself. Our construction of a modification $f$ uses a rigidification,
so let us first establish local results on existence of the latter.

\begin{lemma2}\label{8-rigidlem} Let $X,Z,G$ be as
in \S\ref{8-notationendsec}. Assume that $X=\Spec(A)$ is a
local scheme with closed point $x$, $G=G_\ox$, and $\mu_N\subset
A$, where $N$ is the order of $G$. Then $X$ possesses a strict
rigidification.  \end{lemma2} \begin{proof} Choose $t_1\. t_n\in A$
such that $Z_i=(t_i)$ are the components of $Z$. By $G$-strictness
of $Z$, for any $g\in G$ we have that $Z_i=(gt_i)$ locally at
$x$, in particular, the tangent space to each $Z_i$ at $x$ is
$G$-invariant. Now, we can use averaging by the $G$-action to make
the parameters $G$-equivariant. Namely, $G$ acts on $dt_i$ by a
character $\chi_i$ and replacing $t_i$ with $\frac{1}{|G|}\sum_{g\in
G}\frac{gt_i}{\chi_i(g)}$ we do not change $dt_i$ and achieve that
$gt_i=\chi_i(g)t_i$.

The action of $G$ on the cotangent space at $x$ is diagonalizable
because $G$ is abelian and $\mu_N\subset k(x)$. In particular,
we can complete the family $dt_1\. dt_n$ to a basis $dt_1\. dt_l$
such that $t_i\in\calO_{X,x}$ and $G$ acts on each $dt_i$ by
a character $\chi_i$. Then $t_1\. t_l$ is a regular family of
parameters of $\calO_{X,x}$ and using the same averaging procedure
as above we can make them $G$-equivariant. Take now $\oZ$ to be the
union of all divisors $(t_i)$ with $1\le i\le l$. The action of $G$
on $(X,\oZ)$ is very tame at $x$ because it is the only point of
its log stratum. So, it remains to use Lemma \ref{8-tamelocallem}.
\end{proof}

\subsubsection{Inert morphisms}\label{8-inertsec} Let $(X,Z)$
and $G$ be as in \S\ref{8-notationendsec}. Our next aim is to
find an \'etale cover $f:(Y,T)\to(X,Z)$ which "preserves" the
log-inertia structure of $(X,Z)$ and such that $(Y,T)$ admits a
rigidification. The condition on the log structure is obvious:
we want $f$ to be strict, i.e. $f^{-1}(Z)=T$. Let us introduce a
restriction related to the inertia groups.

Assume that $(Y,T)$ with an action of $H$ is another such triple,
and let $\lambda:H\to G$ be a homomorphism. A $\lambda$-equivariant
morphism $f{\colon}(Y,T)\to(X,Z)$ will be called {\em inert} if
for any point $y\in Y$ with $x\in X$ the induced homomorphism of
inertia groups $G_\oy\to G_\ox$ is an isomorphism. In particular,
the inertia stratification of $Y$ is the preimage of the inertia
stratification of $X$.

\begin{lemma2}\label{8-rigidlem2} Let $X$, $Z$ and $G$ be as above,
and assume that $X$ is qe. Then there exists a $G$-equivariant
surjective \'etale inert strict morphism $h{\colon}(Y,T)\to(X,Z)$
such that $Y$ is affine and $(Y,T)$ possesses a strict
rigidification.  \end{lemma2} \begin{proof} First, we note that the
problem is local on $X$. Namely, it suffices for any point $x\in
X$ to find a $G_\ox$-equivariant \'etale inert strict morphism
$h{\colon}(Y,T)\to(X,Z)$ such that $Y$ is affine, $x\in h(Y)$, any
point $x'\in h(Y)$ satisfies $G_{\ox'}\subseteq G_\ox$, and $(Y,T)$
admits a strict rigidification. Indeed, $h$ can be extended to a
$G$-equivariant morphism $Y\times_X(X\times G/G_\ox)=\coprod_{g\in
G/G_\ox}Y_g\to X$, where each $Y_g$ is isomorphic to $Y$ and
the morphism $Y_g\to X$ is obtained by composing $Y\to X$ with
$g{\colon}X\to X$. Clearly, the latter morphism is \'etale, inert,
and strict, and by quasi-compactness of $X$ we can combine finitely
many such morphism to obtain a required cover of $(X,Z)$.

Now, fix $x\in X$ and consider the $G_\ox$-invariant neighborhood
$X'=X\setminus\cup_{H\nsubseteq G_\ox}X^H$ of $x$. We
will work over $X'$, so the condition  $G_{\ox'}\subseteq
G_{\ox}$ will be automatic. Let $N$ be the order of
$G_\ox$, and consider the $G_\ox$-equivariant morphism
$f{\colon}Y=X'\times\Spec(\ZZ[\frac{1}{N},\mu_N])\to X$ (with
$G_\ox$ acting trivially on the second factor). Let $T=f^{-1}(Z)$
and let $y$ be any lift of $x$. It suffices to show that
$(Y,T)$ admits a rigidification in an affine $G_\ox$-invariant
neighborhood of $y$ (such neighborhoods form a fundamental
family of neighborhoods of $y$). By Lemma \ref{8-rigidlem}, the
localization $Y_y=\Spec(\calO_{Y,y})$ with the restriction $T_y$
of $T$ possesses a strict rigidification $\oT_y$. Clearly, $\oT_y$
extends to a divisor $\oT$ with $T\into\oT\into Y$ and we claim
that it is a rigidification in a neighborhood of $y$. Indeed, $\oT$
is snc at $y$, hence it is snc in a neighborhood of $y$ by Remark
\ref{8-weaksncrem}, and it remains to use Lemma \ref{8-tamelocallem}.
\end{proof}

\subsubsection{Main construction}\label{8-constrsec} Assume,
now, that $X=(X,Z)$ is admissible and admits a rigidification
$\oZ$. We are going to construct a $G$-equivariant modification
$$f_{(G,X,Z,\oZ)}{\colon}(X',Z')\to(X,Z)$$ such that $G$ acts very
tamely on the target and $f_{(G,X,Z,\oZ)}$ is independent of the
rigidification. The latter is a subtle property (missing in the
obvious modification $(X,\oZ)\to(X,Z)$), and it will take us a
couple of pages to establish it.

The quotient log scheme $\oY=(Y,\oT)=(\oX/G,\oZ/G)$ is log regular
by Theorem \refext{VI}{6-tres modere}, hence by Theorem \ref{8-monoidaldesth}
there exists a functorial saturated log blow up tower
$\oh=\tilcalF^\rmlog(\oY){\colon}\oY'=(Y',\oT')\to\oY$ with a regular
and log regular source. Let $\of{\colon}\oX'=(X',\oZ')\to\oX$
be the pullback of $\oh$ (as a saturated log blow up tower, see
\S\ref{8-logpullsec}), then $\alpha'{\colon}\oX'\to\oY'$ is a Kummer
\'etale $G$-cover because $\alpha{\colon}\oX\to\oY$ is so by \refext{VI}{6-tres modere}
as the square $$ \xymatrix{(X',\oZ') \ar[d]^{\alpha'} \ar[r]^\of&
(X,\oZ) \ar[d]^\alpha \\ (Y',\oT') \ar[r]^\oh & (Y,\oT)} $$ is
cartesian in the category of fs log scheme.

Since $G$ acts freely on $U=X-Z$, $V=U/G$ is regular and $\oT|_V$
is snc. In particular, $\oh$ is an isomorphism over $V$ and hence
$\of$ is an isomorphism over $U$.  We claim that the Weil divisor
$T=Z/G$ of $Y$ is $\QQ$-Cartier (it does not have to be Cartier,
as the orbifold case with $X=\AA^2$, $Z=\AA^1$ and $G=\{\pm 1\}$
shows). Indeed, it suffices to check this \'etale-locally at a
point $y\in Y$. In particular, we can assume that $\mu_N\subset
\OO_{Y,y}$, where $N=|G|$. Then, as we showed in the proof of Lemma
\ref{8-rigidlem}, $Z$ can be locally defined by equivariant
parameters, in particular, $Z=V(f)$ where $G$ acts on $f$ by
characters. Therefore, $f^N$ is $G$-fixed, and we obtain that
$T$ is the reduction of the Cartier divisor $C$ of $Y$ given by
$f^N=0$. (The same argument applied to $T\times_YX$ shows that $NT$
is Cartier, so $T$ is $\QQ$-Cartier.) So, $C'=C\times_YY'$ is a
Cartier divisor whose reduction is $T'=\oh^{-1}(T)$. Since $Y'$
is regular and $T'$ lies in the snc divisor $\oT'$, we obtain that
$T'$ is itself an snc divisor.

Let $Z'$ denote the divisor $\alpha'^{-1}(T')=\of^{-1}(Z)$. Since
$X'\to Y'$ is \'etale over $V=Y'-T'$, the morphism of log schemes
$(X',Z')\to(Y',T')$ is a Kummer \'etale $G$-cover. This follows
from a variant of the classical Abhyankar's lemma (IX 2.1), which
is independent of the results of the present expos\'e.

In particular, $X'=(X',Z')$ is log regular and it follows that the
action of $G$ on $X'$ is very tame. We define $f_{(G,X,Z,\oZ)}$
to be the modification $X'\to X$.

\begin{remark2}\label{8-constrrem} (i) Note that $X'\to X$ satisfies
all conditions of Theorem \ref{8-TH. MODIFICATION} because the action
is very tame and $G$ acts freely on $X'\setminus f^{-1}(Z)$. So, we
completed the proof in the case when $(X,Z)$ admits a rigidification
$\oZ$. Our last task will be to get rid of the rigidification.

(ii) The only dependence of our construction on the rigidification
is when we construct the resolution of $(Y,\oT)$. Conjecturally,
it depends only on the scheme $Y$, and then $(Y',T')$, and hence
also $(X',Z')$, would depend only on $(X,Z)$. Recall that we
established in Theorem \ref{8-Morel 2.4.2} the particular case
of this conjecture when all maximal points of the log strata of
$(Y,\oT)$ are of characteristic zero. Hence independence of the
rigidification is unconditional in this case, and, fortunately,
this will suffice.  \end{remark2}

\subsubsection{Finer structure of
$f_{(G,X,Z,\overline{Z})}$}\label{8-finersec} Obviously, the saturated
log blow up tower $\of{\colon}(X',\oZ')\to(X,\oZ)$ depends on the
rigidification, and this is the reason why we prefer to consider
the modification $f_{(G,X,Z,\oZ)}{\colon}(X',Z')\to(X,Z)$
instead. However, there is an additional structure on
$f_{(G,X,Z,\oZ)}$ that has a chance to be independent of $\oZ$, and
which should be taken into account. By \S\ref{8-desingloregsec}, the
modification of schemes $f{\colon}X'\to X$ has a natural structure of
a normalized blow up tower $X_\bullet$ with $X=X_0$ and $X'=X_n$. Note
also that the tower contains no empty blow ups because this is true
for $\tilcalF^\rmlog(X/G,\oZ/G)$ and $f_{(G,X,Z,\oZ)}$ is its strict
transform with respect to the surjective morphism $X\to X/G$.

Note also that the log structure on $(X',Z')$ is reconstructed
uniquely from $f$ because $Z'=f^{-1}(Z)$ and $(X',Z')$ is
saturated and log regular. Therefore, it is safe from now on to
view $f_{(G,X,Z,\oZ)}$ as a normalized blow up tower of $X$, but
the modification of log schemes $(X',Z')\to(X,Z)$ will also be
denoted as $f_{(G,X,Z,\oZ)}$.

\begin{remark2}\label{8-finnormrem} Although we do not assume that
$X$ is qe, all normalizations in the tower $f_{(G,X,Z,\oZ)}$ are
finite. This happens because they underly saturations of fine log
schemes, which are always finite morphisms.  \end{remark2}

\subsection{Functoriality}\label{8-functsec} Clearly, the
construction of $f$ depends canonically on $(G,X,Z,\oZ)$,
i.e. is compatible with any automorphism of such
quadruple. Our next aim is to establish functoriality with
respect to strict inert $\lambda$-equivariant morphisms
$\phi{\colon}(H,Y,T,\oT)\to(G,X,Z,\oZ)$ (i.e. morphisms
that "preserve the combinatorial structure"). For this one
has first to study the quotient morphism of log schemes
$\tilphi{\colon}(Y/H,\oT/H)\to(X/G,\oZ/G)$.

\subsubsection{Log structure of the quotients}\label{8-logquotsec}
Recall the following facts from Proposition \refext{VI}{6-tres modere ponctuel}(b)
and its proof. Assume that $X=(X,M_X)$ is an fs log scheme provided with a
very tame action of a group $G$. After replacing $X$ with its strict
localization at a geometric point $\ox$, it admits an equivariant
chart $X\to\Spec(\Lambda[Q])$, where $\Lambda=\ZZ[1/N,\mu_N]$
for the order $N$ of $G_\ox$, $Q$ is an fs monoid and the
action of $\oG_x$ is via a pairing $\chi{\colon}G_\ox\otimes
Q\to\mu_N$. Moreover, if $P\subseteq Q$ is the maximal submonoid with
$\chi(G_\ox\otimes P)=1$ then $\Spec(\Lambda[Q])\to\Spec(\Lambda[P])$
is a chart of $X\to X/G_\ox$. Now, let us apply this description
to the study of $\tilphi$.

\begin{proposition2}\label{8-logquotprop} Assume that fs log schemes
$X$, $Y$ are provided with admissible very tame actions of groups $G$
and $H$, respectively,  $\lambda{\colon}H\to G$ is a homomorphism,
and $\phi{\colon}Y\to X$ is a strict inert $\lambda$-equivariant
morphism. Then the quotient morphism $\tilphi{\colon}Y/H\to X/G$
is strict.  \end{proposition2} \begin{proof} Fix a geometric point
$\oy$ of $Y$ and let $\ox$ be its image in $X$. It suffices to
show that $\tilphi$ is strict at the image of $\oy$ in $Y/H$. The
morphism $Y/H_\oy\to Y/H$ is strict (and \'etale) over the image of
$\oy$, and the same is true for $X$. Therefore we can replace $H$
and $G$ with $H_\oy\iso G_\ox$, and then we can also replace $X$
and $Y$ with their strict localizations at $\ox$ and $\oy$. Now,
the morphism $X\to\tilX=X/G_\ox$ admits an equivariant chart
$h{\colon}\Spec(\Lambda[Q])\to\Spec(\Lambda[P])$ as explained before
the proposition. Since $\phi$ is strict, the induced morphism
$Y\to\Spec(\Lambda[Q])$ is also a chart and hence $h$ is also a
chart of $Y\to\Y/G_\oy$. Thus, $\tilphi$ is strict.  \end{proof}

\subsubsection{An application to functoriality of $f_\bullet$} Assume
that $(G,X,Z,\oZ)$ is as earlier, and let $(H,Y,T,\oT)$ be another
such quadruple (i.e. $(Y,T)$ with the action of $H$ satisfies
conditions of Steps 1, 2, 3 and $(Y,\oT)$ is its rigidification).

\begin{corollary2}\label{8-logquotcor} Assume that
$\lambda{\colon}H\to G$ is a homomorphism and $\phi{\colon}Y\to
X$ is a $\lambda$-equivariant inert morphism such that
$T=Z\times_XY$ and $\oT=\oZ\times_XY$. Then $f_\bullet$ is
compatible with $\phi$ in the sense that $f_{(H,Y,T,\oT)}$ is
the contraction of $\phi^\st(f_{(G,X,Z,\oZ)})$. In addition,
$\phi^\st(f_{(G,X,Z,\oZ)})=f_{(G,X,Z,\oZ)}\times_XY$.
\end{corollary2} \begin{proof} The morphism
$\psi{\colon}(Y,\oT)\to(X,\oZ)$ is strict, hence the morphism of
quotients is strict by Proposition \ref{8-logquotprop}, and by
functoriality of saturated monoidal desingularization we obtain
that $\tilcalF^\rmlog(Y/H,\oT/H)$ is the contracted pullback
of $\tilcalF^\rmlog(X/G,\oZ/G)$. So, both $f_{(H,Y,T,\oT)}$ and
$f_{(G,X,Z,\oZ)}$ are obtained as the contraction of the strict
transform of $\tilcalF^\rmlog(X/G,\oZ/G)$. The first claim of the
Corollary follows.

Furthermore, $f_{(G,X,Z,\oZ)}$ underlies a log blow
up tower of $(X,\oZ)$ which is the strict transform of
$\tilcalF^\rmlog(X/G,\oZ/G)$, and the same is true for
$f_{(H,Y,T,\oT)}$. Since $\psi$ is strict it follows from Lemma
\ref{8-logregblowlem}(ii) that the strict transform is a pullback,
i.e. $\phi^\st(f_{(G,X,Z,\oZ)})=f_{(G,X,Z,\oZ)}\times_XY$.
\end{proof}

\subsubsection{Localizations and completions}\label{8-generalizsec}
In particular, it follows that the construction of $f_\bullet$
is compatible with localizations and completions. Namely, if
$x\in X$ is a point, $X_x=\Spec(\calO_{X,x})$, $Z_x=Z\times_XX_x$
and $\oZ_z=\oZ\times_XX_x$, then $f_{(G_\ox,X_x,Z_z,\oZ_x)}$ is
the contraction of $f_{(G,X,Z,\oZ)}\times_XX_x$. Similarly, if
$\hatX_x=\Spec(\hatcalO_{X,x})$, $\hatZ_x=Z\times_X\hatX_x$
and $\widehat{\oZ}_x=\oZ\times_X\hatX_x$, then
$f_{(G_\ox,\hatX_x,\widehat{Z}_x,\widehat{\oZ}_x)}$ is the
contraction of $f_{(G,X,Z,\oZ)}\times_X\hatX_x$.

\subsection{Globalization}\label{8-globalsec} To complete the
proof of Theorem \ref{8-TH. MODIFICATION} it suffices to show that
$f_{(G,X,Z,\oZ)}$ is independent of $\oZ$, and hence the local
constructions glue to a global normalized blow up tower. The main
idea is to simultaneously lift two rigidifications to characteristic
zero and apply Theorem \ref{8-Morel 2.4.2}.

\subsubsection{Independence of rigidification}\label{8-indeprigsec}
We start with the case of complete local rings. Then the problem
is solved by lifting to characteristic zero and referencing to
\ref{8-Morel 2.4.2}. The general case will follow rather easily.

\begin{lemma2}\label{8-formallem} Keep assumptions on $(X,Z)$
and $G$ as in \S\ref{8-notationendsec} and assume, in addition,
that $X=\coprod_{i=1}^m\Spec(A_i)$ where each $A_i$ is a complete
noetherian regular local ring with a separably closed residue
field. Then for any pair of rigidifications $\oZ$ and $\oZ'$ the
equality $f_{(G,X,Z,\oZ)}=f_{(G,X,Z,\oZ')}$ holds.  \end{lemma2}
\begin{proof} Almost the whole argument runs independently on
each irreducible component, so assume first that $X=\Spec(A)$
is irreducible. By Remark \ref{8-constrrem}(ii), it suffices
to consider the case when ${\rm char}(k)=p>0$, so let $C(k)$
be a Cohen ring of $k$. Note that we can work with $H=G_\ox$
instead of $G$ because $f_{(H,X,Z,\oZ)}=f_{(G,X,Z,\oZ)}$ by
Corollary \ref{8-logquotcor}. Since $H$ acts trivially on $k$,
for any element $t\in A$ its $H$-averaging is an element of $A^H$
with the same image in the residue field. Hence $k$ is the residue
field of $A^H$ and the usual theory of Cohen rings provides a
homomorphism $C(k)\to A^H$ that lifts $C(k)\to A^H/m_{A^H}$. Note
that $\oZ$ and $\oZ'$ are snc because each $A_i$ is strictly
henselian. Using averaging on the action of $H$ again, we can find
regular families of $H$-equivariant parameters $\uz=(z_1\. z_d)$
and $\uz'=(z'_1\. z'_d)$ such that $Z=V(\prod_{i=1}^lz_i)$,
$z'_i=z_i$ for $1\le i\le l$, $\oZ=V(\prod_{i=1}^nz_i)$ and
$\oZ'=V(\prod_{i=1}^{n'}z'_i)$. Explicitly, the action on $z_i$
(resp. $z'_i$) is by a character $\chi_i{\colon}H\to k(x)^\times$
(resp. $\chi'_i{\colon}H\to k(x)^\times$).

Since the image of $\uz$ is a basis of the cotangent space at $x$,
we obtain a surjective homomorphism $f{\colon}B=C(k)[[t_1\. t_d]]\to
A$ taking $t_i$ to $z_i$. Provide $B$ with the action of $H$
which is trivial on $C(k)$ and acts on $t_i$ via $\chi_i$, in
particular, $f$ is $H$-equivariant. Let us also lift each $z'_i$
to an $H$-equivariant parameter $t'_i\in B$. For $i\le l$ we
take $t'_i=t_i$, and for $i>l$ we first choose any lift and then
replace it with its $\chi_i$-weighted $H$-averaging. Consider
the regular scheme $Y=\Spec(B)$ with $H$-equivariant snc
divisors $T=V(\prod_{i=1}^lz_i)$, $\oT=V(\prod_{i=1}^nz_i)$ and
$\oT'=V(\prod_{i=1}^{n'}z'_i)$.

Since $H$ acts vary tamely on $(X,\oZ)$, it acts trivially on
$V(z_1\. z_n)=\Spec(k[[z_{n+1}\. z_d]])$ and we obtain that
$\chi_i=1$ for $i>n$. Therefore, $H$ also acts trivially on
$\Spec(B/(t_1\. t_n))=\Spec(C(k)[[t_{n+1}\. t_d]])$ and we
obtain that the action on $(Y,\oT)$ is very tame. Since the
closed immersion $j:X\to Y$ is $H$-equivariant and strict, and
$\oZ=\oT\times_YX$, Corollary \ref{8-logquotcor} implies that
$f_\bullet$ is compatible with $j$, i.e., $f_{(H,X,Z,\oZ)}$ is the
contracted strict transform of $f_{(H,Y,T,\oT)}$. The same argument
applies to the rigidifications $\oZ'$ and $\oT'$, so it now suffices
to show that $f_{(H,Y,T,\oT)}=f_{(H,Y,T,\oT')}$. For this we observe
that maximal points of log strata of the log schemes $(Y/H,\oT/H)$
and $(Y/H,\oT'/H)$ are of characteristic zero, hence the latter
equality holds by Theorem \ref{8-Morel 2.4.2}.

Finally, let us explain how one deals with the case of $m>1$. First
one finds an $H$-equivariant strict closed immersion $i:X\to Y$
such that $\oZ$ and $\oZ'$ extend to rigidifications $\oT$ and
$\oT'$ of $(Y,T)$, and the maximal points of the log strata of
$(Y/H,\oT/H)$ and $(Y/H,\oT'/H)$ are of characteristic zero. For
this we apply independently the above construction to the connected
components of $X$. Once $i$ is constructed, the same reference to
\ref{8-Morel 2.4.2} shows that $f_{(H,Y,T,\oT)}=f_{(H,Y,T,\oT')}$
and hence $f_{(H,X,Z,\oZ)}=f_{(H,X,Z,\oZ')}$ \end{proof}

\begin{corollary2}\label{8-indepprop} Let $(X,Z)$ and $G$
be as in \S\ref{8-notationendsec} and assume that the action
is admissible. Then for any choice of rigidifications $\oZ$
and $\oZ'$ we have that $f_{(G,X,Z,\oZ)}=f_{(G,X,Z,\oZ')}$.
\end{corollary2} \begin{proof} For a point $x\in X$ let
$\widehat{\calO^\sh_{X,x}}$ denote the completion of the strict
henselization of $\calO_{X,x}$. It suffices to check that for
any point $x$ the normalized blow up towers $f_{(G,X,Z,\oZ)}$
and $f_{(G,X,Z,\oZ')}$ pull back to the same normalized blow up
towers of $\hatX^\sh_x=\Spec(\widehat{\calO^\sh_{X,x}})$ (with
respect to the morphism $\hatX^\sh_x\to X$). Indeed, any normalized
blow up tower $\calX=(X_\bullet,V_\bullet)$ is uniquely determined by
its centers $V_i$. For each $i$ the morphism $Y_i=\coprod_{x\in
X}X_i\times_X\hatX^\sh_x\to X_i$ is faithfully flat, hence $V_i$
is uniquely determined by $V_i\times_{X_i}Y_i$, which is the center
of $\calX\times_X\coprod_{x\in X}\hatX^\sh_x$.

By \S\ref{8-generalizsec},
$f_{(G_\ox,\hatX^\sh_x,Z\times_X\hatX^\sh_x,\oZ\times_X\hatX^\sh_x)}$
is the contracted pullback of $f_{(G,X,Z,\oZ)}$, and an analogous
result is true for $f_{(G,X,Z,\oZ')}$. Thus the contracted
pullbacks are equal by Lemma \ref{8-formallem}. We have, however,
to worry also for the synchronization, i.e. to establish equality
of non-contracted pullbacks. For this we will use the following
trick. The towers $f_{(G,X,Z,\oZ)}$ and $f_{(G,X,Z,\oZ')}$ are
of finite length, hence there exists a finite subset $S\subset X$
such that the image of any center of either of these towers has a
non-empty intersection with $S$. Set $\hatX^\sh_S=\coprod_{s\in
S}\hatX^\sh_s$, then the pullbacks of $f_{(G,X,Z,\oZ)}$ and
$f_{(G,X,Z,\oZ')}$ to $\hatX^\sh_S$ are already contracted. Now,
in order to compare the pullbacks to $\hatX^\sh_x$, consider the
pullbacks to $\hatX^\sh_x\coprod\hatX^\sh_S$. They are contracted,
so Lemma \ref{8-formallem} (which covers disjoint unions) implies
that these pullbacks are equal. Restricting them onto $\hatX^\sh_x$
we obtain equality of non-contracted pullbacks to $\hatX^\sh_x$.
\end{proof}

\begin{remark2}\label{8-indeprem} (i) The above corollary implies
that the modification $f_{(G,X,Z,\oZ)}$ depends only on $(G,X,Z)$,
so it will be denoted $f_{(G,X,Z)}$ in the sequel. At this stage,
$f_{(G,X,Z)}$ is defined only when $X$ is admissible and $(X,Z)$
admits a rigidification.

(ii) Corollaries \ref{8-logquotcor} and \ref{8-indepprop} imply
that $f_{(G,X,Z)}$ is functorial with respect to equivariant strict
inert morphisms.  \end{remark2}

\subsubsection{Theorem \ref{8-TH. MODIFICATION} -- end of proof
for}\label{8-endproofsec} Let $X=(X,Z)$ be as assumed in
\S\ref{8-notationendsec}, and suppose that $X$ is qe. By Lemma
\ref{8-rigidlem2} there exists a surjective \'etale inert strict
morphism $h{\colon}X_0\to X$ such that $X_0$ is affine and possesses
a rigidification. Then $X_1=X_0\times_XX_0$ is affine and also
admits a rigidification (e.g. the preimage of that of $X_0$ by
one of the canonical projections). By Remark \ref{8-indeprem}(i),
$X_0$ and $X_1$ possess normalized blow up towers $f_{(G,X_0,Z_0)}$
and $f_{(G,X_1,Z_1)}$, which are  compatible with both projections
$X_1\to X_0$ by Remark \ref{8-indeprem}(ii). It follows that
$f_{(G,X_0,Z_0)}$ is induced from a unique normalized blow up
tower of $X$ that we denote as $f_{(G,X,Z)}$. This modification
satisfies all assertions of Theorem \ref{8-TH. MODIFICATION}
because $f_{(G,X_0,Z_0)}$ does so by Remark \ref{8-constrrem}(i).

\subsection{Additional properties of $f_{(G,X,Z)}$}\label{8-addsec}
Finally, let us formulate an addendum to Theorem
\ref{8-TH. MODIFICATION} where we summarize additional properties
of the constructed modification of $(X,Z)$. At this stage
we drop any assumptions on $(X,Z)$ beyond the assumptions of
\ref{8-TH. MODIFICATION}. By $f_{(G,X,Z)}$ we denote below the
entire modification from Theorem \ref{8-TH. MODIFICATION} that also
involves the modifications of Steps 1, 2, 3.

\begin{theorem2}\label{8-addendum} Keep assumptions of Theorem
\ref{8-TH. MODIFICATION}. In addition to assertions of the theorem,
the modifications $f_{(G,X,Z)}$ can be constructed uniformly for all
triples $(G,X,Z)$ such that the following properties are satisfied:

(i) Each $f_{(G,X,Z)}$ is provided with a structure of a normalized
blow up tower and its centers are contained in the preimages of
$Z\cup T$.

(ii) For any homomorphism $\lambda:H\to G$ the construction is
functorial with respect to $\lambda$-equivariant inert strict
morphisms $(Y,T)\to(X,Z)$.
\end{theorem2}
\begin{proof} The total modification $f_{(G,X,Z)}$ is obtained
by composing four modifications $f_1$, $f_2$, $f_3$ and $f_4$:
the modifications from Steps 1, 2, 3 and the modification we have
constructed in \S\ref{8-endsec}. Recall that $f_1$ and $f_4$ are constructed
as normalized blow up towers. Modifications $f_2$ and $f_3$ are
permissible blow up towers, hence they are also normalized blow up
towers with the same centers. This establishes the first part of (i).

Claim (ii) follows from the following functorialities: $f_1$
and $f_2$ are functorial with respect to all strict morphisms,
$f_3$ is functorial with respect to all inert morphisms, and $f_4$
is functorial with respect to strict inert morphisms. To prove
the second part of (i) we use (ii) to restrict $f_{(G,X,Z)}$ onto
$U=X\setminus{Z\cup T}$. Then $U$ is a regular scheme with a trivial
log structure which is acted freely by $G$. It follows from the
definitions of $f_1$, $f_2$, $f_3$ and $f_4$ that they are trivial
for such $U$. So, $f_{(G,U,\emptyset)}$ is the trivial tower, and
hence all centers of $f_{(G,X,Z)}$ are disjoint from the preimage
of $U$.
\end{proof}

\ifx\danslelivre\undefined
\bibliography{../fichiers-configuration-latex/bibliographie}
\bibliographystyle{../fichiers-configuration-latex/smfalpha}
\end{document} \fi
